\documentclass[10pt,a4paper]{article}

% --- PAQUETES DE IDIOMA Y CODIFICACIÓN ---
\usepackage[utf8]{inputenc}
\usepackage[T1]{fontenc}
\usepackage[english,spanish,es-nodecimaldot,es-noshorthands]{babel}

% --- GEOMETRÍA Y TIPOGRAFÍA ---
\usepackage[top=2.0cm,bottom=2.0cm,left=2.0cm,right=2.0cm]{geometry}
\usepackage{lmodern}
\usepackage{microtype}

% --- MATEMÁTICA Y FÍSICA ---
\usepackage{amsmath,amssymb,amsfonts,mathtools,bm}
\usepackage{slashed}

% --- TABLAS, FIGURAS Y COLORES ---
\usepackage{xcolor}
\usepackage{graphicx}
\usepackage{booktabs}
\usepackage{array}
\usepackage{cite}

% --- ESQUEMAS CONCEPTUALES VECTORIALES ---
% (añadido en la revisión de redacción: diagramas de las Figs. 1, 2, 4 y 5)
% Se usan únicamente librerías estándar de PGF/TikZ; no se requiere
% ningún paquete externo adicional.
\usepackage{tikz}
\usetikzlibrary{decorations.pathmorphing,calc,positioning,arrows}

% --- JERARQUÍA VISUAL DE ECUACIONES (dos niveles) ---
% Nivel 1 -> marco (amsmath "boxed"): los resultados que sostienen
%            el hilo argumental completo del artículo.
% Nivel 2 -> fondo sombreado (macro propia): resultados importantes de cada
%            tramo, con realce ligero y sin marco.
\definecolor{colordest}{RGB}{238,242,248}
\definecolor{colorborde}{RGB}{110,130,160}
\newcommand{\dest}[1]{{\setlength{\fboxsep}{4pt}\colorbox{colordest}{$\displaystyle #1$}}}
\newcommand{\lambdabar}{\mathchar'26\mkern-9mu\lambda}

% --- HIPERVÍNCULOS ---
\usepackage{hyperref}
\hypersetup{
    colorlinks=true,
    linkcolor=blue,
    citecolor=red,
    urlcolor=blue,
    pdfauthor={Paul Martínez y equipo de trabajo},
    pdftitle={Del momento magnético anómalo del protón a su radio de carga: estructura electromagnética del nucleón en QED}
}

% --- TÍTULO Y AUTORES ---
\title{\textbf{\Large Del Momento Magnético Anómalo del Protón a su Radio de Carga}\\[1.0ex]
\large Estructura Electromagnética del Nucleón en QED: Identidad de Gordon, Factores de Forma de Sachs y Separación de Rosenbluth}

\author{
    \textbf{Equipo de Trabajo Subatómica}\\[0.5ex]
    \textit{Docente Evaluador:} Prof. Jhon Morales Aponte, Ph.D.\\[0.5ex]
    \small Departamento de Física, Facultad de Ciencias\\
    \small Universidad Nacional de Colombia — Sede Bogotá\\
    \small Asignatura: Introducción a la Física Subatómica (2026-II)
}
\date{\today}

\begin{document}

\maketitle
\vspace{-1.5em}
\begin{quote}
\small
\textbf{Resumen}---En 1933, Otto Stern midió el momento magnético del protón y obtuvo $\mu_p \approx 2.79\,\mu_N$, casi tres veces el valor que la ecuación de Dirac impone sin margen de ajuste a cualquier fermión puntual de espín $1/2$ ($g = 2$). Aquella anomalía fue la primera evidencia de que el protón ocupa volumen, y con ella quedó planteada una pregunta cuantitativa: ¿cómo se convierte una discrepancia magnética en una medida del tamaño del protón, expresada en femtómetros? Este artículo responde esa pregunta reconstruyendo paso a paso, sin saltos algebraicos, la cadena completa que une la electrodinámica cuántica (QED) con un número medible, adoptando como punto de partida y eje articulador la sección eficaz diferencial de la dispersión elástica $e^- p \to e^- p$ que Marshall N. Rosenbluth publicó en 1950. Demostramos primero que la invariancia de Lorentz y la conservación de paridad admiten un vértice hadrónico covariante con tres funciones escalares, y que la identidad de Ward---la conservación de la corriente electromagnética de Noether---anula la tercera de forma unívoca ($C(q^2) \equiv 0$): por complicado que sea el interior del protón, no puede esconder más de \emph{dos} factores de forma, el de Dirac $F_1$ y el de Pauli $F_2$. La pieza algebraica que hace posible esa reducción es la Identidad de Gordon, que deducimos evaluando explícitamente cada conmutador espinorial en la capa de masa. A continuación resolvemos la cinemática de retroceso tipo Compton en el laboratorio y evaluamos las trazas de Dirac de hasta seis matrices gamma que definen los tensores leptónico $L_{\mu\nu}$ y hadrónico $W^{\mu\nu}$; su contracción contiene un término de interferencia $F_1 F_2$ que impide separar experimentalmente la carga de la magnetización. Ese obstáculo es precisamente lo que motiva el sistema de Breit: calculando los productos de espinores de Pauli en el marco donde la transferencia de energía se anula, las densidades estáticas se desacoplan y emergen los factores de Sachs ($G_E$, $G_M$), cuyos términos cruzados se cancelan de forma algebraicamente exacta. Integrando el espacio de fases relativista de dos cuerpos cerramos la demostración de la fórmula de Rosenbluth, verificamos que se reduce literalmente a la Ecuación (1) del artículo original de 1950 y formulamos la técnica de separación lineal que extrae $G_E^2$ del intercepto y $G_M^2$ de la pendiente de una recta. El resultado final llega en el último tramo: probamos que $G_E(Q^2)$ es la transformada de Fourier tridimensional de la densidad espacial estática de carga $\rho(\mathbf{r})$, que su pendiente en el origen fija unívocamente el radio cuadrático medio $\langle r_E^2 \rangle = -6 \left.\frac{dG_E}{dQ^2}\right|_0$, y que el modelo dipolar, invertido analíticamente mediante el teorema de residuos de Cauchy en el plano complejo, corresponde a una distribución exponencial de carga. Concluimos con el desenlace contemporáneo de esta historia: la técnica de transferencia de polarización ($P_t/P_l$) en Jefferson Lab, la discrepancia que reveló frente a la separación de Rosenbluth, su resolución mediante el intercambio de dos fotones (TPE) y el puzzle del radio del protón.
\par\vspace{1ex}
\textbf{Palabras clave:} Sección eficaz de Rosenbluth; Identidad de Gordon; Trazas de Dirac; Factores de Sachs; Marco de Breit; Transformada de Fourier; Teorema del residuo; Radio del protón; Transferencia de polarización; Intercambio de dos fotones (TPE).
\par\vspace{1.5ex}
\textbf{Abstract}---In 1933, Otto Stern measured the magnetic moment of the proton and obtained $\mu_p \approx 2.79\,\mu_N$, nearly three times the value that the Dirac equation imposes, with no adjustable freedom, on any point-like spin-$1/2$ fermion ($g = 2$). That anomaly was the first evidence that the proton occupies volume, and it raised a quantitative question: how does one turn a magnetic discrepancy into a measurement of the size of the proton, expressed in femtometres? This paper answers that question by reconstructing step by step, with no algebraic gaps, the full chain that connects quantum electrodynamics (QED) to a measurable number, taking as point of departure and guiding thread the elastic $e^- p \to e^- p$ differential cross section published by Marshall N. Rosenbluth in 1950. We first show that Lorentz invariance and parity conservation admit a covariant hadronic vertex built from three scalar functions, and that the Ward identity---the conservation of the Noether electromagnetic current---makes the third one vanish uniquely ($C(q^2) \equiv 0$): however complicated the interior of the proton may be, it cannot hide more than \emph{two} form factors, Dirac's $F_1$ and Pauli's $F_2$. The algebraic tool that makes this reduction possible is the Gordon Identity, which we derive by evaluating every spinor commutator explicitly on the mass shell. We then solve the Compton-like recoil kinematics in the laboratory frame and evaluate the Dirac traces of up to six gamma matrices defining the leptonic ($L_{\mu\nu}$) and hadronic ($W^{\mu\nu}$) tensors; their contraction contains an $F_1 F_2$ interference term that prevents any experimental separation of charge from magnetisation. That obstacle is exactly what motivates the Breit frame: by evaluating two-component Pauli spinor products in the frame where the energy transfer vanishes, the static densities decouple and the Sachs form factors ($G_E$, $G_M$) emerge, with their cross terms cancelling exactly. Integrating the relativistic two-body phase space completes the derivation of the Rosenbluth formula, which we verify reduces literally to Equation (1) of the original 1950 paper, and we formulate the linear separation technique that extracts $G_E^2$ from the intercept and $G_M^2$ from the slope of a straight line. The payoff comes in the final stretch: we prove that $G_E(Q^2)$ is the three-dimensional Fourier transform of the static spatial charge density $\rho(\mathbf{r})$, that its slope at the origin uniquely fixes the mean square charge radius $\langle r_E^2 \rangle = -6 \left.\frac{dG_E}{dQ^2}\right|_0$, and that the dipole model, inverted analytically via Cauchy's residue theorem in the complex plane, corresponds to an exponential charge distribution. We close with the contemporary denouement of this story: the polarization transfer technique ($P_t/P_l$) at Jefferson Lab, the discrepancy it revealed against the Rosenbluth separation, its resolution through two-photon exchange (TPE), and the proton radius puzzle.
\par\vspace{1ex}
\textbf{Keywords:} Rosenbluth cross section; Gordon Identity; Dirac traces; Sachs form factors; Breit frame; Fourier transform; Residue theorem; Proton radius; Polarization transfer; Two-photon exchange (TPE).
\end{quote}

%=======================================================================
%  CAJA DE CONVENCIONES Y NOTACIÓN
%=======================================================================
\vspace{0.5em}
\begin{center}
{\setlength{\fboxsep}{9pt}\setlength{\fboxrule}{0.6pt}%
\fcolorbox{colorborde}{colordest}{%
\begin{minipage}{0.955\textwidth}
\small
\begin{center}\textbf{\normalsize Convenciones y notación}\end{center}
\vspace{-0.4ex}
{\color{colorborde}\hrule height 0.4pt}
\vspace{1.0ex}

\textbf{Unidades naturales.} Trabajamos con $\hbar = c = 1$, de modo que energías, momentos y masas comparten unidad (GeV) y las longitudes son inversas de energía (GeV$^{-1}$). La conversión a unidades del laboratorio se hace con $\hbar c = 0.197327\;\text{GeV}\cdot\text{fm}$. La carga se mide mediante la constante de estructura fina $\alpha \equiv e^2/(4\pi) \approx 1/137.036$ (unidades de Heaviside-Lorentz). Designamos por $M$ la masa del protón ($938.272\;\text{MeV}$) y por $m_e$ la del electrón ($0.511\;\text{MeV}$), que se desprecia frente a $E$ salvo indicación explícita.

\vspace{0.8ex}
\textbf{Métrica y índices.} Signatura $g_{\mu\nu} = \mathrm{diag}(+1,-1,-1,-1)$, de manera que $a\cdot b = a^0 b^0 - \mathbf{a}\cdot\mathbf{b}$. Los índices griegos ($\mu, \nu, \alpha, \beta$) recorren $0,1,2,3$ y los latinos ($j,k,l$) las tres direcciones espaciales; los índices repetidos se suman (convenio de Einstein). Los trivectores se escriben en negrita ($\mathbf{q}$, $\boldsymbol{\sigma}$) y los cuadrivectores sin ella ($q^\mu$).

\vspace{0.8ex}
\textbf{Transferencia de momento: $q^2$ frente a $Q^2$.} El cuadrimomento que porta el fotón virtual es $q^\mu \equiv k^\mu - k'^\mu = p'^\mu - p^\mu$. En dispersión elástica $q$ es siempre de tipo espacial, es decir $q^2 \le 0$; por comodidad se trabaja con la cantidad positiva $Q^2 \equiv -q^2 \ge 0$. Ambas notaciones designan la misma magnitud física y los factores de forma se escriben indistintamente $F(q^2)$ o $F(Q^2)$ según convenga. El parámetro adimensional de retroceso es $\tau \equiv Q^2/(4M^2)$.

\vspace{0.8ex}
\textbf{Álgebra de Dirac.} Se adopta la representación estándar de Dirac-Pauli. Las matrices gamma satisfacen el álgebra de Clifford $\{\gamma^\mu,\gamma^\nu\} = 2g^{\mu\nu}\mathbb{I}$; se emplea la notación de Feynman $\slashed{a} \equiv a_\mu \gamma^\mu$, el tensor de espín $\sigma^{\mu\nu} \equiv \tfrac{i}{2}[\gamma^\mu,\gamma^\nu]$ y la matriz quiral $\gamma^5 \equiv i\gamma^0\gamma^1\gamma^2\gamma^3$. Las trazas se toman sobre el espacio espinorial de cuatro dimensiones, con $\mathrm{Tr}[\mathbb{I}] = 4$ y la convención $\epsilon^{0123} = +1$. Todas las identidades de contracción y de trazas utilizadas en el cuerpo del texto se demuestran desde cero en el Apéndice~\ref{sec:apendice_dirac}. Los espinores libres se normalizan covariantemente, $\bar{u}(p)u(p) = 2M$, con relación de completitud $\sum_s u(p,s)\bar{u}(p,s) = \slashed{p} + M$.

\vspace{0.8ex}
\textbf{Dos símbolos que se repiten con significados distintos.} El símbolo $\epsilon$ con índices ($\epsilon_{jkl}$, $\epsilon^{\mu\nu\alpha\beta}$) denota siempre el tensor de Levi-Civita; sin índices (Secciones~\ref{sec:seccion_rosenbluth} y~\ref{sec:radio_fourier}) denota la polarización longitudinal del fotón virtual. Análogamente, $\sigma^{\mu\nu}$ y $\boldsymbol{\sigma}$ son matrices (de Dirac y de Pauli respectivamente), mientras que $\sigma$ y $\sigma_{\text{red}}$ designan secciones eficaces.

\vspace{0.8ex}
\textbf{El símbolo $\dot{=}$.} En la Sección~\ref{sec:cinematica_trazas} escribimos $A \dot{=} B$ para indicar que $A$ y $B$ producen el mismo resultado \emph{al ser contraídos con el tensor leptónico} $L_{\mu\nu}$, que es transversal ($q^\mu L_{\mu\nu} = 0$). Es una igualdad válida bajo la contracción, no una identidad tensorial.

\vspace{0.8ex}
\textbf{Jerarquía de las ecuaciones destacadas.} Las ecuaciones \emph{enmarcadas} son los siete resultados que sostienen el hilo argumental completo del artículo; las de \emph{fondo sombreado} son los resultados importantes de cada tramo; el resto se numera de forma ordinaria. Esta distinción es deliberada: permite releer el documento siguiendo solo los marcos y recuperar el argumento entero.
\end{minipage}}}
\end{center}

\vspace{1.0em}
\section{Punto de Partida: La Sección Eficaz de Rosenbluth y la Estructura Hadrónica}
\label{sec:punto_de_partida}

Un protón dispersa electrones de manera cuantitativamente distinta a como lo haría una partícula puntual de la misma masa y del mismo espín. Todo este artículo es el desarrollo analítico y pedagógico de esa afirmación: qué se mide experimentalmente en el laboratorio, qué principios de simetría obligan a la teoría a admitir exactamente dos funciones de estructura para describirlo, y cómo esas dos funciones se traducen finalmente en un radio de carga medible en femtómetros.

Para guiar la intuición física desde el inicio, consideremos la escala espacial de resolución~\cite{das2003,henley2007}. De acuerdo con el principio de de Broglie, una sonda cuántica que transfiere un momento tridimensional $\mathbf{q}$ sondea el blanco con una longitud de onda reducida $\lambdabar = \hbar / |\mathbf{q}|$. En el régimen de baja energía ($|\mathbf{q}| \ll 200\text{ MeV}$, equivalente a transferencias $Q^2 \ll 0.05\text{ GeV}^2$), la longitud de onda de la sonda satisface $\lambdabar \gg 1\text{ fm}$, excediendo con creces la extensión espacial típica del nucleón. En dicho límite, el fotón virtual no puede resolver subestructuras espaciales y «ve» únicamente una carga neta puntual efectiva. En cambio, cuando la energía del proyectil se incrementa hasta el régimen ultra-relativista y el momento transferido alcanza $|\mathbf{q}| \gtrsim \hbar c / (0.84\text{ fm}) \approx 235\text{ MeV}$ ($Q^2 \gtrsim 0.1\text{ GeV}^2$), la resolución espacial de la sonda desciende por debajo de la dimensión física del protón ($\lambdabar \lesssim R_p$). En este régimen, las amplitudes dispersadas desde distintas regiones de la distribución interna de carga y corriente interfieren destructivamente, provocando una caída pronunciada de la sección eficaz con respecto a la predicción para un fermión puntual.

La estrategia expositiva que adoptamos en este trabajo sigue un esquema riguroso y transparente: presentamos en primer término el punto de llegada fenomenológico. Exponemos la fórmula deducida por Marshall N. Rosenbluth en 1950, la cual condensa toda la información físicamente accesible en la dispersión elástica $e^- p \to e^- p$. De este modo, se dispone desde el comienzo del mapa integral del territorio teórico; las secciones posteriores deducirán analíticamente cada componente de dicho mapa y, al concluir la Sección~\ref{sec:seccion_rosenbluth}, reconstruiremos la fórmula original término por término. La Figura~\ref{fig:puntual_vs_extendido} resume el contraste físico fundamental que motiva este estudio.

\begin{figure}[t]
\centering
\begin{tikzpicture}[>=stealth, line cap=round, font=\small]

%----------------------------------------------------------------------
% PANEL (a): fermión puntual de Dirac
%----------------------------------------------------------------------
\begin{scope}[shift={(0,0)}]
  % equipotenciales coulombianas
  \foreach \r in {0.50,0.90,1.30,1.70}
     \draw[gray!40, dashed, line width=0.3pt] (0,0) circle (\r);
  % líneas de campo radiales
  \foreach \a in {0,30,...,330}
     \draw[gray!60, line width=0.35pt, ->] (\a:0.16) -- (\a:1.82);
  % la carga puntual
  \fill[black] (0,0) circle (2.0pt);
  \node[font=\footnotesize] at (0,-2.18) {sin extensión espacial: $r < 10^{-19}\,\text{m}$};
  \node[font=\small\bfseries] at (0,2.35) {(a) Fermión puntual de Dirac};
\end{scope}

%----------------------------------------------------------------------
% PANEL (b): protón extendido
%----------------------------------------------------------------------
\begin{scope}[shift={(7.4,0)}]
  % distribución de carga difusa (sombreado radial)
  \shade[inner color=blue!42, outer color=blue!2] (0,0) circle (1.80);
  % contorno del radio de carga
  \draw[blue!55, dashed, line width=0.5pt] (0,0) circle (1.25);
  \draw[blue!70, line width=0.5pt, <->] (0,0) -- (-42:1.25);
  \node[blue!70, font=\footnotesize, anchor=west, inner sep=2pt] at (-38:1.32) {$r_p \!\approx\! 0.84\,\text{fm}$};
  % quarks de valencia
  \fill[red!75!black]   (-0.62, 0.42) circle (2.6pt) node[above left,  font=\scriptsize, inner sep=2.5pt] {$u$};
  \fill[red!75!black]   ( 0.18, 0.78) circle (2.6pt) node[above right, font=\scriptsize, inner sep=2.5pt] {$u$};
  \fill[green!45!black] (-0.10,-0.55) circle (2.6pt) node[below right, font=\scriptsize, inner sep=2.5pt] {$d$};
  \node[font=\footnotesize] at (0,-2.18) {densidad de carga y de magnetización $\rho(\mathbf{r})$};
  \node[font=\small\bfseries] at (0,2.35) {(b) Protón real: distribución extendida};
\end{scope}

%----------------------------------------------------------------------
% SONDA: fotón virtual incidente sobre cada panel
%----------------------------------------------------------------------
\draw[decorate, decoration={coil, aspect=0, segment length=4.6pt, amplitude=2.6pt},
      line width=0.6pt, gray!85] (-3.30,-1.55) -- (-1.05,-0.62);
\draw[->, gray!85, line width=0.6pt] (-1.05,-0.62) -- (-0.55,-0.35);
\node[gray!85, font=\footnotesize] at (-3.42,-1.78) {$\gamma^*$};

\draw[decorate, decoration={coil, aspect=0, segment length=4.6pt, amplitude=2.6pt},
      line width=0.6pt, gray!85] (4.10,-1.55) -- (6.35,-0.62);
\draw[->, gray!85, line width=0.6pt] (6.35,-0.62) -- (6.85,-0.35);
\node[gray!85, font=\footnotesize] at (3.98,-1.78) {$\gamma^*$};

%----------------------------------------------------------------------
% RÓTULOS COMPARATIVOS INFERIORES
%----------------------------------------------------------------------
\node[align=center, font=\footnotesize] at (0,-3.15)
  {$g = 2$ \; (exacto) \qquad $\mu = \dfrac{e}{2M}$\\[0.5ex]
   $F_1 \equiv 1$, \; $F_2 \equiv 0$\\[0.3ex]
   $G_E(Q^2) = 1$ \; para todo $Q^2$};

\node[align=center, font=\footnotesize] at (7.4,-3.15)
  {$g_p \approx 5.586$ \qquad $\mu_p = 2.7928\,\mu_N$\\[0.5ex]
   $F_1(Q^2)$, \; $F_2(Q^2)$ \; no triviales\\[0.3ex]
   $G_E(Q^2)$ \; decae al crecer $Q^2$};

% separador vertical
\draw[gray!35, line width=0.4pt] (3.70,-3.75) -- (3.70,2.55);

\end{tikzpicture}
\caption{Contraste físico fundamental entre un blanco puntual y un nucleón extendido. \textbf{(a)} Un fermión elemental puntual de Dirac posee un campo coulombiano puro originado en una singularidad puntual sin estructura ($r < 10^{-19}\text{ m}$), su factor giromagnético es exactamente $g=2$ y su respuesta electromagnética no depende de la resolución del fotón virtual $\gamma^*$ ($F_1 \equiv 1$, $F_2 \equiv 0$, $G_E = G_M = 1$). \textbf{(b)} El protón real es un estado ligado no perturbativo de quarks y gluones en QCD, con una distribución finita de carga y magnetización $\rho(\mathbf{r})$ con radio $r_p \approx 0.84\text{ fm}$. El fotón virtual con momento transferido $Q^2$ resuelve distancias espaciales del orden de $\lambdabar \sim 1/\sqrt{Q^2}$, generando interferencias destructivas que atenúan fuertemente la dispersión respecto a la predicción puntual. Esta atenuación queda rigurosamente codificada en los factores de forma $G_E(Q^2)$ y $G_M(Q^2)$.}
\label{fig:puntual_vs_extendido}
\end{figure}

\subsection{La Fórmula Magistral de Rosenbluth y la Anatomía de sus Términos}
Considérese el proceso fundamental de dispersión elástica de un electrón ultra-relativista sobre un protón en el sistema de referencia del laboratorio, donde el protón actúa como blanco inicialmente en reposo:
\begin{equation}
e^-(k) + p(p) \longrightarrow e^-(k') + p(p').
\label{eq:proceso_ep}
\end{equation}
En agosto de 1950, Marshall N. Rosenbluth publicó en \textit{Physical Review} el cálculo analítico que transformó el estudio de la estructura hadrónica \cite{rosenbluth1950}. En su artículo seminal, Rosenbluth presentó la \textbf{fórmula analítica completa y cerrada} para la sección eficaz elástica diferencial en el laboratorio:
\begin{equation}
\boxed{
\begin{aligned}
\sigma \, d\Omega &= \left(\frac{e' e''}{2E}\right)^2 \cot^2\left(\frac{\theta}{2}\right) \csc^2\left(\frac{\theta}{2}\right) \\
&\quad \times \left\{ \frac{1 + 2\left(\dfrac{E}{M}\right)\sin^2\left(\dfrac{\theta}{2}\right) + \left(\dfrac{E^2}{M^2}\right)\left[ 2(1+\kappa')^2 \tan^2\left(\dfrac{\theta}{2}\right)\sin^2\left(\dfrac{\theta}{2}\right) + {\kappa'}^2 \sin^2\left(\dfrac{\theta}{2}\right) \right]}{\left[1 + 2\left(\dfrac{E}{M}\right)\sin^2\left(\dfrac{\theta}{2}\right)\right]^2} \right\} d\Omega
\end{aligned}
}
\label{eq:rosenbluth_1950_original}
\end{equation}
donde $E$ denota la energía incidente del proyectil leptónico, $M$ la masa en reposo del protón, $\theta$ el ángulo de dispersión en el laboratorio, $e''$ la carga efectiva del electrón que incorpora correcciones radiativas calculadas por Julian Schwinger \cite{schwinger1949}, $e'$ la carga efectiva del protón y $\kappa'$ su momento magnético anómalo efectivo inducido por su estructura interna. Conviene advertir que esta notación histórica de 1950 prescindía de la noción moderna de factores de forma dependientes de $Q^2$, expresando la estructura mediante $e'$ y $\kappa'$, que hoy identificamos con $F_1(Q^2)$ y $F_2(Q^2)$. La reproducción exacta de esta fórmula constituye uno de los objetivos de control de nuestro desarrollo: al finalizar la Sección~\ref{sec:seccion_rosenbluth} volveremos sobre la Ec.~\eqref{eq:rosenbluth_1950_original} y comprobaremos que el formalismo covariante moderno se reduce rigurosamente a ella término por término.

En el marco moderno de la electrodinámica cuántica y la física hadrónica, desplegando todos los factores cinemáticos y sustituyendo el factor de Mott sin abreviaciones, la \textbf{Fórmula Completa y Explícita de Rosenbluth} expresada en términos de los factores de forma de Sachs $G_E(Q^2)$ y $G_M(Q^2)$ adopta la forma canónica:
\begin{equation}
\dest{
\left(\frac{d\sigma}{d\Omega}\right) = \frac{\alpha^2 \cos^2\left(\dfrac{\theta}{2}\right)}{4 E^2 \sin^4\left(\dfrac{\theta}{2}\right) \left[ 1 + \dfrac{2E}{M}\sin^2\left(\dfrac{\theta}{2}\right) \right]} \left[ \frac{G_E^2(Q^2) + \tau G_M^2(Q^2)}{1+\tau} + 2\tau G_M^2(Q^2) \tan^2\left(\frac{\theta}{2}\right) \right]
}
\label{eq:rosenbluth_master_intro}
\end{equation}
o equivalentemente, en la representación covariante canónica de los factores de forma de Dirac $F_1(Q^2)$ y Pauli $F_2(Q^2)$:
\begin{equation}
\dest{
\begin{aligned}
\left(\frac{d\sigma}{d\Omega}\right) = \frac{\alpha^2 \cos^2\left(\dfrac{\theta}{2}\right)}{4 E^2 \sin^4\left(\dfrac{\theta}{2}\right) \left[ 1 + \dfrac{2E}{M}\sin^2\left(\dfrac{\theta}{2}\right) \right]} &\left\{ F_1^2(Q^2) + \tau \left[ F_2^2(Q^2) \right.\right. \\
&\quad \left.\left. +\, 2\Big( F_1(Q^2) + F_2(Q^2) \Big)^2 \tan^2\left(\frac{\theta}{2}\right) \right] \right\}
\end{aligned}
}
\label{eq:rosenbluth_dirac_pauli_intro}
\end{equation}
donde el parámetro adimensional de retroceso relativista $\tau \equiv Q^2/(4M^2)$ se escribe explícitamente en función de las variables de laboratorio como:
\begin{equation}
\tau = \frac{E^2 \sin^2(\theta/2)}{M^2 \left[1 + \dfrac{2E}{M}\sin^2(\theta/2)\right]}.
\label{eq:tau_explicit_intro}
\end{equation}

Para comprender la arquitectura física de la Ec.~\eqref{eq:rosenbluth_master_intro}, conviene desglosar minuciosamente cada uno de sus componentes:
\begin{enumerate}
    \item \textbf{La sección eficaz de Mott con retroceso relativista:} El factor cinemático base está dado por:
    \begin{equation}
    \left(\frac{d\sigma}{d\Omega}\right)_{\text{Mott}} = \frac{\alpha^2}{4 E^2 \sin^4(\theta/2)} \left(\frac{E'}{E}\right) \cos^2\left(\frac{\theta}{2}\right),
    \label{eq:mott_recoil_intro}
    \end{equation}
    donde $\alpha \equiv e^2/(4\pi) \approx 1/137.036$ es la constante de estructura fina de QED, $E$ y $E'$ son las energías incidente y dispersada del electrón, y $\theta$ es el ángulo de dispersión en el laboratorio. Este factor describe la dispersión elástica sobre un blanco puntual de espín 0 o la línea de base para un fermión sin estructura interna ($r < 10^{-19}\text{ m}$) con factor giromagnético puro $g=2$.
    \item \textbf{El parámetro adimensional de retroceso cinemático ($\tau$):} Definido formalmente como:
    \begin{equation}
    \tau \equiv \frac{Q^2}{4M^2},
    \label{eq:tau_def_intro}
    \end{equation}
    donde $Q^2 \equiv -q^2 > 0$ es el cuadrimomento transferido de tipo espacial y $M$ es la masa en reposo del protón. El parámetro $\tau$ cuantifica la importancia de los efectos relativistas de retroceso en escalas de la masa nucleónica; en el régimen no relativista ($\tau \ll 1$), el retroceso es insignificante, mientras que para $\tau \sim 1$ las interacciones dinámicas de espín y transferencia magnética dominan cuantitativamente la sección eficaz.
    \item \textbf{Los factores de forma de Sachs ($G_E$ y $G_M$):} Las funciones escalares $G_E(Q^2)$ (\textbf{factor de forma eléctrico}) y $G_M(Q^2)$ (\textbf{factor de forma magnético}) parametrizan el desconocimiento microscópico de la distribución espacial interna de carga eléctrica y de densidad magnetizante del nucleón. Son las dos únicas incógnitas de la ecuación, y la teoría demuestra que no pueden existir más funciones independientes.
    \item \textbf{La dispersión magnética pura ($\tan^2(\theta/2)$):} A ángulos traseros extremos ($\theta \to 180^\circ$, donde $\cos(\theta/2) \to 0$ y $\tan(\theta/2) \to \infty$), el acoplamiento electrostático se suprime cinemáticamente por conservación de helicidad y la dispersión ocurre exclusivamente a través de la interacción magnética dipolar con volteo de espín (\textit{spin-flip}) del nucleón, gobernada por $G_M^2(Q^2)$.
    \item \textbf{El desacoplamiento cuadrático:} En la Ec.~\eqref{eq:rosenbluth_master_intro} no existe ningún término de interferencia cruzado del tipo $G_E(Q^2) G_M(Q^2)$. Ambos observables aparecen estrictamente desacoplados como combinaciones cuadráticas independientes ($G_E^2$ y $G_M^2$), permitiendo una separación experimental limpia mediante la técnica de regresión lineal ideada por Rosenbluth. Esta propiedad algebraicamente notable es consecuencia directa de formular la interacción en el sistema de referencia de Breit, como demostraremos rigurosamente en la Sección~\ref{sec:marco_breit}.
\end{enumerate}

Para situar históricamente y metodológicamente la fórmula de Rosenbluth dentro del andamiaje de la física de colisiones, la Tabla~\ref{tab:modelos_dispersion} presenta un despiece comparativo paralelo de los cuatro regímenes canónicos de dispersión electromagnética, desde la mecánica clásica no relativista hasta la electrodinámica cuántica con estructura hadrónica.

\begin{table}[t]
\centering
\small
\renewcommand{\arraystretch}{1.3}
\caption{Comparativa sistemática de los modelos de dispersión electromagnética elástica. La progresión muestra cómo cada avance teórico introduce sucesivamente el espín del proyectil, el retroceso relativista del blanco y, finalmente, la estructura espacial interna del nucleón codificada en los factores de forma de Rosenbluth.}
\label{tab:modelos_dispersion}
\resizebox{0.98\textwidth}{!}{%
\begin{tabular}{l c c c p{7.0cm} p{4.6cm}}
\toprule
\textbf{Modelo / Régimen} & \textbf{Espín $s_e / s_p$} & \textbf{Estructura Blanco} & \textbf{Cinemática} & \textbf{Sección Eficaz $\dfrac{d\sigma}{d\Omega}$} & \textbf{Física Clave / Observación} \\
\midrule
\textbf{Rutherford (1911)}~\cite{rutherford1911} & $0 \;/\; 0$ & Puntual ($r=0$) & Estática ($M \to \infty$) &
$\displaystyle \left(\frac{\alpha}{4E}\right)^2 \frac{1}{\sin^4(\theta/2)}$ &
Dispersión coulombiana clásica pura; divergencia hacia adelante ($\theta \to 0$). \\
\addlinespace[0.6ex]
\textbf{Mott (1929)}~\cite{mott1929} & $\frac{1}{2} \;/\; 0$ & Puntual ($r=0$) & Estática ($M \to \infty$) &
$\displaystyle \left(\frac{d\sigma}{d\Omega}\right)_{\text{Ruth}} \cos^2\left(\frac{\theta}{2}\right)$ &
Electrón de Dirac relativista; conservación de helicidad anula la dispersión a $\theta = 180^\circ$. \\
\addlinespace[0.6ex]
\textbf{Dirac con Retroceso}~\cite{dirac1928} & $\frac{1}{2} \;/\; \frac{1}{2}$ & Puntual ($g=2$) & Relativista ($E'/E$) &
$\displaystyle \left(\frac{d\sigma}{d\Omega}\right)_{\text{Mott}} \left[ 1 + 2\tau\tan^2\left(\frac{\theta}{2}\right) \right]$ &
Blanco fermiónico puntual con espín; el término $2\tau\tan^2(\theta/2)$ surge del acoplamiento magnético normal. \\
\addlinespace[0.6ex]
\textbf{Rosenbluth (1950)}~\cite{rosenbluth1950} & $\frac{1}{2} \;/\; \frac{1}{2}$ & Extendido ($g \neq 2$) & Relativista ($E'/E$) &
$\displaystyle \left(\frac{d\sigma}{d\Omega}\right)_{\text{Mott}} \left[ \frac{G_E^2 + \tau G_M^2}{1+\tau} + 2\tau G_M^2\tan^2\left(\frac{\theta}{2}\right) \right]$ &
Blanco hadrónico con estructura interna; separación cuadrática de Sachs sin interferencias cruzadas $G_E G_M$. \\
\bottomrule
\end{tabular}%
}%
\end{table}

La presencia ineludible de las funciones $G_E(Q^2)$ y $G_M(Q^2)$ en la Ec.~\eqref{eq:rosenbluth_master_intro} plantea una pregunta física primaria: ¿por qué la dispersión sobre el protón no puede describirse mediante el factor de Mott puntual puro o el modelo de Dirac puntual con retroceso? La respuesta no es una conjetura teórica, sino un hecho experimental contundente anterior en casi dos décadas a la fórmula de Rosenbluth, y a su deducción dedicamos las subsecciones inmediatas.

\subsection{La Ruptura de la Puntualidad: Del Teorema de Dirac (\texorpdfstring{$g=2$}{g=2}) a la Anomalía de Stern}
Antes de aceptar que el protón tiene estructura interna es preciso saber, con toda precisión y sin ningún margen de ajuste, qué predeciría la teoría si no la tuviera. Esa es la función de la presente subsección. No introducimos la ecuación de Dirac por completitud formal ni por respeto histórico: la introducimos porque de ella se desprende una predicción numérica cerrada---el factor giromagnético $g = 2$---que no contiene un solo parámetro libre y que, en consecuencia, puede ser refutada por un único experimento. El cálculo que sigue consiste en reducir la ecuación de Dirac a su límite no relativista y comprobar que el acoplamiento magnético del espín aparece \emph{solo}, sin haber sido postulado en ningún punto.

Al lector que encuentra aquí por primera vez la ecuación de Dirac le bastan tres hechos para seguir el desarrollo sin pérdida. Primero: se trata de la ecuación de onda relativista de una partícula de espín $1/2$, es decir, el análogo de la ecuación de Schrödinger cuando se exige compatibilidad estricta con la relatividad especial. Segundo: su incógnita $\psi$ no es una función escalar sino un objeto de cuatro componentes (un biespinor), porque describe simultáneamente los dos estados de espín de la partícula y los dos de su antipartícula. Tercero: los símbolos $\gamma^\mu$ son cuatro matrices $4\times 4$ fijas, y la única propiedad suya que usaremos en esta subsección es su forma explícita en bloques $2\times 2$; el Apéndice~\ref{sec:apendice_dirac} demuestra desde cero todas las identidades algebraicas que necesitaremos más adelante. El punto conceptual que conviene retener es que el espín no se añade a mano en esta formulación: viene incorporado en la estructura matricial de la ecuación, y el momento magnético que estamos a punto de deducir es una consecuencia inevitable de esa estructura.

En electrodinámica cuántica relativista~\cite{dirac1928,griffiths2008,thomson2013}, un fermión elemental puntual de espín $1/2$, masa $M$ y carga $e$, interactúa con el cuadripotencial electromagnético $A^\mu = (\Phi, \mathbf{A})$ según la ecuación covariante de Dirac con acoplamiento mínimo:
\begin{equation}
\left[ \gamma^\mu (i\partial_\mu - e A_\mu) - M \right] \psi = 0.
\label{eq:dirac_min_coupling}
\end{equation}
En la métrica de Minkowski $\eta_{\mu\nu} = \text{diag}(+1, -1, -1, -1)$, el cuadrivector derivada y el cuadripotencial se escriben como $\partial_\mu = (\partial_0, \boldsymbol{\nabla})$ y $A_\mu = (\Phi, -\mathbf{A})$. Por tanto, el operador diferencial covariante se descompone explícitamente en su parte temporal y espacial:
\begin{equation}
\gamma^\mu (i\partial_\mu - e A_\mu) = \gamma^0 (i\partial_0 - e\Phi) - \boldsymbol{\gamma} \cdot (-i\boldsymbol{\nabla} - e\mathbf{A}) = \gamma^0 (i\partial_0 - e\Phi) - \boldsymbol{\gamma} \cdot (\mathbf{p} - e\mathbf{A}),
\end{equation}
donde hemos identificado el operador de momento lineal $\mathbf{p} \equiv -i\boldsymbol{\nabla}$.

En la representación estándar de Dirac-Pauli, las matrices $\gamma^\mu$ se definen en bloques de $2 \times 2$ como:
\begin{equation}
\gamma^0 = \begin{pmatrix} \mathbb{I} & 0 \\ 0 & -\mathbb{I} \end{pmatrix}, \qquad \boldsymbol{\gamma} = \begin{pmatrix} 0 & \boldsymbol{\sigma} \\ -\boldsymbol{\sigma} & 0 \end{pmatrix},
\end{equation}
donde $\mathbb{I} = \begin{pmatrix} 1 & 0 \\ 0 & 1 \end{pmatrix}$ es la matriz identidad $2\times 2$ y $\boldsymbol{\sigma} = (\sigma_1, \sigma_2, \sigma_3)$ son las matrices hermíticas de Pauli:
\begin{equation}
\sigma_1 = \begin{pmatrix} 0 & 1 \\ 1 & 0 \end{pmatrix}, \qquad \sigma_2 = \begin{pmatrix} 0 & -i \\ i & 0 \end{pmatrix}, \qquad \sigma_3 = \begin{pmatrix} 1 & 0 \\ 0 & -1 \end{pmatrix}.
\end{equation}
La estrategia del cálculo es la siguiente: al ser $\psi$ un objeto de cuatro componentes, la Ec.~\eqref{eq:dirac_min_coupling} equivale a un sistema de dos ecuaciones acopladas para dos espinores de dos componentes cada uno. En el régimen no relativista una de esas dos piezas resulta ser mucho menor que la otra, lo que permite eliminarla y quedarse con una única ecuación de Schrödinger efectiva. Todo el contenido magnético queda entonces al descubierto. Descomponiendo el bioespinor de cuatro componentes en dos espinores de Pauli de dos componentes asociados a un estado estacionario de energía total $E$:
\begin{equation}
\psi(\mathbf{x}, t) = e^{-i E t} \begin{pmatrix} \phi(\mathbf{x}) \\ \chi(\mathbf{x}) \end{pmatrix},
\end{equation}
el operador temporal actúa directamente como $i\partial_0 \psi = i(-iE)\psi = E\psi$. Sustituyendo en la Ec.~\eqref{eq:dirac_min_coupling}:
\begin{equation}
\left[ \begin{pmatrix} \mathbb{I} & 0 \\ 0 & -\mathbb{I} \end{pmatrix} (E - e\Phi) - \begin{pmatrix} 0 & \boldsymbol{\sigma} \\ -\boldsymbol{\sigma} & 0 \end{pmatrix} \cdot (\mathbf{p} - e\mathbf{A}) - M \begin{pmatrix} \mathbb{I} & 0 \\ 0 & \mathbb{I} \end{pmatrix} \right] \begin{pmatrix} \phi \\ \chi \end{pmatrix} = \begin{pmatrix} 0 \\ 0 \end{pmatrix}.
\end{equation}
Agrupando las matrices en un solo bloque de $2 \times 2$:
\begin{equation}
\begin{pmatrix} (E - e\Phi - M)\mathbb{I} & -\boldsymbol{\sigma} \cdot (\mathbf{p} - e\mathbf{A}) \\ \boldsymbol{\sigma} \cdot (\mathbf{p} - e\mathbf{A}) & -(E - e\Phi + M)\mathbb{I} \end{pmatrix} \begin{pmatrix} \phi \\ \chi \end{pmatrix} = \begin{pmatrix} 0 \\ 0 \end{pmatrix}.
\end{equation}
Efectuando la multiplicación matricial bloque a bloque, obtenemos el sistema lineal acoplado exacto:
\begin{align}
(E - e\Phi - M)\phi - \boldsymbol{\sigma} \cdot (\mathbf{p} - e\mathbf{A}) \chi &= 0, \label{eq:spinor_phi} \\
\boldsymbol{\sigma} \cdot (\mathbf{p} - e\mathbf{A}) \phi - (E - e\Phi + M)\chi &= 0. \label{eq:spinor_chi}
\end{align}

Hasta aquí no se ha hecho ninguna aproximación: las Ecs.~\eqref{eq:spinor_phi} y \eqref{eq:spinor_chi} son exactamente equivalentes a la ecuación de Dirac de partida. La aproximación entra ahora, y es la única de todo el cálculo. En el régimen no relativista, la energía total está dominada por la masa en reposo del fermión:
\begin{equation}
E = M + E_{\text{NR}}, \qquad \text{con } E_{\text{NR}} \ll M \quad \text{y} \quad |e\Phi| \ll M.
\end{equation}
Sustituyendo $E$ en el coeficiente de la Ec.~\eqref{eq:spinor_chi}:
\begin{equation}
E - e\Phi + M = M + E_{\text{NR}} - e\Phi + M = 2M + (E_{\text{NR}} - e\Phi) \approx 2M.
\end{equation}
Por tanto, la Ec.~\eqref{eq:spinor_chi} permite despejar rigurosamente el espinor inferior (pequeño) $\chi$ en términos del espinor superior (grande) $\phi$:
\begin{equation}
2M \chi \approx \boldsymbol{\sigma} \cdot (\mathbf{p} - e\mathbf{A}) \phi \implies \chi \approx \frac{\boldsymbol{\sigma} \cdot (\mathbf{p} - e\mathbf{A})}{2M} \phi.
\label{eq:chi_approx}
\end{equation}
Sustituyendo la Ec.~\eqref{eq:chi_approx} en la primera ecuación del sistema, Ec.~\eqref{eq:spinor_phi}, y notando que $E - M = E_{\text{NR}}$:
\begin{equation}
(E_{\text{NR}} - e\Phi)\phi - [\boldsymbol{\sigma} \cdot (\mathbf{p} - e\mathbf{A})] \left[ \frac{\boldsymbol{\sigma} \cdot (\mathbf{p} - e\mathbf{A})}{2M} \phi \right] = 0,
\end{equation}
lo que conduce a:
\begin{equation}
(E_{\text{NR}} - e\Phi)\phi = \frac{1}{2M} \left[ \boldsymbol{\sigma} \cdot (\mathbf{p} - e\mathbf{A}) \right] \left[ \boldsymbol{\sigma} \cdot (\mathbf{p} - e\mathbf{A}) \right] \phi.
\label{eq:pauli_eq_step1}
\end{equation}

Nos queda ya una sola ecuación, con la forma de una ecuación de Schrödinger. Todo depende ahora de evaluar el miembro derecho, que contiene el producto de dos operadores idénticos del tipo $\boldsymbol{\sigma}\cdot\boldsymbol{\pi}$. Podría suponerse que ese producto es simplemente $\boldsymbol{\pi}^2$, es decir, la energía cinética habitual; no lo es, y la diferencia entre ambos es exactamente el término magnético que estamos buscando. Para calcular analíticamente el miembro derecho, demostremos primero la \textbf{Identidad Vectorial Fundamental de Pauli} para dos vectores de operadores arbitrarios $\mathbf{a}$ y $\mathbf{b}$. Las matrices de Pauli satisfacen el conmutador y el anticonmutador:
\begin{equation}
[\sigma_j, \sigma_k] = 2i \epsilon_{jkl} \sigma_l, \qquad \{\sigma_j, \sigma_k\} = 2\delta_{jk} \mathbb{I}.
\end{equation}
Sumando ambas relaciones miembro a miembro:
\begin{equation}
\sigma_j \sigma_k + \sigma_k \sigma_j + \sigma_j \sigma_k - \sigma_k \sigma_j = 2\delta_{jk}\mathbb{I} + 2i\epsilon_{jkl}\sigma_l \implies 2\sigma_j \sigma_k = 2(\delta_{jk}\mathbb{I} + i\epsilon_{jkl}\sigma_l).
\end{equation}
Dividiendo entre 2, obtenemos el producto de dos matrices de Pauli cualesquiera:
\begin{equation}
\sigma_j \sigma_k = \delta_{jk}\mathbb{I} + i\epsilon_{jkl}\sigma_l.
\label{eq:pauli_product_rule}
\end{equation}
Contrayendo la Ec.~\eqref{eq:pauli_product_rule} con las componentes $a_j b_k$:
\begin{align}
(\boldsymbol{\sigma} \cdot \mathbf{a})(\boldsymbol{\sigma} \cdot \mathbf{b}) &= (\sigma_j a_j)(\sigma_k b_k) = \sigma_j \sigma_k a_j b_k \nonumber \\
&= (\delta_{jk}\mathbb{I} + i\epsilon_{jkl}\sigma_l) a_j b_k = (\delta_{jk} a_j b_k)\mathbb{I} + i \sigma_l (\epsilon_{jkl} a_j b_k) \nonumber \\
&= (\mathbf{a} \cdot \mathbf{b})\mathbb{I} + i \sigma_l (\mathbf{a} \times \mathbf{b})_l = (\mathbf{a} \cdot \mathbf{b})\mathbb{I} + i \boldsymbol{\sigma} \cdot (\mathbf{a} \times \mathbf{b}).
\label{eq:pauli_vector_identity}
\end{align}
El primer sumando es el producto escalar esperado; el segundo, proporcional al producto vectorial, es el que no tiene análogo clásico y el que va a generar el magnetismo.

Definiendo el operador momento cinemático gauge-invariante $\boldsymbol{\pi} \equiv \mathbf{p} - e\mathbf{A}$, aplicamos la Ec.~\eqref{eq:pauli_vector_identity} haciendo $\mathbf{a} = \mathbf{b} = \boldsymbol{\pi}$:
\begin{equation}
(\boldsymbol{\sigma} \cdot \boldsymbol{\pi})^2 = \boldsymbol{\pi}^2 \, \mathbb{I} + i \boldsymbol{\sigma} \cdot (\boldsymbol{\pi} \times \boldsymbol{\pi}).
\label{eq:pi_squared}
\end{equation}
Para un vector ordinario, el producto vectorial de un vector consigo mismo se anula idénticamente y el segundo término desaparecería. Aquí no ocurre así, porque las componentes de $\boldsymbol{\pi}$ son operadores que no conmutan entre sí en presencia de un campo magnético. Evaluemos entonces minuciosamente la componente $i$-ésima del producto vectorial $(\boldsymbol{\pi} \times \boldsymbol{\pi})$:
\begin{align}
(\boldsymbol{\pi} \times \boldsymbol{\pi})_i &= \epsilon_{ijk} \pi_j \pi_k = \epsilon_{ijk} (p_j - e A_j)(p_k - e A_k) \nonumber \\
&= \epsilon_{ijk} \left[ p_j p_k - e p_j A_k - e A_j p_k + e^2 A_j A_k \right].
\end{align}
Dado que las derivadas espaciales conmutan ($p_j p_k = p_k p_j$) y las componentes del potencial clásico conmutan ($A_j A_k = A_k A_j$), sus contracciones con el tensor completamente antisimétrico de Levi-Civita se anulan idénticamente:
\begin{equation}
\epsilon_{ijk} p_j p_k = 0, \qquad \epsilon_{ijk} A_j A_k = 0.
\end{equation}
Por consiguiente:
\begin{equation}
(\boldsymbol{\pi} \times \boldsymbol{\pi})_i = -e \epsilon_{ijk} (p_j A_k + A_j p_k).
\end{equation}
Reetiquetando los índices mudos en el segundo término ($j \leftrightarrow k$) y utilizando la propiedad antisimétrica $\epsilon_{ikj} = -\epsilon_{ijk}$:
\begin{equation}
-e \epsilon_{ijk} A_j p_k = -e \epsilon_{ikj} A_k p_j = +e \epsilon_{ijk} A_k p_j.
\end{equation}
Sumando ambos términos:
\begin{equation}
(\boldsymbol{\pi} \times \boldsymbol{\pi})_i = -e \epsilon_{ijk} (p_j A_k - A_k p_j) = -e \epsilon_{ijk} [p_j, A_k].
\end{equation}
Para calcular el conmutador cuántico $[p_j, A_k]$, lo aplicamos sobre una función de onda de prueba $\phi(\mathbf{x})$ utilizando la regla de derivación de Leibniz del producto:
\begin{align}
[p_j, A_k]\phi &= p_j(A_k \phi) - A_k(p_j \phi) = -i\partial_j(A_k \phi) - A_k(-i\partial_j \phi) \nonumber \\
&= -i (\partial_j A_k)\phi - i A_k (\partial_j \phi) + i A_k (\partial_j \phi) = -i (\partial_j A_k)\phi.
\end{align}
Dado que la función de prueba $\phi$ es arbitraria, el conmutador es:
\begin{equation}
[p_j, A_k] = -i \partial_j A_k = -i \frac{\partial A_k}{\partial x^j}.
\end{equation}
Sustituyendo este conmutador:
\begin{equation}
(\boldsymbol{\pi} \times \boldsymbol{\pi})_i = -e \epsilon_{ijk} (-i \partial_j A_k) = i e \epsilon_{ijk} \partial_j A_k = i e (\boldsymbol{\nabla} \times \mathbf{A})_i.
\end{equation}
Recordando que el campo de inducción magnética es $\mathbf{B} = \boldsymbol{\nabla} \times \mathbf{A}$, obtenemos con exactitud matemática:
\begin{equation}
\boldsymbol{\pi} \times \boldsymbol{\pi} = i e \mathbf{B}.
\end{equation}
Este es el paso decisivo de toda la subsección, y conviene detenerse en él: el campo magnético ha aparecido no porque lo hayamos introducido, sino porque las componentes del momento cinemático dejan de conmutar cuando hay un potencial vectorial. De esa falta de conmutatividad---y de nada más---va a salir el acoplamiento del espín al campo. Sustituyendo este conmutador en la Ec.~\eqref{eq:pi_squared}:
\begin{equation}
(\boldsymbol{\sigma} \cdot \boldsymbol{\pi})^2 = (\mathbf{p} - e\mathbf{A})^2 + i \boldsymbol{\sigma} \cdot (i e \mathbf{B}) = (\mathbf{p} - e\mathbf{A})^2 + i^2 e (\boldsymbol{\sigma} \cdot \mathbf{B}) = (\mathbf{p} - e\mathbf{A})^2 - e \boldsymbol{\sigma} \cdot \mathbf{B}.
\label{eq:pauli_identity_expanded}
\end{equation}
Insertando la Ec.~\eqref{eq:pauli_identity_expanded} en la Ec.~\eqref{eq:pauli_eq_step1}:
\begin{equation}
(E_{\text{NR}} - e\Phi)\phi = \frac{1}{2M} \left[ (\mathbf{p} - e\mathbf{A})^2 - e \boldsymbol{\sigma} \cdot \mathbf{B} \right] \phi.
\end{equation}
Despejando $E_{\text{NR}}\phi$, arribamos a la ecuación cuántica estacionaria de Pauli-Schrödinger:
\begin{equation}
E_{\text{NR}} \phi = \left[ \frac{(\mathbf{p} - e\mathbf{A})^2}{2M} + e\Phi - \frac{e}{2M} \boldsymbol{\sigma} \cdot \mathbf{B} \right] \phi.
\label{eq:pauli_schrodinger_final}
\end{equation}
Los dos primeros términos del corchete son la energía cinética y la energía potencial electrostática que ya aparecerían en la mecánica cuántica no relativista ordinaria. El tercero es nuevo, y no fue introducido en ningún momento: es un residuo inevitable de haber partido de una ecuación relativista.

El hamiltoniano de interacción magnética es el término de acoplamiento Zeeman:
\begin{equation}
H_{\text{mag}} = - \frac{e}{2M} \boldsymbol{\sigma} \cdot \mathbf{B}.
\end{equation}
En mecánica cuántica, la energía potencial de interacción de un dipolo magnético $\boldsymbol{\mu}$ inmerso en un campo magnético externo $\mathbf{B}$ se define por $H_{\text{mag}} = - \boldsymbol{\mu} \cdot \mathbf{B}$. Por consiguiente, el operador momento dipolar magnético predicho por Dirac es:
\begin{equation}
\boldsymbol{\mu} = \frac{e}{2M} \boldsymbol{\sigma}.
\end{equation}
Dado que el operador de momento angular de espín mecánico para una partícula de espín $1/2$ es $\mathbf{S} = \frac{1}{2}\boldsymbol{\sigma} \iff \boldsymbol{\sigma} = 2\mathbf{S}$, reescribimos el momento magnético en la forma estándar de Landé:
\begin{equation}
\boldsymbol{\mu} = \frac{e}{2M} (2\mathbf{S}) = 2 \left( \frac{e}{2M} \right) \mathbf{S} \equiv g \left( \frac{e}{2M} \right) \mathbf{S}.
\end{equation}
Comparando directamente ambos miembros, la teoría relativista de Dirac fija de forma unívoca y no negociable el factor giromagnético de cualquier fermión puntual de espín $1/2$:
\begin{equation}
\dest{g_{\text{Dirac}} \equiv 2 \implies \mu_{\text{Dirac}} = \frac{e}{2M}}
\label{eq:g_dirac_exact}
\end{equation}

Conviene subrayar la naturaleza lógica de lo que acabamos de obtener, porque de ella depende todo el resto del artículo. No se ha ajustado ningún parámetro ni se ha elegido ninguna constante: el valor $g = 2$ es una consecuencia forzosa de exigir invariancia relativista a una partícula de espín $1/2$, y en este sentido es una predicción falsable y no una convención. La implicación es inmediata y muy fuerte: cualquier medida de $g$ distinta de $2$ en una partícula de espín $1/2$ constituye una afirmación sobre su estructura interna, y no una imprecisión del experimento ni una libertad del modelo. Con la predicción así fijada, el experimento puede ahora pronunciarse.

\subsection{El Quiebre Experimental de Stern (1933)}
Para el protón ($M_p \approx 938.272\text{ MeV}$), se define el magnetón nuclear:
\begin{equation}
\mu_N \equiv \frac{e}{2M_p}.
\end{equation}
La predicción puntual de Dirac exigía estrictamente:
\begin{equation}
\mu_p^{\text{Dirac}} = 1.00\,\mu_N, \qquad \mu_n^{\text{Dirac}} = 0.00\,\mu_N.
\end{equation}
Sin embargo, en 1933, Otto Stern, I. Estermann y R. Frisch midieron la deflexión de haces moleculares de $H_2$ en campos magnéticos de gradiente extremo \cite{stern1933,griffiths2008}, obteniendo:
\begin{equation}
\mu_p = +2.7928473446(8)\,\mu_N \implies g_p \approx 5.586 \neq 2.
\end{equation}
Posteriormente, para el neutrón (partícula neutra, $Q=0$):
\begin{equation}
\mu_n = -1.9130427(5)\,\mu_N \implies g_n \approx -3.826 \neq 0.
\end{equation}
Definiendo el momento magnético anómalo $\kappa$ mediante:
\begin{equation}
\mu \equiv \left( \frac{Q}{e} + \kappa \right) \mu_N,
\end{equation}
se obtienen anomalías nucleónicas enormes:
\begin{equation}
\kappa_p \approx +1.7928, \qquad \kappa_n \approx -1.9130.
\end{equation}

El resultado del neutrón merece un comentario aparte, porque es el más difícil de conciliar con una imagen ingenua y el más revelador sobre la naturaleza del problema. Una partícula estrictamente neutra y sin estructura no puede tener momento magnético alguno; sin embargo, el neutrón lo tiene, y de magnitud comparable a la del protón. La explicación es que su carga eléctrica \emph{total} es nula---lo que impone $\int \rho_n(\mathbf{r})\,d^3\mathbf{r} = 0$ y, como veremos, $F_1^n(0) = G_E^n(0) = 0$---pero su carga eléctrica \emph{local} no lo es: el neutrón está compuesto por tres quarks de valencia ($udd$) de cargas fraccionarias $+2/3$, $-1/3$ y $-1/3$, distribuidos de manera no homogénea en el espacio. Al ponderar por la distancia las corrientes internas asociadas a esa distribución, las contribuciones positivas y negativas ya no se cancelan, y el momento magnético resultante es no nulo y negativo. Es decir: la anomalía magnética del neutrón es, en sí misma, una medida indirecta de que existe una distribución espacial interna. Esta observación anticipa la distinción entre carga integrada y distribución de carga que estructurará la Sección~\ref{sec:radio_fourier}.

A diferencia de la anomalía leptónica del electrón ($\kappa_e = \frac{\alpha}{2\pi} \approx 0.00116$), explicada cuantitativamente por la fluctuación del vacío a un bucle de QED (Schwinger 1948), las anomalías del protón y neutrón son de orden de magnitud unidad ($\mathcal{O}(1)$). Esta comparación es el argumento decisivo: una anomalía de una parte en mil admite una explicación perturbativa dentro de la propia QED sin renunciar a la puntualidad; una anomalía de orden unidad, en cambio, no puede provenir de correcciones radiativas y obliga a admitir grados de libertad internos. Esta es la prueba empírica incontestable de que el protón y el neutrón poseen una compleja estructura espacial interna constituida por quarks y gluones \cite{henley2007,das2003,pdg2022}.


\subsection{Hoja de Ruta Analítica: De QED a la Transformada de Fourier y el Radio del Protón}
La demostración experimental de que el protón es un objeto extendido impone un doble desafío teórico y fenomenológico que articula el desarrollo de este artículo. Conviene enunciarlo como dos preguntas, porque cada una de ellas organiza la mitad del documento: primero, ¿cuántas funciones desconocidas hacen falta para describir un nucleón con estructura, y de dónde sale exactamente la fórmula de Rosenbluth?; segundo, una vez medidas esas funciones, ¿cómo se traducen en una distribución espacial y en un radio?
\begin{enumerate}
    \item \textbf{Primer Objetivo (Demostración microscópica en QED):} Partiendo del lagrangiano fundamental de QED, deducir analíticamente y paso a paso la sección eficaz de Rosenbluth, Ec.~\eqref{eq:rosenbluth_master_intro}. Esto requiere construir el operador de vértice covariante más general mediante simetrías de Lorentz y gauge (Sección~\ref{sec:vertice_gordon}), calcular las trazas de Dirac de los tensores leptónico y hadrónico en el laboratorio (Sección~\ref{sec:cinematica_trazas}), y resolver el problema de las interferencias cruzadas transformando al sistema de Breit (Sección~\ref{sec:marco_breit}), para finalmente integrar el espacio de fases relativista de dos cuerpos y formular la técnica de separación lineal de Rosenbluth (Sección~\ref{sec:seccion_rosenbluth}).
    \item \textbf{Segundo Objetivo (Transformada de Fourier y Radio del Protón):} Una vez aisladas cuantitativamente las funciones $G_E(Q^2)$ y $G_M(Q^2)$, tender el puente matemático hacia el espacio de coordenadas reales. En la Sección~\ref{sec:radio_fourier}, demostramos que en el sistema de Breit $G_E(Q^2)$ es la transformada de Fourier tridimensional de la densidad espacial de carga $\rho(r)$, probamos mediante expansión de Taylor que su pendiente en el origen determina unívocamente el radio cuadrático medio de carga $\langle r_E^2 \rangle = -6 \left.\frac{dG_E}{dQ^2}\right|_0$, e invertimos analíticamente el modelo dipolar mediante cálculo de residuos de Cauchy en el plano complejo para obtener la densidad exponencial del protón y su radio numérico, confrontándolo con las fronteras modernas de la física subatómica (transferencia de polarización en Jefferson Lab, intercambio de dos fotones TPE y el puzzle del radio del protón).
\end{enumerate}
El Apéndice~\ref{sec:apendice_dirac}, por su parte, reúne y demuestra desde cero todo el aparato algebraico de matrices gamma y de trazas que el cuerpo del texto utilizará repetidamente; el lector que no haya trabajado antes con el álgebra de Clifford puede consultarlo en cualquier momento sin perder el hilo del argumento principal.

Para emprender esta demostración desde la teoría cuántica de campos, el primer paso fundamental consiste en determinar cómo se acopla microscópicamente el fotón virtual con un barión compuesto de espín $1/2$. No podemos asumir a priori la existencia de dos factores de forma; es indispensable construir el operador de vértice más general permitido por la covarianza de Lorentz y las simetrías discretas, y aplicar rigurosamente la conservación de la corriente electromagnética de Noether para restringir sus grados de libertad.

\section{Parametrización de la Corriente Hadrónica y la Identidad de Gordon}
\label{sec:vertice_gordon}

La sección anterior dejó establecido un hecho y una carencia. El hecho: el protón no es puntual, y la anomalía de su momento magnético lo demuestra sin apelación posible. La carencia: no sabemos todavía con cuántas funciones desconocidas hay que describirlo. Es tentador resolver esa carencia escribiendo directamente el vértice de Dirac-Pauli y llamándolo definición, que es lo que hacen la mayoría de los tratamientos introductorios. Aquí procederemos en sentido contrario. Partiremos de que \emph{no sabemos absolutamente nada} sobre el interior del nucleón y preguntaremos cuánto puede afirmarse aun así. La respuesta es sorprendentemente fuerte: dos principios de simetría---la invariancia de Lorentz y la conservación de la corriente electromagnética---bastan para demostrar que el número de funciones escalares desconocidas es exactamente dos, ni una más. En ningún punto de la demostración intervendrán los quarks, los gluones ni modelo alguno de estructura interna. Ese es precisamente el valor del argumento: cualquier teoría futura de la estructura del nucleón, por exótica que sea, tendrá que seguir describiéndose con dos factores de forma.

\subsection{Construcción Axiomática del Vértice General}
Para restringir el objeto desconocido hay que empezar por situarlo dentro de la amplitud completa del proceso. Esta subsección hace exactamente eso: identifica en qué lugar del diagrama de Feynman de la dispersión elástica aparece la estructura del protón (Figura~\ref{fig:feynman_tree}), y a continuación enumera sistemáticamente qué estructuras matriciales admite la covarianza de Lorentz en ese lugar.

\begin{figure}[t]
\centering
\begin{tikzpicture}[>=stealth, line width=0.7pt, font=\small]

%----------------------------------------------------------------------
% LÍNEA LEPTÓNICA (arriba): electrón puntual
%----------------------------------------------------------------------
\draw[->] (-2.70, 1.95) -- (-1.45, 1.78);
\draw     (-1.45, 1.78) -- ( 0.00, 1.60);
\draw[->] ( 0.00, 1.60) -- ( 1.45, 1.78);
\draw     ( 1.45, 1.78) -- ( 2.70, 1.95);
\node[anchor=east, inner sep=3pt] at (-2.76, 1.95) {$e^-(k)$};
\node[anchor=west, inner sep=3pt] at ( 2.76, 1.95) {$e^-(k')$};

% vértice leptónico: puntual y exactamente conocido
\fill (0,1.60) circle (2.2pt);
\node[anchor=south east, font=\footnotesize, inner sep=3pt] at (-0.10, 1.66) {$-ie\,\gamma^\mu$};

%----------------------------------------------------------------------
% FOTÓN VIRTUAL
%----------------------------------------------------------------------
\draw[decorate, decoration={coil, aspect=0, segment length=5.2pt, amplitude=3.2pt}]
      (0,1.54) -- (0,-1.36);
\node[anchor=east, font=\footnotesize, inner sep=4pt] at (-0.16, 0.42) {$\dfrac{-i\,g_{\mu\nu}}{q^2}$};
\node[anchor=west, font=\footnotesize, inner sep=4pt, align=left] at (0.18, 0.20)
      {$\gamma^*$\\[0.4ex] $q^\mu = k^\mu - k'^\mu$\\[0.2ex] $\phantom{q^\mu} = p'^\mu - p^\mu$};

%----------------------------------------------------------------------
% VÉRTICE HADRÓNICO: estructura desconocida
%----------------------------------------------------------------------
\shade[inner color=orange!60, outer color=orange!10] (0,-1.60) circle (0.46);
\draw[orange!65!black, line width=0.5pt] (0,-1.60) circle (0.46);
\node[font=\footnotesize] at (0,-1.60) {$ie\,\Gamma^\nu$};

%----------------------------------------------------------------------
% LÍNEA HADRÓNICA (abajo): protón con estructura
%----------------------------------------------------------------------
\draw[->, line width=1.1pt] (-2.70,-1.95) -- (-1.55,-1.83);
\draw[line width=1.1pt]     (-1.55,-1.83) -- (-0.44,-1.70);
\draw[line width=1.1pt]     ( 0.44,-1.70) -- ( 1.55,-1.83);
\draw[->, line width=1.1pt] ( 1.55,-1.83) -- ( 2.70,-1.95);
\node[anchor=east, inner sep=3pt] at (-2.76,-1.95) {$p(p)$};
\node[anchor=west, inner sep=3pt] at ( 2.76,-1.95) {$p(p')$};

%----------------------------------------------------------------------
% ANOTACIONES LATERALES
%----------------------------------------------------------------------
\node[anchor=north west, font=\footnotesize, align=left, text width=3.9cm] at (4.30, 2.15)
     {\textbf{Conocido exactamente.} El electrón es un fermión elemental: su acoplamiento al fotón es $\gamma^\mu$, sin ninguna función libre.};
\node[anchor=north west, font=\footnotesize, align=left, text width=3.9cm] at (4.30, 0.15)
     {\textbf{Desconocido.} Toda la estructura interna del nucleón se concentra aquí. La Sección~\ref{sec:vertice_gordon} demuestra que cabe en dos funciones escalares.};

% eje temporal
\draw[->, gray!60, line width=0.5pt] (-4.15,-2.45) -- (-4.15, 2.45);
\node[gray!70, font=\scriptsize, rotate=90, anchor=south] at (-4.32, 0) {tiempo};

\end{tikzpicture}
\caption{Diagrama de Feynman a nivel árbol del proceso $e^-(k) + p(p) \to e^-(k') + p(p')$ en la aproximación de Born de intercambio de un solo fotón virtual. La lectura del diagrama es la siguiente: el electrón incidente de cuadrimomento $k$ emite un fotón virtual $\gamma^*$ que transporta el cuadrimomento $q^\mu = k^\mu - k'^\mu$, y ese mismo fotón es absorbido por el protón, que retrocede con $p'^\mu = p^\mu + q^\mu$. La asimetría entre los dos extremos de la línea de fotón es el contenido físico de toda esta sección. En el vértice superior sabemos exactamente qué escribir---el electrón es un fermión elemental de Dirac y su acoplamiento es la matriz $\gamma^\mu$ multiplicada por la carga, sin ninguna función libre---y lo mismo ocurre con el propagador del fotón, $-ig_{\mu\nu}/q^2$, que es un objeto de QED pura. En el vértice inferior, en cambio, no sabemos qué escribir: el protón es un estado ligado de quarks y gluones, y toda esa complejidad desconocida queda encapsulada en el operador matricial $\Gamma^\nu(p,p')$, representado aquí como una región difusa. El problema que resolveremos a continuación no es calcular $\Gamma^\nu$ a partir de la dinámica interna---eso exigiría cromodinámica cuántica no perturbativa---sino algo más modesto y a la vez más robusto: determinar, usando solo simetrías, cuántos grados de libertad independientes puede contener.}
\label{fig:feynman_tree}
\end{figure}

Considérese el proceso de dispersión elástica $e^-(k) + p(p) \to e^-(k') + p(p')$. En la teoría de perturbaciones de QED a primer orden (aproximación de Born de un solo fotón virtual $\gamma^*$)~\cite{halzen1984,thomson2013}, el elemento de la matriz $S$ se deriva a partir de la fórmula de reducción LSZ o la serie de Dyson:
\begin{equation}
S_{fi} = -i \int d^4x \, \mathcal{H}_{\text{int}}(x) = -i \int d^4x \, j^\mu_{\text{lept}}(x) A_\mu(x),
\end{equation}
donde el fotón virtual conecta la corriente leptónica del electrón con la corriente electromagnética hadrónica del protón a través del propagador de Feynman en el gauge de Feynman:
\begin{equation}
D_{\mu\nu}(x - y) = \int \frac{d^4q}{(2\pi)^4} \frac{-i g_{\mu\nu}}{q^2 + i\epsilon} e^{-i q \cdot (x - y)}.
\end{equation}
Efectuando las integrales espacio-temporales en $d^4x$ y $d^4y$, la conservación del cuadrimomento total $(2\pi)^4 \delta^{(4)}(k + p - k' - p')$ factoriza la amplitud invariante de Feynman:
\begin{equation}
\mathcal{M} = \left[ \bar{u}(k') (-ie\gamma_\mu) u(k) \right] \left( \frac{-i g^{\mu\nu}}{q^2} \right) \langle p(p') | j_\nu(0) | p(p) \rangle,
\label{eq:feynman_amplitude_full}
\end{equation}
donde el cuadrimomento transferido por el fotón virtual es:
\begin{equation}
q^\mu \equiv k^\mu - k'^\mu = p'^\mu - p^\mu.
\end{equation}
Los dos primeros corchetes de la Ec.~\eqref{eq:feynman_amplitude_full} son enteramente calculables: corresponden al vértice superior y a la línea de fotón de la Figura~\ref{fig:feynman_tree}. Todo lo que ignoramos está en el tercer factor, el elemento de matriz de la corriente hadrónica, y es sobre él que debemos operar.

El elemento de matriz hadrónico entre estados asintóticos fermiónicos libres on-shell ($p^2 = p'^2 = M^2$) se parametriza mediante un operador matricial de vértice $\Gamma^\mu(p, p')$ acoplado a la carga elemental $e$:
\begin{equation}
\langle p(p') | j^\mu(0) | p(p) \rangle = e\,\bar{u}(p')\,\Gamma^\mu(p, p')\,u(p),
\label{eq:hadron_current_matrix_element}
\end{equation}
donde los espinores de Dirac libres satisfacen las ecuaciones de movimiento:
\begin{equation}
(\slashed{p} - M) u(p) = 0 \iff \slashed{p} u(p) = M u(p), \qquad \bar{u}(p') (\slashed{p}' - M) = 0 \iff \bar{u}(p') \slashed{p}' = M \bar{u}(p').
\label{eq:dirac_free_eqs}
\end{equation}
Conviene insistir en que la Ec.~\eqref{eq:hadron_current_matrix_element} no es un resultado sino una \emph{definición}: hemos bautizado $\Gamma^\mu$ a lo que no sabemos, sin ganar todavía ninguna información. La ganancia empieza ahora, cuando preguntemos qué restricciones pesan sobre ese objeto por el solo hecho de vivir en un espacio-tiempo de Minkowski y de describir una interacción electromagnética.

Definiendo la combinación simétrica de momentos $P^\mu \equiv p'^\mu + p^\mu$, evaluamos su producto escalar con la transferencia $q^\mu$:
\begin{equation}
P \cdot q = (p' + p)_\mu (p' - p)^\mu = p'^2 - p' \cdot p + p \cdot p' - p^2 = p'^2 - p^2 = M^2 - M^2 = 0.
\label{eq:P_dot_q_zero}
\end{equation}
Por consiguiente, $P$ y $q$ son ortogonales en el espacio de Minkowski, y la única variable cinemática escalar Lorentz-invariante independiente es la masa virtual del fotón:
\begin{equation}
q^2 \equiv -Q^2 \le 0.
\end{equation}
Esta ortogonalidad tiene una consecuencia práctica que conviene retener desde ahora: de las dos combinaciones de momentos disponibles solo una produce un escalar independiente, de modo que las funciones desconocidas que aparezcan en el vértice no podrán depender de dos variables sino de una sola. Toda la estructura del protón, cualquiera que sea, será función únicamente de $q^2$. Esto es lo que hace que la medida experimental sea un problema abordable: basta con barrer un parámetro.

Para determinar la estructura matricial de $\Gamma^\mu(p, p')$, imponemos los principios de simetría del Modelo Estándar~\cite{henley2007,griffiths2008}:
\begin{enumerate}
    \item \textbf{Invarianza de Lorentz:} $\Gamma^\mu$ debe transformar como un cuadrivector bajo transformaciones ortócronas de Lorentz. Los únicos cuadrivectores independientes disponibles en el sistema son $\gamma^\mu$, $P^\mu = (p+p')^\mu$ y $q^\mu = (p'-p)^\mu$.
    \item \textbf{Conservación de Paridad ($\mathcal{P}$):} La interacción electromagnética conserva estrictamente la paridad espacial. Bajo paridad, un campo espinorial transforma según $\psi(\mathbf{x}, t) \to \psi'(\mathbf{x}', t) = \gamma^0 \psi(-\mathbf{x}, t)$. La corriente dipolar fermiónica transforma como:
    \begin{equation}
    \bar{\psi} \gamma^\mu \psi \xrightarrow{\mathcal{P}} \bar{\psi} \gamma^0 \gamma^\mu \gamma^0 \psi = \begin{cases} +\bar{\psi}\gamma^0\psi & (\mu=0) \\ -\bar{\psi}\gamma^j\psi & (\mu=j) \end{cases} = (\gamma^\mu)_\nu (\bar{\psi}\gamma^\nu\psi),
    \end{equation}
    es decir, como un cuadrivector polar genuino. En contraste, cualquier término bilineal que contenga la matriz pseudoescalar $\gamma^5$ transforma con signo opuesto:
    \begin{equation}
    \bar{\psi} \gamma^\mu \gamma^5 \psi \xrightarrow{\mathcal{P}} -\begin{cases} +\bar{\psi}\gamma^0\gamma^5\psi & (\mu=0) \\ -\bar{\psi}\gamma^j\gamma^5\psi & (\mu=j) \end{cases}
    \end{equation}
    (comportamiento pseudovectorial o axial). Por tanto, la conservación de paridad excluye de forma rigurosa combinaciones como $\gamma^\mu\gamma^5$, $P^\mu\gamma^5$, $q^\mu\gamma^5$, y contracciones con el pseudotensor de Levi-Civita $\epsilon^{\mu\nu\alpha\beta}P_\nu q_\alpha \gamma_\beta$.
\end{enumerate}
Por ende, la combinación lineal covariante y polar más general posible para el vértice hadrónico depende únicamente de tres factores escalares de forma $A(q^2)$, $B(q^2)$ y $C(q^2)$:
\begin{equation}
\Gamma^\mu(p, p') = A(q^2) \gamma^\mu + B(q^2) (p + p')^\mu + C(q^2) q^\mu.
\label{eq:vertice_general_abc}
\end{equation}

Merece la pena detenerse en lo que acaba de obtenerse, porque es un resultado más fuerte de lo que su brevedad sugiere. No hemos supuesto nada sobre el protón salvo que tiene espín $1/2$ y que interactúa electromagnéticamente conservando paridad. Con eso solo, toda la complejidad concebible de su interior---tres quarks de valencia, un mar de pares quark-antiquark, un campo gluónico no abeliano, cualquier cosa que la cromodinámica cuántica tenga que decir---ha quedado comprimida en tres funciones de una sola variable. Tres funciones son, sin embargo, una más de las que finalmente necesitaremos, y las dos subsecciones siguientes se dedican a eliminar la sobrante. El camino tiene dos etapas de naturaleza muy distinta: la primera es puramente algebraica y consiste en reconocer que el término de convección $(p+p')^\mu$ no es realmente independiente de $\gamma^\mu$ cuando se lo intercala entre espinores en capa de masa; la segunda es física y apela a la conservación de la carga eléctrica.

\subsection{Deducción Paso a Paso de la Identidad de Gordon}
La Identidad de Gordon es un teorema fundamental del álgebra de Dirac~\cite{gordon1928,halzen1984} que permite sustituir el término de convección $(p+p')^\mu$ en favor del término tensorial de espín $\sigma^{\mu\nu} q_\nu$. Formulado así suena a tecnicismo prescindible, pero su papel en el argumento es sustantivo, y conviene entenderlo antes de emprender el cálculo. Lo que la identidad proporciona es un \emph{cambio de base} en el espacio de estructuras matriciales admisibles: en lugar de describir el vértice con el par $\{\gamma^\mu,\,(p+p')^\mu\}$, lo describiremos con el par $\{\gamma^\mu,\,\sigma^{\mu\nu}q_\nu\}$. La razón para preferir la segunda base es física y no estética. El término de convección $(p+p')^\mu$ mezcla dos efectos que quisiéramos separar---el arrastre de la carga por el movimiento del nucleón y su momento magnético intrínseco---mientras que la pareja $\{\gamma^\mu,\,\sigma^{\mu\nu}q_\nu\}$ los aísla limpiamente: $\gamma^\mu$ es exactamente la corriente de una carga sin estructura, y $\sigma^{\mu\nu}q_\nu$ tiene exactamente la forma de un acoplamiento dipolar magnético. En esta segunda base, y solo en ella, los coeficientes que acompañen a cada estructura tendrán interpretación física directa: uno medirá carga y el otro magnetización anómala.

Antes de empezar, una advertencia que el lector debe retener porque volverá a ser relevante. La identidad que vamos a demostrar depende de manera esencial de las ecuaciones de Dirac libres, Ec.~\eqref{eq:dirac_free_eqs}, y por tanto es válida únicamente cuando ambos nucleones están en capa de masa ($p^2 = p'^2 = M^2$), es decir, cuando son partículas reales entrantes y salientes. Si el nucleón fuese una línea interna virtual---propagándose dentro de un bucle, o en un estado fuera de la capa de masa---aparecerían términos adicionales proporcionales a $(\slashed{p} - M)$ que rompen la equivalencia algebraica entre $(p+p')^\mu$ y $\sigma^{\mu\nu}q_\nu$, y el vértice requeriría factores de forma \emph{off-shell} adicionales. En la dispersión elástica a nivel árbol que nos ocupa, la condición de capa de masa se satisface exactamente, y por eso podemos emplear la identidad sin reservas.

Partimos de la definición del conmutador tensorial de Dirac:
\begin{equation}
\sigma^{\mu\nu} \equiv \frac{i}{2} [\gamma^\mu, \gamma^\nu] = \frac{i}{2} (\gamma^\mu \gamma^\nu - \gamma^\nu \gamma^\mu).
\end{equation}
Multiplicando por $+i$ en ambos lados y recordando que $i^2 = -1$:
\begin{equation}
i \sigma^{\mu\nu} = -\frac{1}{2} (\gamma^\mu \gamma^\nu - \gamma^\nu \gamma^\mu).
\end{equation}
Contrayendo con el cuadrimomento transferido $q_\nu = (p' - p)_\nu$:
\begin{align}
i \sigma^{\mu\nu} q_\nu &= -\frac{1}{2} (\gamma^\mu \gamma^\nu - \gamma^\nu \gamma^\mu)(p'_\nu - p_\nu) \nonumber \\
&= -\frac{1}{2} \left[ \gamma^\mu \gamma^\nu p'_\nu - \gamma^\mu \gamma^\nu p_\nu - \gamma^\nu \gamma^\mu p'_\nu + \gamma^\nu \gamma^\mu p_\nu \right] \nonumber \\
&= -\frac{1}{2} \left[ \gamma^\mu \slashed{p}' - \gamma^\mu \slashed{p} - \slashed{p}' \gamma^\mu + \slashed{p} \gamma^\mu \right].
\label{eq:i_sigma_expanded}
\end{align}
La estrategia del cálculo es ahora transparente. Cada uno de los cuatro sumandos de la Ec.~\eqref{eq:i_sigma_expanded} contiene una matriz $\slashed{p}$ o $\slashed{p}'$ que, al intercalarse entre espinores en capa de masa, puede ser \emph{absorbida}: si está adyacente al espinor correspondiente, las ecuaciones de Dirac la convierten directamente en el número $M$; si no lo está, basta anticonmutarla a través de $\gamma^\mu$ usando el álgebra de Clifford para dejarla adyacente. Intercalemos entonces el operador entre los espinores libres $\bar{u}(p')$ y $u(p)$:
\begin{equation}
\bar{u}(p') i\sigma^{\mu\nu} q_\nu u(p) = -\frac{1}{2} \left[ \mathcal{T}_1 + \mathcal{T}_2 + \mathcal{T}_3 + \mathcal{T}_4 \right],
\end{equation}
donde definimos los cuatro términos matriciales:
\begin{align}
\mathcal{T}_1 &\equiv \bar{u}(p') \gamma^\mu \slashed{p}' u(p), \qquad \mathcal{T}_2 \equiv -\bar{u}(p') \gamma^\mu \slashed{p} u(p), \\
\mathcal{T}_3 &\equiv -\bar{u}(p') \slashed{p}' \gamma^\mu u(p), \qquad \mathcal{T}_4 \equiv \bar{u}(p') \slashed{p} \gamma^\mu u(p).
\end{align}

Evaluemos algebraicamente cada término paso a paso:
\begin{itemize}
    \item \textbf{Cálculo de $\mathcal{T}_1$:}
    Utilizamos la relación de anticonmutación de Clifford $\{\gamma^\mu, \gamma^\nu\} = 2g^{\mu\nu}\mathbb{I}$:
    \begin{equation}
    \gamma^\mu \slashed{p}' = \gamma^\mu \gamma^\nu p'_\nu = (2 g^{\mu\nu}\mathbb{I} - \gamma^\nu \gamma^\mu) p'_\nu = 2 p'^\mu \mathbb{I} - \slashed{p}' \gamma^\mu.
    \end{equation}
    Sustituyendo entre los espinores:
    \begin{equation}
    \mathcal{T}_1 = \bar{u}(p') [2 p'^\mu - \slashed{p}' \gamma^\mu] u(p) = 2 p'^\mu \bar{u}(p') u(p) - [\bar{u}(p') \slashed{p}'] \gamma^\mu u(p).
    \end{equation}
    Aplicando la ecuación adjunta de Dirac en la capa de masa $\bar{u}(p') \slashed{p}' = M \bar{u}(p')$:
    \begin{equation}
    \mathcal{T}_1 = 2 p'^\mu \bar{u}(p') u(p) - M \bar{u}(p') \gamma^\mu u(p).
    \label{eq:T1_eval}
    \end{equation}

    \item \textbf{Cálculo de $\mathcal{T}_2$:}
    Aplicamos directamente la ecuación de Dirac hacia la derecha sobre $u(p)$, $\slashed{p} u(p) = M u(p)$:
    \begin{equation}
    \mathcal{T}_2 = -\bar{u}(p') \gamma^\mu [\slashed{p} u(p)] = - M \bar{u}(p') \gamma^\mu u(p).
    \label{eq:T2_eval}
    \end{equation}

    \item \textbf{Cálculo de $\mathcal{T}_3$:}
    Aplicamos directamente la ecuación de Dirac hacia la izquierda sobre $\bar{u}(p')$, $[\bar{u}(p') \slashed{p}'] = M \bar{u}(p')$:
    \begin{equation}
    \mathcal{T}_3 = - [\bar{u}(p') \slashed{p}'] \gamma^\mu u(p) = - M \bar{u}(p') \gamma^\mu u(p).
    \label{eq:T3_eval}
    \end{equation}

    \item \textbf{Cálculo de $\mathcal{T}_4$:}
    Anticonmutamos el operador $\slashed{p}$ a través de $\gamma^\mu$:
    \begin{equation}
    \slashed{p} \gamma^\mu = p_\nu \gamma^\nu \gamma^\mu = p_\nu (2 g^{\nu\mu}\mathbb{I} - \gamma^\mu \gamma^\nu) = 2 p^\mu \mathbb{I} - \gamma^\mu \slashed{p}.
    \end{equation}
    Sustituyendo y aplicando la ecuación de Dirac $\slashed{p} u(p) = M u(p)$:
    \begin{align}
    \mathcal{T}_4 &= \bar{u}(p') [2 p^\mu - \gamma^\mu \slashed{p}] u(p) = 2 p^\mu \bar{u}(p') u(p) - \bar{u}(p') \gamma^\mu [\slashed{p} u(p)] \nonumber \\
    &= 2 p^\mu \bar{u}(p') u(p) - M \bar{u}(p') \gamma^\mu u(p).
    \label{eq:T4_eval}
    \end{align}
\end{itemize}
Obsérvese el patrón: los cuatro términos han producido entre todos exactamente dos tipos de objeto, la densidad escalar $\bar{u}(p')u(p)$ multiplicada por momentos, y la corriente vectorial $\bar{u}(p')\gamma^\mu u(p)$ multiplicada por $M$. No podía ser de otro modo, y eso es lo que hace que la identidad exista.

Sumando algebraicamente las cuatro contribuciones Ecs.~\eqref{eq:T1_eval}, \eqref{eq:T2_eval}, \eqref{eq:T3_eval} y \eqref{eq:T4_eval}:
\begin{align}
\sum_{k=1}^4 \mathcal{T}_k &= \left[ 2 p'^\mu \bar{u}(p') u(p) - M \bar{u}(p') \gamma^\mu u(p) \right] - M \bar{u}(p') \gamma^\mu u(p) \nonumber \\
&\quad - M \bar{u}(p') \gamma^\mu u(p) + \left[ 2 p^\mu \bar{u}(p') u(p) - M \bar{u}(p') \gamma^\mu u(p) \right] \nonumber \\
&= 2 (p'^\mu + p^\mu) \bar{u}(p') u(p) - (M + M + M + M) \bar{u}(p') \gamma^\mu u(p) \nonumber \\
&= 2 (p + p')^\mu \bar{u}(p') u(p) - 4 M \bar{u}(p') \gamma^\mu u(p).
\end{align}
Multiplicando por el prefactor $-1/2$:
\begin{align}
\bar{u}(p') i\sigma^{\mu\nu} q_\nu u(p) &= -\frac{1}{2} \left[ 2 (p + p')^\mu \bar{u}(p') u(p) - 4 M \bar{u}(p') \gamma^\mu u(p) \right] \nonumber \\
&= -(p + p')^\mu \bar{u}(p') u(p) + 2M \bar{u}(p') \gamma^\mu u(p).
\end{align}
Despejando el término de convección $(p+p')^\mu$, arribamos a la formulación canónica de la \textbf{Identidad de Gordon}:
\begin{equation}
\boxed{\bar{u}(p') (p + p')^\mu u(p) = \bar{u}(p') \left[ 2M \gamma^\mu - i \sigma^{\mu\nu} q_\nu \right] u(p)}
\label{eq:gordon_identity_full}
\end{equation}
Equivalentemente, sumando $i\sigma^{\mu\nu}q_\nu$ en ambos miembros y dividiendo entre $2M$, despejamos la corriente de Dirac:
\begin{equation}
\dest{\bar{u}(p') \gamma^\mu u(p) = \frac{1}{2M} \bar{u}(p') \left[ (p + p')^\mu + i \sigma^{\mu\nu} q_\nu \right] u(p)}
\label{eq:gordon_gamma_form}
\end{equation}

La segunda forma, Ec.~\eqref{eq:gordon_gamma_form}, admite una lectura física que conviene señalar porque ilumina todo lo que sigue: afirma que la corriente de una partícula de Dirac puntual, aparentemente elemental, se descompone espontáneamente en una parte de convección proporcional a $(p+p')^\mu$---el arrastre de la carga por el movimiento, que es lo único que tendría una partícula de espín cero---y una parte de espín proporcional a $\sigma^{\mu\nu}q_\nu$, que es el magnetismo. El momento magnético de Dirac no es, pues, un añadido: es una de las dos mitades en que la propia corriente $\gamma^\mu$ se separa. La forma que usaremos de inmediato, en cambio, es la primera, Ec.~\eqref{eq:gordon_identity_full}, que permite eliminar $(p+p')^\mu$ del vértice general en cuanto sepamos que el tercer término $C(q^2)q^\mu$ desaparece. Demostrar esa desaparición es lo que hacemos a continuación, y requiere abandonar el álgebra pura e invocar un principio físico.

\subsection{Identidad de Ward y Anulación de \texorpdfstring{$C(q^2)$}{C(q\^2)}}
El término longitudinal $C(q^2)q^\mu$ no puede eliminarse por manipulación algebraica: es una estructura covariante perfectamente legítima y ninguna identidad de matrices gamma la hace desaparecer. Su eliminación exige un ingrediente nuevo, y ese ingrediente es físico: la conservación de la carga eléctrica. La cadena lógica que recorreremos tiene cuatro eslabones, y vale la pena tenerla a la vista antes de escribir una sola ecuación. La invariancia de norma local $U(1)$ de la lagrangiana de QED implica, por el teorema de Noether, la conservación local de la corriente electromagnética; tomar el elemento de matriz de esa ley de conservación entre estados de protón la traduce al espacio de momentos como la \emph{condición de Ward}; y contraer el vértice general con $q_\mu$ bajo esa condición anula unívocamente el coeficiente $C$. El resultado será que la estructura longitudinal está prohibida no por accidente sino por la misma razón por la que la carga del protón no cambia durante la colisión.

En la teoría cuántica de campos, la invarianza de norma $U(1)$ local de la lagrangiana de QED garantiza, vía el Teorema de Noether, que la corriente electromagnética se conserva de forma local y exacta~\cite{ward1950,thomson2013}:
\begin{equation}
\partial_\mu j^\mu(x) = 0.
\end{equation}
Tomando el elemento de matriz de transición hadrónico entre estados asintóticos:
\begin{equation}
\langle p(p') | \partial_\mu j^\mu(x) | p(p) \rangle = 0.
\end{equation}
Mediante el operador de traslación espacio-temporal $j^\mu(x) = e^{i \hat{P} \cdot x} j^\mu(0) e^{-i \hat{P} \cdot x}$, donde $\hat{P}^\mu |p\rangle = p^\mu |p\rangle$:
\begin{equation}
\langle p(p') | j^\mu(x) | p(p) \rangle = e^{i(p' - p) \cdot x} \langle p(p') | j^\mu(0) | p(p) \rangle = e^{i q \cdot x} \langle p(p') | j^\mu(0) | p(p) \rangle.
\end{equation}
Aplicando la derivada parcial espacio-temporal $\partial_\mu$:
\begin{equation}
\partial_\mu \langle p(p') | j^\mu(x) | p(p) \rangle = i q_\mu e^{i q \cdot x} \langle p(p') | j^\mu(0) | p(p) \rangle = 0.
\end{equation}
Evaluando en $x=0$ y dividiendo entre la unidad imaginaria $i$, obtenemos la \textbf{Condición de Ward} en el espacio de momentos:
\begin{equation}
q_\mu \langle p(p') | j^\mu(0) | p(p) \rangle = 0 \iff q_\mu \, \bar{u}(p') \Gamma^\mu(p, p') u(p) = 0.
\label{eq:ward_condition}
\end{equation}
La traducción ha convertido una ecuación diferencial en el espacio-tiempo en una condición algebraica en el espacio de momentos: la corriente hadrónica debe ser \emph{transversal} a la transferencia de momento. Aplicarla es ahora una cuestión de contraer y contar.

Contrayamos ahora término a término el vértice general de la Ec.~\eqref{eq:vertice_general_abc} con $q_\mu$:
\begin{equation}
q_\mu \bar{u}(p') \Gamma^\mu u(p) = A(q^2) \left[ q_\mu \bar{u}(p') \gamma^\mu u(p) \right] + B(q^2) \left[ q_\mu (p + p')^\mu \bar{u}(p') u(p) \right] + C(q^2) \left[ q_\mu q^\mu \bar{u}(p') u(p) \right].
\end{equation}
Evaluemos minuciosamente cada uno de los tres términos escalares:
\begin{enumerate}
    \item \textbf{Término de Dirac ($q_\mu \gamma^\mu$):}
    \begin{align}
    q_\mu \bar{u}(p') \gamma^\mu u(p) &= \bar{u}(p') \slashed{q} u(p) = \bar{u}(p') (\slashed{p}' - \slashed{p}) u(p) \nonumber \\
    &= [\bar{u}(p') \slashed{p}'] u(p) - \bar{u}(p') [\slashed{p} u(p)] \nonumber \\
    &= M \bar{u}(p') u(p) - M \bar{u}(p') u(p) \equiv 0.
    \end{align}
    ¡El término de Dirac se anula idénticamente en virtud de las ecuaciones de movimiento libres!

    \item \textbf{Término de convección ($q_\mu (p+p')^\mu$):}
    \begin{equation}
    q_\mu (p + p')^\mu = (p' - p)_\mu (p' + p)^\mu = p'^2 - p^2 = M^2 - M^2 \equiv 0.
    \end{equation}
    ¡El término de convección se anula idénticamente por conservación de la masa invariante on-shell!

    \item \textbf{Término tensorial ($q_\mu \sigma^{\mu\nu} q_\nu$):}
    Aunque la descomposición original de la Ec.~\eqref{eq:vertice_general_abc} usa $(p+p')^\mu$, al emplear la identidad de Gordon surge el término $\sigma^{\mu\nu} q_\nu$. Su contracción con $q_\mu$ es:
    \begin{equation}
    q_\mu (i \sigma^{\mu\nu} q_\nu) = i q_\mu q_\nu \sigma^{\mu\nu}.
    \end{equation}
    Demostremos formalmente que esta contracción es cero. Dado que $\mu$ y $\nu$ son índices mudos de suma, reetiquetamos $\mu \leftrightarrow \nu$:
    \begin{equation}
    \mathcal{I}_{\text{tens}} \equiv q_\mu q_\nu \sigma^{\mu\nu} = q_\nu q_\mu \sigma^{\nu\mu}.
    \end{equation}
    Puesto que las componentes del cuadrimomento conmutan ($q_\nu q_\mu = q_\mu q_\nu$) y el tensor de espín es estrictamente antisimétrico ($\sigma^{\nu\mu} = -\sigma^{\mu\nu}$):
    \begin{equation}
    \mathcal{I}_{\text{tens}} = q_\mu q_\nu (-\sigma^{\mu\nu}) = - q_\mu q_\nu \sigma^{\mu\nu} = - \mathcal{I}_{\text{tens}}.
    \end{equation}
    Sumando $\mathcal{I}_{\text{tens}}$ en ambos miembros:
    \begin{equation}
    2 \mathcal{I}_{\text{tens}} = 0 \implies \mathcal{I}_{\text{tens}} \equiv 0 \implies q_\mu (i \sigma^{\mu\nu} q_\nu) = 0.
    \end{equation}
    Conviene notar la naturaleza de este último resultado, que lo distingue de los dos anteriores: se ha obtenido contrayendo un tensor simétrico ($q_\mu q_\nu$) con uno antisimétrico ($\sigma^{\mu\nu}$), sin utilizar en ningún momento la capa de masa ni condición cinemática alguna. Es una anulación puramente estructural y, por tanto, válida en cualquier circunstancia, dentro o fuera de la capa de masa.

    \item \textbf{Término longitudinal ($q_\mu q^\mu$):}
    \begin{equation}
    q_\mu [C(q^2) q^\mu \bar{u}(p') u(p)] = C(q^2) q^2 \bar{u}(p') u(p).
    \end{equation}
    Este es el único de los cuatro que \emph{no} se anula por sí solo, y en esa asimetría reside todo el argumento.
\end{enumerate}

Sustituyendo estas cuatro contracciones en la condición de Ward, Ec.~\eqref{eq:ward_condition}:
\begin{equation}
A(q^2) \cdot 0 + B(q^2) \cdot 0 + C(q^2) q^2 \bar{u}(p') u(p) = 0 \implies C(q^2) q^2 \bar{u}(p') u(p) = 0.
\end{equation}
En la dispersión física elástica, la transferencia de momento es estrictamente espacial y no nula ($q^2 = -Q^2 < 0$), y la densidad escalar de espinores en el límite físico no se anula ($\bar{u}(p') u(p) \approx 2M \neq 0$). Para que la igualdad se satisfaga idénticamente en todo el dominio cinemático, el coeficiente escalar debe anularse sin excepción:
\begin{equation}
\dest{C(q^2) \equiv 0}
\label{eq:C_zero_exact}
\end{equation}
Nótese la fuerza del argumento: no se ha demostrado que $C$ sea pequeño, ni que resulte despreciable en algún régimen, sino que se anula idénticamente para todo valor de la transferencia de momento. Los coeficientes $A(q^2)$ y $B(q^2)$, en cambio, quedaron multiplicados por cero en la condición de Ward y por tanto permanecen enteramente libres: la conservación de la corriente no dice nada sobre ellos, y son ellos los dos grados de libertad que sobreviven.

Sustituyendo la Identidad de Gordon de la Ec.~\eqref{eq:gordon_identity_full} para eliminar el término de convección $(p+p')^\mu$ en favor del tensor dipolar de Pauli $\sigma^{\mu\nu} q_\nu$, el operador de vértice electromagnético de Dirac-Pauli queda formalmente establecido:
\begin{equation}
\boxed{\Gamma^\mu(q^2) = \gamma^\mu F_1(q^2) + \frac{i \sigma^{\mu\nu} q_\nu}{2M} F_2(q^2)}
\label{eq:vertice_dirac_pauli_final}
\end{equation}
donde $F_1(q^2)$ es el \textbf{factor de forma de Dirac} y $F_2(q^2)$ es el \textbf{factor de forma de Pauli}.

Este es el primer resultado verdaderamente central del artículo, y merece leerse despacio. Lo que la Ec.~\eqref{eq:vertice_dirac_pauli_final} afirma es que el acoplamiento electromagnético de cualquier partícula de espín $1/2$, elemental o compuesta, conocida o por descubrir, tiene forzosamente esta forma y contiene exactamente dos funciones escalares desconocidas. La región difusa del vértice inferior de la Figura~\ref{fig:feynman_tree} ha quedado caracterizada sin haberla resuelto. Conviene ser explícito sobre lo que ha costado: nada, salvo tres supuestos que ninguna teoría razonable del nucleón pondría en duda---que el protón tiene espín $1/2$, que la interacción electromagnética conserva paridad y que la carga eléctrica se conserva. No ha intervenido la cromodinámica cuántica, ni un modelo de quarks, ni una función de onda hadrónica. El precio de esa generalidad es igualmente claro: la teoría nos dice cuántos factores de forma hay y cómo entran en la amplitud, pero no cuánto valen. Determinar $F_1(Q^2)$ y $F_2(Q^2)$ es una tarea experimental, y construir la fórmula que permite medirlos es el objeto del resto de este trabajo.

\subsection{Consistencia Dimensional, Relaciones de Base y Verificaciones Físicas}
\label{subsec:consistencia_bases}
Queda por justificar un detalle de la Ec.~\eqref{eq:vertice_dirac_pauli_final} que suele pasar desapercibido: ¿de dónde sale el factor $1/(2M)$ que acompaña al término de Pauli, y por qué esa escala de masa en concreto? La respuesta tiene dos capas, una obligatoria y otra convencional, y separarlas evita confusiones conceptuales.

En unidades naturales ($\hbar = c = 1$), la acción es adimensional. La matriz de Dirac $\gamma^\mu$ es adimensional ($[\gamma^\mu] = [M]^0$), mientras que el cuadrimomento transferido posee dimensión de masa ($[q_\nu] = [M]^1$). Por consiguiente, el operador tensorial $i\sigma^{\mu\nu} q_\nu$ tiene dimensiones de masa $[M]^1$. Para que ambos factores de forma $F_1$ y $F_2$ sean funciones genuinamente adimensionales, es forzoso introducir un factor normalizador con dimensiones de masa inversa $[M]^{-1}$ en el término dipolar de Pauli. La elección de la escala $1/(2M)$ coincide con la definición canónica del magnetón nuclear $\mu_N = e/(2M)$. Con esta normalización, el valor estático $F_2(0)$ coincide exactamente con el momento magnético anómalo medido en unidades de $\mu_N$.

Aunque la base de Dirac-Pauli $\{F_1, F_2\}$ surge de manera natural de la teoría cuántica de campos, no es la más conveniente para el análisis experimental ni para la interpretación intuitiva en el espacio de configuración. En 1962, Ernst J. Sachs y colaboradores introdujeron las combinaciones lineales desacopladas conocidas como los \textbf{factores de forma de Sachs} \cite{sachs1962}:
\begin{align}
G_E(Q^2) &\equiv F_1(Q^2) - \tau F_2(Q^2), \label{eq:sachs_ge_def} \\
G_M(Q^2) &\equiv F_1(Q^2) + F_2(Q^2), \label{eq:sachs_gm_def}
\end{align}
cuyas transformaciones inversas están dadas por:
\begin{equation}
F_1(Q^2) = \frac{G_E(Q^2) + \tau G_M(Q^2)}{1 + \tau}, \qquad F_2(Q^2) = \frac{G_M(Q^2) - G_E(Q^2)}{1 + \tau}.
\label{eq:sachs_inversas}
\end{equation}

Para fijar con claridad la distinción y el rol operativo de cada una de estas representaciones, la Tabla~\ref{tab:comparacion_factores_forma} presenta un despiece comparativo de las bases de Dirac-Pauli y de Sachs.

\begin{table}[t]
\centering
\small
\renewcommand{\arraystretch}{1.3}
\caption{Comparativa sistemática entre las representaciones de factores de forma del nucleón. La base covariante de Dirac-Pauli surge directamente del lagrangiano y de las reglas de Feynman en QED, mientras que la base de Sachs desacopla los canales multipolares eliminando los términos de interferencia en la sección eficaz y admite una interpretación espacial directa en el marco de Breit.}
\label{tab:comparacion_factores_forma}
\resizebox{0.98\textwidth}{!}{%
\begin{tabular}{p{3.2cm} p{6.5cm} p{6.5cm}}
\toprule
\textbf{Criterio / Propiedad} & \textbf{Base de Dirac-Pauli $\{F_1(Q^2), F_2(Q^2)\}$} & \textbf{Base de Sachs $\{G_E(Q^2), G_M(Q^2)\}$} \\
\midrule
\textbf{Significado Microscópico} & $F_1$: Corriente vectorial / Dirac puntual vestida.\newline $F_2$: Densidad tensorial de Pauli (anomalía magnética). & $G_E$: Distribución neta de carga eléctrica.\newline $G_M$: Distribución de corrientes magnéticas. \\
\addlinespace[0.6ex]
\textbf{Estructura en el Vértice} & $\Gamma^\mu = \gamma^\mu F_1 + \dfrac{i\sigma^{\mu\nu}q_\nu}{2M} F_2$ (covarianza manifiesta). & Combinación lineal dependiente del parámetro cinemático de retroceso $\tau = Q^2/(4M^2)$. \\
\addlinespace[0.6ex]
\textbf{Valores Estáticos ($Q^2 \to 0$)} & \textbf{Protón:} $F_1^p(0) = 1$, \; $F_2^p(0) = \kappa_p \approx +1.7928$.\newline \textbf{Neutrón:} $F_1^n(0) = 0$, \; $F_2^n(0) = \kappa_n \approx -1.9130$. & \textbf{Protón:} $G_E^p(0) = 1$, \; $G_M^p(0) = \mu_p \approx 2.7928$.\newline \textbf{Neutrón:} $G_E^n(0) = 0$, \; $G_M^n(0) = \mu_n \approx -1.9130$. \\
\addlinespace[0.6ex]
\textbf{Interferencias en $d\sigma/d\Omega$} & Contiene términos cruzados proporcionales a $F_1 F_2$ en el sistema del laboratorio. & \textbf{Desacoplamiento cuadrático exacto:} solo aparecen $G_E^2$ y $G_M^2$, sin interferencias cruzadas $G_E G_M$. \\
\addlinespace[0.6ex]
\textbf{Interpretación Espacial} & Opacada en espacio de coordenadas por términos mixtos de acoplamiento espín-órbita relativistas. & Transformada de Fourier directa de densidades tridimensionales $\rho_E(\mathbf{r})$ y $\rho_M(\mathbf{r})$ en el marco de Breit. \\
\bottomrule
\end{tabular}%
}%
\end{table}

Antes de avanzar hacia el cálculo de las trazas, someteremos las parametrizaciones obtenidas a tres \textbf{verificaciones físicas multidimensionales} (\textit{sanity checks}):
\begin{enumerate}
    \item \textbf{Límite de fermión puntual de Dirac:} Si el protón careciera de estructura interna y fuera un fermión elemental puntual idéntico al electrón, su corriente dipolar anómala se anularía idénticamente ($F_2 \equiv 0$) y su carga puntual sería fija ($F_1 \equiv 1$). En dicho caso:
    \begin{equation}
    G_E(Q^2) = 1 - \tau(0) = 1, \qquad G_M(Q^2) = 1 + 0 = 1, \quad \forall Q^2.
    \end{equation}
    Sustituyendo estos valores en la Ec.~\eqref{eq:rosenbluth_master_intro}, el corchete de estructura se reduce a:
    \begin{equation}
    \frac{1 + \tau}{1+\tau} + 2\tau \tan^2\left(\frac{\theta}{2}\right) = 1 + 2\tau \tan^2\left(\frac{\theta}{2}\right),
    \end{equation}
    recuperando exactamente la sección eficaz de Dirac para un blanco puntual con espín y retroceso, como se consignó en la tercera fila de la Tabla~\ref{tab:modelos_dispersion}.
    \item \textbf{Límite estático de baja energía ($Q^2 \to 0$):} Cuando $Q^2 \to 0$, el parámetro de retroceso se anula ($\tau \to 0$). Las Ecs.~\eqref{eq:sachs_ge_def} y \eqref{eq:sachs_gm_def} se reducen a:
    \begin{equation}
    G_E(0) = F_1(0) = \frac{Q_{\text{total}}}{e}, \qquad G_M(0) = F_1(0) + F_2(0) = \frac{Q_{\text{total}}}{e} + \kappa = \frac{\mu}{\mu_N}.
    \end{equation}
    Para el protón, $G_E^p(0) = 1$ y $G_M^p(0) = 1 + 1.7928 = 2.7928$. Para el neutrón, $G_E^n(0) = 0$ y $G_M^n(0) = 0 - 1.9130 = -1.9130$. El formalismo reproduce de manera exacta los momentos magnéticos estáticos medidos experimentalmente.
    \item \textbf{Resolución de la aparente paradoja del neutrón:} Que el neutrón posea $G_E^n(0) = 0$ y simultáneamente un momento magnético $G_M^n(0) \neq 0$ deja de ser una contradicción en la teoría de campos. La carga neta integrada en todo el espacio es nula ($\int \rho_n(\mathbf{r}) d^3\mathbf{r} = 0$), pero al estar formado por quarks de valencia $u(+2/3)$ y dos $d(-1/3)$, la densidad de carga local $\rho_n(r)$ es positiva en el centro (dominada por el quark $u$) y negativa en la periferia (dominada por los quarks $d$). Al ponderar con $r^2$, el radio cuadrático medio resulta negativo ($\langle r_n^2 \rangle < 0$), lo que genera corrientes magnetizantes netas no compensadas.
\end{enumerate}

Habiendo demostrado que la invariancia de Lorentz y la conservación de corriente restringen el acoplamiento hadrónico a exactamente dos factores de forma escalares ($F_1$ y $F_2$), disponemos formalmente de la amplitud de transición invariante $\mathcal{M}$. No obstante, esta amplitud contiene espinores de Dirac libres que no son directamente observables en el laboratorio: un detector no mide espinores, mide cuentas por unidad de ángulo sólido. Para conectar este elemento de matriz con una sección eficaz diferencial medible en un detector real, debemos promediar sobre los estados de polarización iniciales, sumar sobre los estados finales, evaluar algebraicamente las trazas tensoriales de Dirac resultantes ($L_{\mu\nu}$ y $W^{\mu\nu}$) y resolver la cinemática relativista del retroceso elástico del protón en el sistema de referencia del laboratorio. Ese es el trabajo de la sección siguiente, la más intensiva en álgebra de todo el artículo y, por esa misma razón, la que más conviene abordar sabiendo de antemano hacia dónde se dirige.
\section{Cinemática Relativista y Trazas Tensoriales de Dirac}
\label{sec:cinematica_trazas}

Esta es, con diferencia, la sección más intensiva en álgebra de todo el artículo, y precisamente por eso conviene entrar en ella con el mapa desplegado. El objetivo es una única cantidad: el módulo al cuadrado de la amplitud, $\overline{|\mathcal{M}|^2}$, promediado sobre las polarizaciones iniciales y sumado sobre las finales, porque es esa y no la amplitud lo que un detector puede medir. El cálculo se descompone en cuatro etapas de dificultad creciente pero de lógica muy simple. Primero (Sección~\ref{sec:cinematica_trazas}.1) fijamos la cinemática: qué relaciones ligan la energía incidente $E$, la dispersada $E'$, el ángulo $\theta$ y la transferencia $Q^2$ en el sistema del laboratorio. Después separamos el cálculo cuántico en dos mitades independientes, una por cada extremo de la línea de fotón de la Figura~\ref{fig:feynman_tree}: el \emph{tensor leptónico} $L_{\mu\nu}$, que describe el vértice del electrón y es completamente calculable desde QED~\cite{halzen1984,thomson2013}, y el \emph{tensor hadrónico} $W^{\mu\nu}$, donde entran los dos factores de forma desconocidos~\cite{das2003,henley2007}. Finalmente los contraemos entre sí (Sección~\ref{sec:cinematica_trazas}.4).

Conviene anunciar desde ahora el desenlace, porque no es el que cabría esperar y es lo que da sentido a la sección siguiente. El resultado final de este cálculo, la Ec.~\eqref{eq:contraccion_invariante_f1_f2}, será correcto y exacto, pero contendrá un término proporcional al producto $F_1 F_2$ que mezcla la carga con la magnetización de forma inseparable. Es decir: habremos calculado bien y, aun así, no podremos extraer $F_1$ y $F_2$ por separado de una medida. Ese fracaso parcial no es un error de cálculo sino un hecho físico sobre el sistema de laboratorio, y entenderlo es lo que motivará el cambio de marco de referencia de la Sección~\ref{sec:marco_breit}.

\subsection{Sistema del Laboratorio y Retroceso Elástico}
Antes de calcular nada cuántico hay que fijar la geometría del proceso, y la pregunta de esta subsección es puramente cinemática. En un experimento real se dispara un haz de energía conocida $E$ contra un blanco de hidrógeno en reposo y se detecta el electrón dispersado midiendo su energía $E'$ y su ángulo $\theta$. Pero en una colisión \emph{elástica} esas tres cantidades no son independientes: el requisito de que el protón salga siendo un protón---es decir, que conserve su masa---impone una relación exacta entre ellas. Determinar esa relación es lo que hacemos aquí, y su consecuencia práctica es considerable: bastará medir el ángulo para conocer toda la cinemática del suceso. La Figura~\ref{fig:cinematica_lab} fija las convenciones.

\begin{figure}[t]
\centering
\begin{tikzpicture}[>=stealth, line width=0.7pt, font=\small]

% --- eje z (dirección del haz) ---
\draw[gray!55, dashed, line width=0.5pt] (-3.60,0) -- (3.55,0);
\node[gray!70, font=\scriptsize, anchor=west] at (3.60,0) {eje $z$};

% --- electrón incidente ---
\draw[->] (-3.40,0) -- (-1.85,0);
\draw     (-1.85,0) -- (-0.32,0);
\node[anchor=south, font=\footnotesize, inner sep=3pt] at (-2.30,0.04) {$e^-(k)$,\; $E$};

% --- protón blanco en reposo ---
\shade[inner color=orange!58, outer color=orange!8] (0,0) circle (0.34);
\draw[orange!65!black, line width=0.5pt] (0,0) circle (0.34);
\node[anchor=north east, align=right, font=\footnotesize, inner sep=2pt] at (-0.30,-0.44)
      {$p^\mu = (M,\mathbf{0})$\\[0.2ex] {\scriptsize\color{gray}blanco en reposo}};

% --- electrón dispersado ---
\draw[->] (34:0.42) -- (34:2.10);
\draw     (34:2.10) -- (34:3.35);
\node[anchor=west, font=\footnotesize, inner sep=3pt] at (34:3.40) {$e^-(k')$,\; $E'$};

% --- protón de retroceso ---
\draw[->, line width=1.1pt] (-26:0.42) -- (-26:1.75);
\draw[line width=1.1pt]     (-26:1.75) -- (-26:2.75);
\node[anchor=west, font=\footnotesize, inner sep=3pt] at (-26:2.80) {$p(p')$};

% --- ángulo de dispersión ---
\draw[gray!75, line width=0.5pt] (1.15,0) arc (0:34:1.15);
\node[gray!85, font=\footnotesize] at (17:1.42) {$\theta$};

% --- caja de relaciones cinemáticas ---
\node[anchor=north west, font=\footnotesize, align=left, inner sep=0pt] at (4.95, 1.72)
 {Transferencia de momento:\\[0.7ex]
  $Q^2 = 4EE'\sin^2\!\left(\dfrac{\theta}{2}\right)$\\[1.4ex]
  Retroceso elástico:\\[0.7ex]
  $E' = \dfrac{E}{1 + \dfrac{2E}{M}\sin^2\!\left(\dfrac{\theta}{2}\right)}$};

% --- el rasgo decisivo del marco de laboratorio ---
\node[anchor=north west, font=\footnotesize, align=left, inner sep=4pt,
      fill=colordest, draw=colorborde, line width=0.4pt] at (4.85,-1.15)
 {$q^0 = E - E' \neq 0$\\[0.3ex] \scriptsize el blanco absorbe energía};

\end{tikzpicture}
\caption{Cinemática de la dispersión elástica $e^- p \to e^- p$ en el sistema del laboratorio, con el protón inicialmente en reposo. El electrón incide con energía $E$ a lo largo del eje $z$, es dispersado con energía $E'$ formando un ángulo $\theta$, y el protón retrocede absorbiendo el cuadrimomento transferido $q^\mu = k^\mu - k'^\mu$. La condición de elasticidad---que el protón final siga teniendo masa $M$, esto es, $p'^2 = M^2$---liga las tres cantidades $E$, $E'$ y $\theta$ mediante la relación de retroceso indicada a la derecha, formalmente idéntica a la fórmula del efecto Compton con el protón en el papel del electrón. En consecuencia, fijada la energía del haz, \emph{una sola} variable es independiente: medir el ángulo de dispersión determina simultáneamente $E'$ y $Q^2$. El rasgo señalado en el recuadro inferior es el que conviene retener para más adelante: en este marco de referencia el protón absorbe energía, de modo que la componente temporal de la transferencia no se anula ($q^0 = E - E' \neq 0$). Ese hecho, aparentemente inocuo, será el responsable de que las densidades de carga y de magnetización aparezcan mezcladas al final de esta sección, y la Figura~\ref{fig:marco_breit} mostrará el marco de referencia que lo remedia.}
\label{fig:cinematica_lab}
\end{figure}

En el sistema de laboratorio, el protón se encuentra inicialmente en reposo actuando como blanco estático:
\begin{equation}
p^\mu = (M, \mathbf{0}) = (M, 0, 0, 0) \implies p^2 = p_\mu p^\mu = M^2 - |\mathbf{0}|^2 = M^2.
\end{equation}
Para electrones incidentes en el régimen ultra-relativista de alta energía ($E \gg m_e \approx 0.511\text{ MeV} \approx 0$), la relación de dispersión se simplifica a $|\mathbf{k}| = \sqrt{E^2 - m_e^2} \approx E$. Alineando el eje $z$ en la dirección del haz incidente y situando el plano de dispersión en el plano $xz$:
\begin{align}
k^\mu &= (E, 0, 0, E), \\
k'^\mu &= (E', E'\sin\theta, 0, E'\cos\theta).
\end{align}
El producto escalar invariante de Lorentz entre los cuadrimomentos leptónicos inicial y final es:
\begin{align}
k \cdot k' &= k^0 k'^0 - \mathbf{k} \cdot \mathbf{k}' = E E' - (0 \cdot E'\sin\theta + 0 \cdot 0 + E \cdot E'\cos\theta) \nonumber \\
&= E E' - E E' \cos\theta = E E' (1 - \cos\theta).
\end{align}
Aplicando la identidad trigonométrica fundamental del ángulo medio deducida de $\cos(2\alpha) = \cos^2\alpha - \sin^2\alpha = 1 - 2\sin^2\alpha \implies 1 - \cos\theta = 2\sin^2(\theta/2)$:
\begin{equation}
\dest{k \cdot k' = 2 E E' \sin^2\left(\frac{\theta}{2}\right)}
\label{eq:k_dot_kprime}
\end{equation}
La aparición del ángulo mitad, $\theta/2$, no es un capricho notacional: acompañará a todas las expresiones que siguen y terminará estructurando la propia fórmula de Rosenbluth, donde $\cos^2(\theta/2)$ gobierna la respuesta eléctrica y $\tan^2(\theta/2)$ la magnética.

El cuadrimomento transferido por el fotón virtual mediador es $q^\mu = k^\mu - k'^\mu$. Evaluando su cuadrado invariante:
\begin{align}
q^2 &= (k - k')^2 = k^2 - 2 k \cdot k' + k'^2 = m_e^2 - 2(k \cdot k') + m_e^2 \nonumber \\
&\approx 0 - 2 \left[ 2 E E' \sin^2\left(\frac{\theta}{2}\right) \right] + 0 = - 4 E E' \sin^2\left(\frac{\theta}{2}\right) \equiv -Q^2 \le 0.
\label{eq:q2_def}
\end{align}
Definimos $Q^2 \equiv -q^2 = 4 E E' \sin^2(\theta/2) > 0$ como la transferencia espacial de momento al cuadrado. El signo negativo de $q^2$ no es un accidente sino una propiedad general de la dispersión: el fotón intercambiado es de tipo espacial, no puede materializarse como una partícula real, y su «masa virtual» negativa es lo que mide la resolución espacial del sondeo.

Imponiendo que el protón en el estado final sea una partícula libre en su capa de masa física ($p'^2 = M^2$):
\begin{align}
p'^2 &= (p + q)^2 = p^2 + 2 p \cdot q + q^2 \nonumber \\
M^2 &= M^2 + 2 p \cdot q - Q^2 \implies 2 p \cdot q = Q^2.
\label{eq:onshell_recoil}
\end{align}
Esta única condición---el protón sigue siendo un protón---es la que va a producir toda la relación cinemática. Evaluemos el producto escalar $p \cdot q$ en el sistema del laboratorio:
\begin{align}
p \cdot q &= p_\mu q^\mu = p_0 q^0 - \mathbf{p} \cdot \mathbf{q} = M (E - E') - \mathbf{0} \cdot (\mathbf{k} - \mathbf{k}') \nonumber \\
&= M (E - E').
\end{align}
Sustituyendo en la condición on-shell de la Ec.~\eqref{eq:onshell_recoil}:
\begin{equation}
2 M (E - E') = Q^2 = 4 E E' \sin^2\left(\frac{\theta}{2}\right).
\end{equation}
Dividiendo ambos miembros de la ecuación entre la cantidad estrictamente positiva $2 M E E'$:
\begin{equation}
\frac{2 M (E - E')}{2 M E E'} = \frac{4 E E' \sin^2(\theta/2)}{2 M E E'} \implies \frac{E - E'}{E E'} = \frac{2}{M} \sin^2\left(\frac{\theta}{2}\right).
\end{equation}
Descomponiendo la fracción en el miembro izquierdo:
\begin{equation}
\frac{E}{E E'} - \frac{E'}{E E'} = \frac{1}{E'} - \frac{1}{E} = \frac{2}{M} \sin^2\left(\frac{\theta}{2}\right).
\end{equation}
Sumando $1/E$ en ambos lados y extrayendo factor común $1/E$:
\begin{equation}
\frac{1}{E'} = \frac{1}{E} + \frac{2}{M} \sin^2\left(\frac{\theta}{2}\right) = \frac{1}{E} \left[ 1 + \frac{2E}{M} \sin^2\left(\frac{\theta}{2}\right) \right].
\end{equation}
Invirtiendo algebraicamente ambos miembros, obtenemos la fórmula analítica exacta del \textbf{retroceso cinemático relativista de Compton}:
\begin{equation}
\dest{E' = \frac{E}{1 + \frac{2E}{M} \sin^2\left(\frac{\theta}{2}\right)}}
\label{eq:recoil_energy_exact}
\end{equation}
La estructura de esta relación merece un comentario. Es formalmente idéntica a la fórmula del efecto Compton que el lector ya conoce, con el protón desempeñando el papel que allí tenía el electrón. En el límite de blanco infinitamente pesado ($M \to \infty$) el denominador tiende a la unidad y $E' \to E$: la dispersión sería elástica sin pérdida de energía, como el rebote de una pelota contra una pared. El término que rompe esa idealización, $(2E/M)\sin^2(\theta/2)$, mide cuánta energía se lleva el protón al retroceder, y crece tanto con la energía del haz como con el ángulo. La consecuencia experimental es la anunciada al principio: fijada $E$, medir $\theta$ determina $E'$ y por tanto $Q^2$, de modo que cada suceso registrado lleva su propia etiqueta cinemática.

Definimos el parámetro cinemático adimensional de retroceso de Sachs:
\begin{equation}
\dest{\tau \equiv \frac{Q^2}{4M^2} = \frac{2M(E - E')}{4M^2} = \frac{E - E'}{2M}}
\label{eq:tau_def}
\end{equation}
Las dos igualdades de la derecha muestran que $\tau$ admite dos lecturas equivalentes: como transferencia de momento medida en unidades de la masa del nucleón, o como fracción de energía cedida al blanco. Ambas serán útiles, y la segunda hace transparente por qué $\tau$ acabará controlando la importancia relativa de los efectos magnéticos: son los que requieren que el blanco se mueva.

\subsection{Cálculo Detallado del Tensor Leptónico (\texorpdfstring{$L_{\mu\nu}$}{L\_munu})}
Con la cinemática resuelta pasamos a la parte cuántica. El problema es el siguiente: la amplitud $\mathcal{M}$ de la Ec.~\eqref{eq:feynman_amplitude_full} contiene espinores, que dependen de las orientaciones de espín de las cuatro partículas involucradas. Un experimento ordinario no prepara ni mide esas orientaciones: el haz no está polarizado y el detector no distingue el espín del electrón saliente. Lo que se mide, por tanto, es un promedio sobre los espines iniciales y una suma sobre los finales. Esta operación tiene una virtud algebraica decisiva que es el motor de todo el cálculo: al sumar sobre espines, los productos de espinores se convierten en \emph{trazas de productos de matrices}, y una traza es un número, no un objeto con índices espinoriales ocultos. El lector que no haya manejado trazas de Dirac puede tomarlas literalmente por lo que son---la suma de los elementos diagonales de un producto de matrices $4\times 4$---y consultar el Apéndice~\ref{sec:apendice_dirac}, donde se demuestran desde cero las dos únicas identidades que necesitaremos aquí: que las trazas de un número impar de matrices gamma se anulan, y el valor de la traza de cuatro de ellas.

El cálculo se separa además en dos mitades independientes, una por cada vértice, porque el promedio sobre espines del electrón y el del protón son sumas distintas y factorizan. Empezamos por la mitad fácil: la del electrón, donde el vértice es exactamente $\gamma^\mu$ y no hay ninguna función desconocida.

A partir de la amplitud de Feynman invariante de la Ec.~\eqref{eq:feynman_amplitude_full}~\cite{halzen1984,thomson2013,das2003}:
\begin{equation}
\mathcal{M} = \frac{e^2}{q^2} \left[ \bar{u}(k', s') \gamma_\mu u(k, s) \right] \left[ \bar{u}(p', r') \Gamma^\mu u(p, r) \right],
\end{equation}
el módulo al cuadrado promediado sobre los estados de espín iniciales del electrón ($s$) y del protón ($r$) (pesos estadísticos $1/2 \times 1/2 = 1/4$) y sumado sobre los espines finales ($s', r'$) es:
\begin{equation}
\overline{|\mathcal{M}|^2} = \frac{1}{4} \sum_{s, s', r, r'} |\mathcal{M}|^2 = \frac{e^4}{q^4} \left( \frac{1}{2} \sum_{s, s'} j_\mu^e (j_\nu^e)^* \right) \left( \frac{1}{2} \sum_{r, r'} j_p^\mu (j_p^\nu)^* \right) \equiv \frac{e^4}{Q^4} L_{\mu\nu} W^{\mu\nu}.
\end{equation}
La última igualdad es la factorización anunciada: toda la información del vértice leptónico queda en $L_{\mu\nu}$ y toda la del hadrónico en $W^{\mu\nu}$, de modo que pueden calcularse por separado y contraerse al final.

Evaluemos minuciosamente el tensor leptónico $L_{\mu\nu}$. El primer paso es reescribir la corriente conjugada en una forma que permita reconocer la traza, y para ello se usa la propiedad de hermiticidad de las matrices gamma:
\begin{align}
[j_\nu^e]^* &= [\bar{u}(k', s') \gamma_\nu u(k, s)]^* = [u^\dagger(k', s') \gamma^0 \gamma_\nu u(k, s)]^\dagger \nonumber \\
&= u^\dagger(k, s) \gamma_\nu^\dagger (\gamma^0)^\dagger u(k', s') = u^\dagger(k, s) \gamma_\nu^\dagger \gamma^0 u(k', s').
\end{align}
Insertando la identidad $\mathbb{I} = (\gamma^0)^2$ entre $u^\dagger$ y $\gamma_\nu^\dagger$, y usando la relación de hermiticidad $\gamma^0 \gamma_\nu^\dagger \gamma^0 = \gamma_\nu$:
\begin{equation}
[j_\nu^e]^* = u^\dagger(k, s) \gamma^0 (\gamma^0 \gamma_\nu^\dagger \gamma^0) u(k', s') = \bar{u}(k, s) \gamma_\nu u(k', s').
\end{equation}
El resultado es que el conjugado de la corriente tiene la misma forma que la corriente, con los espinores intercambiados. Eso permite encadenar los cuatro espinores en un ciclo cerrado, que es justamente lo que produce una traza. Escribiendo las componentes espinoriales con índices explícitos $\alpha, \beta, \gamma, \delta \in \{1, 2, 3, 4\}$:
\begin{align}
L_{\mu\nu} &= \frac{1}{2} \sum_{s, s'} [\bar{u}(k', s') \gamma_\mu u(k, s)] [\bar{u}(k, s) \gamma_\nu u(k', s')] \nonumber \\
&= \frac{1}{2} \sum_{s, s'} \bar{u}_\alpha(k', s') (\gamma_\mu)_{\alpha\beta} u_\beta(k, s) \, \bar{u}_\gamma(k, s) (\gamma_\nu)_{\gamma\delta} u_\delta(k', s').
\end{align}
Reagrupando los factores escalares:
\begin{equation}
L_{\mu\nu} = \frac{1}{2} (\gamma_\mu)_{\alpha\beta} \left( \sum_s u_\beta(k, s) \bar{u}_\gamma(k, s) \right) (\gamma_\nu)_{\gamma\delta} \left( \sum_{s'} u_\delta(k', s') \bar{u}_\alpha(k', s') \right).
\end{equation}
Aplicando las relaciones de completitud de espinores relativistas libres para fermiones de masa $m_e \to 0$:
\begin{equation}
\sum_s u(k, s) \bar{u}(k, s) = \slashed{k} + m_e \approx \slashed{k}, \qquad \sum_{s'} u(k', s') \bar{u}(k', s') = \slashed{k}' + m_e \approx \slashed{k}'.
\end{equation}
Aquí es donde la suma sobre espines hace su trabajo: convierte dos objetos espinoriales en una matriz ordinaria construida solo con cuadrimomentos. Sustituyendo:
\begin{equation}
L_{\mu\nu} = \frac{1}{2} (\gamma_\mu)_{\alpha\beta} (\slashed{k})_{\beta\gamma} (\gamma_\nu)_{\gamma\delta} (\slashed{k}')_{\delta\alpha} = \frac{1}{2} (\gamma_\mu \slashed{k} \gamma_\nu \slashed{k}')_{\alpha\alpha}.
\end{equation}
La contracción de los índices primero y último representa la definición de la traza de Dirac:
\begin{equation}
L_{\mu\nu} = \frac{1}{2} \text{Tr} \left[ \gamma_\mu \slashed{k} \gamma_\nu \slashed{k}' \right] = \frac{1}{2} \text{Tr} \left[ \slashed{k}' \gamma_\mu \slashed{k} \gamma_\nu \right],
\end{equation}
donde hemos aplicado la propiedad de invarianza cíclica de la traza: $\text{Tr}[A B C D] = \text{Tr}[D A B C]$. El problema cuántico ha quedado reducido a un problema de álgebra de matrices, y lo que resta es puramente mecánico.

Para evaluar analíticamente esta traza, escribimos $\slashed{k}' = k'^\alpha \gamma_\alpha$ y $\slashed{k} = k^\beta \gamma_\beta$:
\begin{equation}
\text{Tr}[\slashed{k}' \gamma_\mu \slashed{k} \gamma_\nu] = k'^\alpha k^\beta \text{Tr}[\gamma_\alpha \gamma_\mu \gamma_\beta \gamma_\nu].
\end{equation}
Aplicando el teorema fundamental de trazas de cuatro matrices de Dirac (demostrado formalmente en el Apéndice~\ref{sec:apendice_dirac}):
\begin{equation}
\text{Tr}[\gamma_\alpha \gamma_\mu \gamma_\beta \gamma_\nu] = 4 \left( g_{\alpha\mu} g_{\beta\nu} - g_{\alpha\beta} g_{\mu\nu} + g_{\alpha\nu} g_{\beta\mu} \right).
\end{equation}
Contrayendo con $k'^\alpha k^\beta$:
\begin{align}
\text{Tr}[\slashed{k}' \gamma_\mu \slashed{k} \gamma_\nu] &= 4 \left[ (k'^\alpha g_{\alpha\mu})(k^\beta g_{\beta\nu}) - (k'^\alpha k^\beta g_{\alpha\beta}) g_{\mu\nu} + (k'^\alpha g_{\alpha\nu})(k^\beta g_{\beta\mu}) \right] \nonumber \\
&= 4 \left[ k'_\mu k_\nu - (k' \cdot k) g_{\mu\nu} + k'_\nu k_\mu \right] = 4 \left[ k_\mu k'_\nu + k'_\mu k_\nu - (k \cdot k') g_{\mu\nu} \right].
\end{align}
Multiplicando por el factor prefactor $1/2$:
\begin{equation}
\dest{L_{\mu\nu} = 2 \left( k_\mu k'_\nu + k'_\mu k_\nu - (k \cdot k') g_{\mu\nu} \right)}
\label{eq:L_munu_exact}
\end{equation}

Este es el tensor leptónico completo, y conviene leer lo que dice antes de seguir. Es una expresión cerrada construida exclusivamente con los cuadrimomentos del electrón y la métrica: no contiene ninguna función desconocida, ningún parámetro libre y ninguna traza sin evaluar. Toda la información cuántica del vértice electrónico está en esos tres términos. Además es manifiestamente \emph{simétrico} en el intercambio $\mu \leftrightarrow \nu$, y posee una segunda propiedad que no es evidente a simple vista pero que resultará extraordinariamente útil: es transversal a la transferencia de momento. Demostrémoslo, porque vamos a explotarlo de inmediato.

Demostremos algebraicamente que la corriente leptónica satisface rigurosamente la condición de conservación de corriente de Ward ($q^\mu L_{\mu\nu} = 0$). Contrayendo con $q^\mu = k^\mu - k'^\mu$:
\begin{align}
q^\mu L_{\mu\nu} &= 2 (k^\mu - k'^\mu) \left[ k_\mu k'_\nu + k'_\mu k_\nu - (k \cdot k') g_{\mu\nu} \right] \nonumber \\
&= 2 \left[ (k \cdot k) k'_\nu + (k \cdot k') k_\nu - (k \cdot k') k_\nu \right. \nonumber \\
&\quad \left. - (k' \cdot k) k'_\nu - (k' \cdot k') k_\nu + (k \cdot k') k'_\nu \right].
\end{align}
Cancelando los términos opuestos $(k \cdot k') k_\nu - (k \cdot k') k_\nu = 0$ y $-(k' \cdot k) k'_\nu + (k \cdot k') k'_\nu = 0$:
\begin{equation}
q^\mu L_{\mu\nu} = 2 \left[ k^2 k'_\nu - k'^2 k_\nu \right].
\end{equation}
En el límite ultra-relativista $k^2 = k'^2 = m_e^2 \approx 0$:
\begin{equation}
\dest{q^\mu L_{\mu\nu} = 0, \qquad q^\nu L_{\mu\nu} = 0}
\end{equation}
Esta transversalidad es la misma condición de Ward que en la Sección~\ref{sec:vertice_gordon} aniquiló el término longitudinal del vértice hadrónico, reapareciendo ahora en el lado leptónico; no es coincidencia, sino la misma conservación de la corriente electromagnética manifestándose en los dos extremos de la línea de fotón. Su utilidad inmediata es de carácter económico: garantiza que \emph{cualquier} término del tensor hadrónico que sea proporcional a $q^\mu$ o a $q^\nu$ desaparecerá al contraerse con $L_{\mu\nu}$. Podremos por tanto descartar esos términos sin calcularlos, lo que ahorrará una cantidad considerable de álgebra en la subsección siguiente.

\subsection{Cálculo Detallado del Tensor Hadrónico (\texorpdfstring{$W^{\mu\nu}$}{W\^munu})}
Abordamos ahora la mitad difícil. La estructura del cálculo es idéntica a la anterior---promediar sobre espines convierte los espinores del protón en una traza---pero con dos complicaciones. La primera es que el protón no puede tratarse como no masivo, de modo que las relaciones de completitud conservan el término $+M$. La segunda, y más grave, es que el vértice ya no es la matriz simple $\gamma^\mu$ sino el operador de Dirac-Pauli de la Ec.~\eqref{eq:vertice_dirac_pauli_final}, que contiene el tensor $\sigma^{\mu\alpha}$; como cada $\sigma$ es a su vez un conmutador de dos matrices gamma, una evaluación directa generaría trazas de hasta ocho matrices, un cálculo largo y propenso a errores. Afortunadamente disponemos de una herramienta para evitarlo, y es precisamente la identidad de Gordon demostrada en la Sección~\ref{sec:vertice_gordon}: reescribiendo el término de Pauli en términos de $\gamma^\mu$ y del momento $P^\mu$, todas las trazas que aparecen se reducen a cuatro matrices como máximo.

El tensor hadrónico no polarizado se define mediante la traza de espines del nucleón:
\begin{equation}
W^{\mu\nu} = \frac{1}{2} \sum_{r, r'} \left[ \bar{u}(p', r') \Gamma^\mu u(p, r) \right] \left[ \bar{u}(p', r') \Gamma^\nu u(p, r) \right]^* = \frac{1}{2} \text{Tr} \left[ (\slashed{p}' + M) \Gamma^\mu (\slashed{p} + M) \bar{\Gamma}^\nu \right],
\end{equation}
donde el conjugado de Dirac del vértice hadrónico es $\bar{\Gamma}^\nu \equiv \gamma^0 (\Gamma^\nu)^\dagger \gamma^0$.

Para evitar el álgebra tediosa de productos cruzados de hasta ocho matrices gamma provenientes de los conmutadores $\sigma^{\mu\alpha}$, utilizamos la Identidad de Gordon demostrada en la Ec.~\eqref{eq:gordon_gamma_form}:
\begin{equation}
\frac{i\sigma^{\mu\alpha} q_\alpha}{2M} u(p) = \left[ \gamma^\mu - \frac{P^\mu}{2M} \right] u(p), \qquad P^\mu \equiv p'^\mu + p^\mu.
\end{equation}
Sustituyendo esta forma equivalente en el vértice de Dirac-Pauli:
\begin{align}
\Gamma^\mu_{\text{eff}} &= \gamma^\mu F_1(q^2) + \left[ \gamma^\mu - \frac{P^\mu}{2M} \right] F_2(q^2) \nonumber \\
&= [F_1(q^2) + F_2(q^2)] \gamma^\mu - \frac{F_2(q^2)}{2M} P^\mu \equiv G_M(q^2) \gamma^\mu - \frac{F_2(q^2)}{2M} P^\mu.
\label{eq:gamma_eff}
\end{align}
El cambio ha valido la pena: el vértice efectivo ya no contiene tensores $\sigma$, sino solo una matriz gamma y un cuadrivector ordinario multiplicando la identidad. Nótese además que la combinación $F_1 + F_2$ ha aparecido por sí sola, sin que la hayamos buscado, y le hemos dado ya el nombre $G_M$ que justificaremos en la Sección~\ref{sec:marco_breit}; esta aparición espontánea es el primer indicio de que esa combinación, y no $F_1$ o $F_2$ por separado, es la magnitud natural del problema.

Evaluemos el conjugado de Dirac del vértice efectivo:
\begin{align}
\bar{\Gamma}^\nu_{\text{eff}} &= \gamma^0 (\Gamma^\nu_{\text{eff}})^\dagger \gamma^0 = \gamma^0 \left[ G_M (\gamma^\nu)^\dagger - \frac{F_2}{2M} P^\nu \mathbb{I}^\dagger \right] \gamma^0 \nonumber \\
&= G_M (\gamma^0 \gamma^{\nu\dagger} \gamma^0) - \frac{F_2}{2M} P^\nu (\gamma^0 \mathbb{I} \gamma^0) = G_M \gamma^\nu - \frac{F_2}{2M} P^\nu.
\end{align}
¡El vértice efectivo $\Gamma^\mu_{\text{eff}}$ es autoconjugado en el sentido de Dirac! Esto simplifica el cálculo de forma apreciable, pues evita tener que llevar dos expresiones distintas a través de la traza.

Insertando $\Gamma^\mu_{\text{eff}}$ y $\bar{\Gamma}^\nu_{\text{eff}}$ en la traza hadrónica:
\begin{equation}
W^{\mu\nu} = \frac{1}{2} \text{Tr} \left[ (\slashed{p}' + M) \left( G_M \gamma^\mu - \frac{F_2}{2M} P^\mu \right) (\slashed{p} + M) \left( G_M \gamma^\nu - \frac{F_2}{2M} P^\nu \right) \right].
\end{equation}
Expandiendo la multiplicación por linealidad en tres términos tensoriales:
\begin{equation}
W^{\mu\nu} = G_M^2 W_{\text{Dirac}}^{\mu\nu} - \frac{G_M F_2}{2M} \left( P^\nu W_{\text{cruz1}}^\mu + P^\mu W_{\text{cruz2}}^\nu \right) + \frac{F_2^2}{4M^2} P^\mu P^\nu W_{\text{conv}}.
\label{eq:W_three_parts}
\end{equation}
Los tres términos corresponden a las tres formas en que los dos factores del vértice pueden combinarse: dos matrices gamma, una gamma y un momento, o dos momentos. Cada uno da lugar a una traza de naturaleza distinta---tensorial, vectorial y escalar respectivamente---y los evaluamos por separado.

Evaluemos minuciosamente cada una de las tres trazas independientes:

\begin{enumerate}
    \item \textbf{Traza de Dirac pura ($W_{\text{Dirac}}^{\mu\nu}$):}
    \begin{equation}
    W_{\text{Dirac}}^{\mu\nu} \equiv \frac{1}{2} \text{Tr} \left[ (\slashed{p}' + M) \gamma^\mu (\slashed{p} + M) \gamma^\nu \right].
    \end{equation}
    Expandiendo el producto dentro de la traza:
    \begin{equation}
    (\slashed{p}' + M) \gamma^\mu (\slashed{p} + M) \gamma^\nu = \slashed{p}' \gamma^\mu \slashed{p} \gamma^\nu + M \slashed{p}' \gamma^\mu \gamma^\nu + M \gamma^\mu \slashed{p} \gamma^\nu + M^2 \gamma^\mu \gamma^\nu.
    \end{equation}
    Por el teorema de anulación de trazas impares ($\text{Tr}[\gamma^{\mu_1}\dots\gamma^{\mu_{2n+1}}] \equiv 0$):
    \begin{equation}
    \text{Tr}[\slashed{p}' \gamma^\mu \gamma^\nu] = 0, \qquad \text{Tr}[\gamma^\mu \slashed{p} \gamma^\nu] = 0.
    \end{equation}
    De los cuatro sumandos sobreviven, pues, solo dos; este teorema de anulación es el que mantiene manejable todo el cálculo. Las trazas pares restantes son:
    \begin{align}
    \text{Tr}[\slashed{p}' \gamma^\mu \slashed{p} \gamma^\nu] &= 4 \left( p'^\mu p^\nu + p^\mu p'^\nu - (p' \cdot p) g^{\mu\nu} \right), \\
    \text{Tr}[M^2 \gamma^\mu \gamma^\nu] &= M^2 \text{Tr}[\gamma^\mu \gamma^\nu] = 4 M^2 g^{\mu\nu}.
    \end{align}
    Por consiguiente:
    \begin{align}
    W_{\text{Dirac}}^{\mu\nu} &= \frac{1}{2} \left[ 4 (p'^\mu p^\nu + p^\mu p'^\nu - (p' \cdot p) g^{\mu\nu}) + 4 M^2 g^{\mu\nu} \right] \nonumber \\
    &= 2 \left( p'^\mu p^\nu + p^\mu p'^\nu \right) - 2 (p' \cdot p - M^2) g^{\mu\nu}.
    \end{align}
    Conviene ahora expresar ambos bloques en términos de las variables del problema. Calculemos analíticamente el escalar $p' \cdot p - M^2$:
    \begin{equation}
    q^2 = (p' - p)^2 = p'^2 - 2 p' \cdot p + p^2 = 2M^2 - 2 p' \cdot p \implies 2 (p' \cdot p - M^2) = -q^2 = Q^2.
    \end{equation}
    Expresándolo en términos de $\tau \equiv Q^2/(4M^2)$:
    \begin{equation}
    2 (p' \cdot p - M^2) = 4 M^2 \tau.
    \end{equation}
    Para la combinación simétrica de momentos $p'^\mu p^\nu + p^\mu p'^\nu$, notemos las identidades:
    \begin{align}
    P^\mu P^\nu &= (p' + p)^\mu (p' + p)^\nu = p'^\mu p'^\nu + p^\mu p^\nu + (p'^\mu p^\nu + p^\mu p'^\nu), \\
    q^\mu q^\nu &= (p' - p)^\mu (p' - p)^\nu = p'^\mu p'^\nu + p^\mu p^\nu - (p'^\mu p^\nu + p^\mu p'^\nu).
    \end{align}
    Restando ambas ecuaciones miembro a miembro:
    \begin{equation}
    P^\mu P^\nu - q^\mu q^\nu = 2 (p'^\mu p^\nu + p^\mu p'^\nu) \implies p'^\mu p^\nu + p^\mu p'^\nu = \frac{1}{2} (P^\mu P^\nu - q^\mu q^\nu).
    \end{equation}
    Es aquí donde rentabilizamos la transversalidad demostrada en la subsección anterior. Puesto que el tensor leptónico es transversal ($q^\mu L_{\mu\nu} = 0$), cualquier término proporcional a $q^\mu$ o $q^\nu$ se anula idénticamente al contraer con $L_{\mu\nu}$, y podemos descartarlo sin arrastrarlo por el resto del cálculo. Denotando con el símbolo $\dot{=}$ la equivalencia física bajo la contracción con $L_{\mu\nu}$:
    \begin{equation}
    2 (p'^\mu p^\nu + p^\mu p'^\nu) = P^\mu P^\nu - q^\mu q^\nu \dot{=} P^\mu P^\nu.
    \end{equation}
    Por tanto, la traza de Dirac adquiere la forma desacoplada exacta:
    \begin{equation}
    \dest{W_{\text{Dirac}}^{\mu\nu} \dot{=} P^\mu P^\nu - 4 M^2 \tau g^{\mu\nu}}
    \label{eq:W_dirac_eval}
    \end{equation}

    \item \textbf{Trazas vectoriales cruzadas ($W_{\text{cruz1}}^\mu$ y $W_{\text{cruz2}}^\nu$):}
    \begin{align}
    W_{\text{cruz1}}^\mu &\equiv \frac{1}{2} \text{Tr} \left[ (\slashed{p}' + M) \gamma^\mu (\slashed{p} + M) \right] \nonumber \\
    &= \frac{1}{2} \text{Tr} \left[ \slashed{p}' \gamma^\mu \slashed{p} + M \slashed{p}' \gamma^\mu + M \gamma^\mu \slashed{p} + M^2 \gamma^\mu \right].
    \end{align}
    Las trazas de 3 matrices ($\text{Tr}[\slashed{p}'\gamma^\mu\slashed{p}]$) y de 1 matriz ($\text{Tr}[M^2\gamma^\mu]$) son de orden impar y se anulan idénticamente. Las trazas pares de dos matrices son:
    \begin{align}
    \text{Tr}[M \slashed{p}' \gamma^\mu] &= M p'_\alpha \text{Tr}[\gamma^\alpha \gamma^\mu] = M p'_\alpha (4 g^{\alpha\mu}) = 4 M p'^\mu, \\
    \text{Tr}[M \gamma^\mu \slashed{p}] &= M p_\beta \text{Tr}[\gamma^\mu \gamma^\beta] = M p_\beta (4 g^{\mu\beta}) = 4 M p^\mu.
    \end{align}
    Sumando y multiplicando por $1/2$:
    \begin{equation}
    W_{\text{cruz1}}^\mu = \frac{1}{2} (4 M p'^\mu + 4 M p^\mu) = 2 M (p' + p)^\mu = 2 M P^\mu.
    \end{equation}
    Por simetría cíclica, la segunda traza cruzada es idéntica:
    \begin{equation}
    W_{\text{cruz2}}^\nu = 2 M P^\nu.
    \end{equation}
    Sustituyendo en el término cruzado de la Ec.~\eqref{eq:W_three_parts}:
    \begin{equation}
    -\frac{G_M F_2}{2M} \left[ P^\nu (2 M P^\mu) + P^\mu (2 M P^\nu) \right] = -\frac{G_M F_2}{2M} \left[ 4 M P^\mu P^\nu \right] = -2 G_M F_2 P^\mu P^\nu.
    \label{eq:W_cruz_eval}
    \end{equation}

    \item \textbf{Traza de convección escalar ($W_{\text{conv}}$):}
    \begin{align}
    W_{\text{conv}} &\equiv \frac{1}{2} \text{Tr} \left[ (\slashed{p}' + M)(\slashed{p} + M) \right] = \frac{1}{2} \text{Tr} \left[ \slashed{p}' \slashed{p} + M \slashed{p}' + M \slashed{p} + M^2 \mathbb{I} \right] \nonumber \\
    &= \frac{1}{2} \left[ 4 (p' \cdot p) + 0 + 0 + 4 M^2 \right] = 2 (p' \cdot p + M^2).
    \end{align}
    Recordando que $p' \cdot p - M^2 = 2 M^2 \tau \iff p' \cdot p = M^2 (1 + 2\tau)$:
    \begin{equation}
    p' \cdot p + M^2 = M^2 (1 + 2\tau) + M^2 = 2 M^2 (1 + \tau).
    \end{equation}
    Por tanto:
    \begin{equation}
    W_{\text{conv}} = 2 \left[ 2 M^2 (1 + \tau) \right] = 4 M^2 (1 + \tau).
    \end{equation}
    Multiplicando por su prefactor en la Ec.~\eqref{eq:W_three_parts}:
    \begin{equation}
    \frac{F_2^2}{4M^2} P^\mu P^\nu \left[ 4 M^2 (1 + \tau) \right] = (1 + \tau) F_2^2 P^\mu P^\nu.
    \label{eq:W_conv_eval}
    \end{equation}
\end{enumerate}

Las tres trazas están calculadas y solo resta ensamblarlas. Sumando algebraicamente las tres contribuciones Ecs.~\eqref{eq:W_dirac_eval}, \eqref{eq:W_cruz_eval} y \eqref{eq:W_conv_eval}:
\begin{align}
W^{\mu\nu} &\dot{=} G_M^2 \left[ P^\mu P^\nu - 4 M^2 \tau g^{\mu\nu} \right] - 2 G_M F_2 P^\mu P^\nu + (1 + \tau) F_2^2 P^\mu P^\nu \nonumber \\
&= \left[ G_M^2 - 2 G_M F_2 + (1 + \tau) F_2^2 \right] P^\mu P^\nu - 4 M^2 \tau G_M^2 g^{\mu\nu}.
\end{align}
Analicemos el coeficiente que multiplica al tensor $P^\mu P^\nu$:
\begin{align}
\mathcal{C}_P &\equiv G_M^2 - 2 G_M F_2 + (1 + \tau) F_2^2 = G_M^2 - 2 G_M F_2 + F_2^2 + \tau F_2^2 \nonumber \\
&= (G_M - F_2)^2 + \tau F_2^2.
\end{align}
Por la definición del factor magnético de Sachs $G_M \equiv F_1 + F_2$, se tiene idénticamente $G_M - F_2 = F_1$. Por lo tanto:
\begin{equation}
\mathcal{C}_P = F_1^2 + \tau F_2^2.
\end{equation}
Sustituyendo $G_M = F_1 + F_2$ en el término métrico, el tensor hadrónico adquiere su forma analítica exacta:
\begin{equation}
\dest{W^{\mu\nu} \dot{=} (F_1^2 + \tau F_2^2) P^\mu P^\nu - 4 M^2 \tau (F_1 + F_2)^2 g^{\mu\nu}}
\label{eq:W_munu_exact}
\end{equation}

Compárese la estructura de este resultado con la del tensor leptónico, Ec.~\eqref{eq:L_munu_exact}: ambos son simétricos y se construyen con un término de momentos y un término métrico. La diferencia decisiva es que los coeficientes del hadrónico ya no son números sino funciones desconocidas, $F_1^2 + \tau F_2^2$ y $(F_1+F_2)^2$. Nótese que estas dos combinaciones aparecen aquí por necesidad algebraica y no por elección: el cálculo no produce $F_1$ y $F_2$ aisladamente, sino exactamente estas dos agrupaciones. Cuáles sean las combinaciones que la naturaleza permite medir, y no cuáles nos gustaría medir, es una cuestión que quedará zanjada al contraer los dos tensores.

\subsection{Contracción Invariante \texorpdfstring{$L_{\mu\nu} W^{\mu\nu}$}{L\_munu W\^munu}}
Disponemos ya de las dos mitades del problema en forma cerrada. Contraerlas entre sí es la última operación del cálculo y produce un escalar de Lorentz---un número, sin índices---que es lo único que puede aparecer en una sección eficaz. El trabajo es mecánico: hay exactamente dos contracciones escalares que evaluar, y ambas se calculan en el sistema del laboratorio usando la cinemática de la Sección~\ref{sec:cinematica_trazas}.1.

Contraemos ahora el tensor leptónico $L_{\mu\nu}$ con el tensor hadrónico $W^{\mu\nu}$:
\begin{equation}
L_{\mu\nu} W^{\mu\nu} = (F_1^2 + \tau F_2^2) (L_{\mu\nu} P^\mu P^\nu) - 4 M^2 \tau (F_1 + F_2)^2 (L_{\mu\nu} g^{\mu\nu}).
\label{eq:LW_master}
\end{equation}

Evaluemos minuciosamente cada una de las dos contracciones escalares en el sistema del laboratorio:

\textbf{1. Contracción escalar métrica ($L_{\mu\nu} g^{\mu\nu}$):}
\begin{align}
L_{\mu\nu} g^{\mu\nu} &= 2 \left( k_\mu k'_\nu + k'_\mu k_\nu - (k \cdot k') g_{\mu\nu} \right) g^{\mu\nu} \nonumber \\
&= 2 \left( k_\mu k'^\mu + k'_\mu k^\mu - (k \cdot k') g_\mu^\mu \right).
\end{align}
Dado que $g_\mu^\mu = \delta_\mu^\mu = 4$ en cuatro dimensiones espacio-temporales:
\begin{equation}
L_{\mu\nu} g^{\mu\nu} = 2 \left( k \cdot k' + k' \cdot k - 4 (k \cdot k') \right) = 2 \left( 2 k \cdot k' - 4 k \cdot k' \right) = -4 (k \cdot k').
\end{equation}
De la Ec.~\eqref{eq:k_dot_kprime}, $k \cdot k' = Q^2/2 = 2 M^2 \tau$. Por tanto:
\begin{equation}
\dest{L_{\mu\nu} g^{\mu\nu} = -4 \left(\frac{Q^2}{2}\right) = -2 Q^2 = -8 M^2 \tau}
\label{eq:L_dot_g}
\end{equation}

\textbf{2. Contracción escalar de convección ($L_{\mu\nu} P^\mu P^\nu$):}
\begin{align}
L_{\mu\nu} P^\mu P^\nu &= 2 \left[ (k_\mu P^\mu)(k'_\nu P^\nu) + (k'_\mu P^\mu)(k_\nu P^\nu) - (k \cdot k')(P_\mu P^\mu) \right] \nonumber \\
&= 2 \left[ (k \cdot P)(k' \cdot P) + (k' \cdot P)(k \cdot P) - (k \cdot k') P^2 \right] \nonumber \\
&= 4 (k \cdot P)(k' \cdot P) - 2 (k \cdot k') P^2.
\label{eq:L_dot_PP_start}
\end{align}
Evaluemos los productos escalares $(k \cdot P)$ y $(k' \cdot P)$ expresando $P^\mu$ en términos de las variables del laboratorio:
\begin{equation}
P^\mu = p^\mu + p'^\mu = 2 p^\mu + q^\mu = 2 p^\mu + k^\mu - k'^\mu.
\end{equation}
Contrayendo con el cuadrimomento incidente $k$:
\begin{align}
k \cdot P &= 2 (k \cdot p) + k^2 - k \cdot k' = 2 (k_0 p^0 - \mathbf{k} \cdot \mathbf{0}) + m_e^2 - k \cdot k' \nonumber \\
&\approx 2 M E - k \cdot k'.
\end{align}
Contrayendo con el cuadrimomento dispersado $k'$:
\begin{align}
k' \cdot P &= 2 (k' \cdot p) + k' \cdot k - k'^2 = 2 (k'_0 p^0 - \mathbf{k}' \cdot \mathbf{0}) + k \cdot k' - m_e^2 \nonumber \\
&\approx 2 M E' + k \cdot k'.
\end{align}
Recordando de la cinemática de Compton que $k \cdot k' = M (E - E')$:
\begin{align}
k \cdot P &= 2 M E - M (E - E') = 2 M E - M E + M E' = M (E + E'), \\
k' \cdot P &= 2 M E' + M (E - E') = 2 M E' + M E - M E' = M (E + E').
\end{align}
¡Ambos productos escalares son rigurosamente idénticos! Su producto es:
\begin{equation}
(k \cdot P)(k' \cdot P) = [M (E + E')]^2 = M^2 (E + E')^2.
\end{equation}

Para el segundo término, calculamos la norma cuadrada $P^2$:
\begin{align}
P^2 &= (p' + p)^2 = p'^2 + 2 p' \cdot p + p^2 = M^2 + 2 [M^2(1 + 2\tau)] + M^2 \nonumber \\
&= 2 M^2 + 2 M^2 + 4 M^2 \tau = 4 M^2 (1 + \tau) = 4 M^2 + Q^2.
\end{align}
Sustituyendo estos resultados en la Ec.~\eqref{eq:L_dot_PP_start}:
\begin{align}
L_{\mu\nu} P^\mu P^\nu &= 4 M^2 (E + E')^2 - 2 \left( \frac{Q^2}{2} \right) \left[ 4 M^2 + Q^2 \right] \nonumber \\
&= 4 M^2 (E + E')^2 - 4 M^2 Q^2 - Q^4.
\label{eq:L_PP_step2}
\end{align}
Expandamos ahora algebraicamente el binomio $(E + E')^2$ en términos de la suma y diferencia:
\begin{equation}
(E + E')^2 = (E - E')^2 + 4 E E'.
\end{equation}
Multiplicando por $4 M^2$:
\begin{equation}
4 M^2 (E + E')^2 = 4 M^2 (E - E')^2 + 16 M^2 E E'.
\end{equation}
De la cinemática de Compton, $E - E' = \frac{Q^2}{2M}$, por lo que:
\begin{equation}
4 M^2 (E - E')^2 = 4 M^2 \left( \frac{Q^2}{2M} \right)^2 = 4 M^2 \left( \frac{Q^4}{4 M^2} \right) = Q^4.
\end{equation}
Por consiguiente:
\begin{equation}
4 M^2 (E + E')^2 = Q^4 + 16 M^2 E E'.
\end{equation}
Sustituyendo esta identidad en la Ec.~\eqref{eq:L_PP_step2}:
\begin{align}
L_{\mu\nu} P^\mu P^\nu &= \left[ Q^4 + 16 M^2 E E' \right] - 4 M^2 Q^2 - Q^4 \nonumber \\
&= 16 M^2 E E' - 4 M^2 Q^2.
\end{align}
¡Los términos de orden superior $Q^4$ se cancelan con absoluta exactitud matemática!\\
Sustituyendo $Q^2 = 4 E E' \sin^2(\theta/2)$:
\begin{align}
L_{\mu\nu} P^\mu P^\nu &= 16 M^2 E E' - 4 M^2 \left[ 4 E E' \sin^2\left(\frac{\theta}{2}\right) \right] \nonumber \\
&= 16 M^2 E E' \left[ 1 - \sin^2\left(\frac{\theta}{2}\right) \right].
\end{align}
Utilizando la identidad trigonométrica fundamental $1 - \sin^2(\theta/2) = \cos^2(\theta/2)$:
\begin{equation}
\dest{L_{\mu\nu} P^\mu P^\nu = 16 M^2 E E' \cos^2\left(\frac{\theta}{2}\right)}
\label{eq:L_dot_PP_final}
\end{equation}
La aparición de $\cos^2(\theta/2)$ en esta contracción no es un detalle trigonométrico: es el origen del factor angular característico de la sección eficaz de Mott, y explica por qué la respuesta eléctrica se apaga a ángulos traseros.

Sustituyendo las contracciones exactas Ecs.~\eqref{eq:L_dot_g} y \eqref{eq:L_dot_PP_final} en la ecuación maestra de contracción Ec.~\eqref{eq:LW_master}:
\begin{align}
L_{\mu\nu} W^{\mu\nu} &= (F_1^2 + \tau F_2^2) \left[ 16 M^2 E E' \cos^2\left(\frac{\theta}{2}\right) \right] - 4 M^2 \tau (F_1 + F_2)^2 \left[ -8 M^2 \tau \right] \nonumber \\
&= 16 M^2 E E' \cos^2\left(\frac{\theta}{2}\right) [F_1^2 + \tau F_2^2] + 32 M^4 \tau^2 [F_1 + F_2]^2.
\end{align}
Reescribamos el segundo sumando utilizando la cinemática en el laboratorio:
\begin{equation}
8 M^2 \tau = 2 Q^2 = 8 E E' \sin^2\left(\frac{\theta}{2}\right),
\end{equation}
por lo que:
\begin{equation}
32 M^4 \tau^2 = 4 M^2 \tau (8 M^2 \tau) = 4 M^2 \tau \left[ 8 E E' \sin^2\left(\frac{\theta}{2}\right) \right] = 32 M^2 E E' \tau \sin^2\left(\frac{\theta}{2}\right).
\end{equation}
Por lo tanto:
\begin{equation}
L_{\mu\nu} W^{\mu\nu} = 16 M^2 E E' \cos^2\left(\frac{\theta}{2}\right) [F_1^2 + \tau F_2^2] + 16 M^2 E E' \left[ 2 \tau \sin^2\left(\frac{\theta}{2}\right) \right] [F_1 + F_2]^2.
\end{equation}
Factorizando el término común global $16 M^2 E E'$:
\begin{align}
\dest{
\begin{aligned}
L_{\mu\nu} W^{\mu\nu} = 16 M^2 E E' &\left\{ \cos^2\left(\frac{\theta}{2}\right) \left[ F_1^2(Q^2) + \tau F_2^2(Q^2) \right] \right. \\
&+ \left. 2\tau \sin^2\left(\frac{\theta}{2}\right) \left[ F_1(Q^2) + F_2(Q^2) \right]^2 \right\}
\end{aligned}
}
\label{eq:contraccion_invariante_f1_f2}
\end{align}

Hemos llegado al final del cálculo anunciado, y el resultado es correcto y completo: la Ec.~\eqref{eq:contraccion_invariante_f1_f2} contiene toda la dinámica cuántica del proceso y solo falta dividir por el flujo e integrar el espacio de fases para obtener una sección eficaz. La estructura angular es además la esperada: el primer término, gobernado por $\cos^2(\theta/2)$, domina a ángulos frontales, mientras que el segundo, con $\sin^2(\theta/2)$ y multiplicado por $\tau$, domina a ángulos traseros y en el régimen relativista. Sin embargo, al expandir el binomio al cuadrado $[F_1 + F_2]^2 = F_1^2 + 2 F_1 F_2 + F_2^2$, surge un término de interferencia cruzado cuántico:
\begin{equation}
\mathcal{I}_{\text{cruz}} = + 4 \tau \sin^2\left(\frac{\theta}{2}\right) F_1(Q^2) F_2(Q^2),
\end{equation}
el cual mezcla de forma inseparable los efectos de la carga con los de la magnetización, imposibilitando un desacoplamiento experimental directo en el marco del laboratorio.


\begin{figure}[t]
\centering
\includegraphics[width=\columnwidth]{figuras/cross_sections_comparison.pdf}
\caption{Comparación analítica y cuantitativa de las secciones eficaces diferenciales $e^- p \to e^- p$ en función del ángulo $\theta$ a $E = 1.2\text{ GeV}$: Rutherford clásica (sin retroceso ni espín), Mott relativista con retroceso, Dirac puro ($F_1=1, F_2=0$) y Rosenbluth experimental. A grandes ángulos ($\theta \approx 120^\circ$), la estructura espacial interna del nucleón amortigua la sección eficaz en más de tres órdenes de magnitud respecto a la partícula puntual de Dirac.}
\label{fig:cross_sections}
\end{figure}


La contracción covariante completa $L_{\mu\nu} W^{\mu\nu}$ ha revelado toda la dinámica cuántica de la dispersión, pero nos enfrenta a un severo obstáculo algebraico y físico: la presencia del término de interferencia cruzado cuántico $\mathcal{I}_{\text{cruz}} = 4\tau \sin^2(\theta/2) F_1 F_2$. En esta forma tensorial en el laboratorio, los efectos electrostáticos ($F_1$) y los dipolares magnéticos ($F_2$) aparecen irremediablemente entremezclados, debido a que el protón absorbe energía y sufre retroceso relativista ($q^0 \neq 0$). Conviene ser preciso sobre la naturaleza del problema, porque no es de cálculo sino de principio: dado un conjunto de medidas de la sección eficaz a distintos ángulos, el ajuste no puede distinguir una contribución eléctrica grande con una magnética pequeña de la combinación opuesta, porque solo accede a la suma en la que ambas aparecen multiplicadas entre sí. El origen físico de esa mezcla es exactamente el rasgo que señalamos en la Figura~\ref{fig:cinematica_lab}: en el laboratorio el protón absorbe energía, de modo que $q^0 \neq 0$ y el fotón virtual induce simultáneamente una respuesta de carga y una de corriente. Para transformar esta expresión en la fórmula limpia de Rosenbluth presentada en la Ec.~\eqref{eq:rosenbluth_master_intro} y eliminar algebraicamente las interferencias, resulta imperativo trasladar el análisis a un sistema de referencia privilegiado donde la transferencia de energía se anule estrictamente.
\section{El Marco de Breit y el Desacoplamiento de Sachs}
\label{sec:marco_breit}

La sección anterior terminó con un diagnóstico preciso: el cálculo es correcto pero las dos funciones desconocidas aparecen multiplicadas entre sí, y por tanto no son separables a partir de una medida de la sección eficaz. El remedio que emplearemos puede parecer desproporcionadamente barato en relación con el problema: no vamos a cambiar la teoría, ni a proponer un experimento distinto, ni a introducir ninguna aproximación. Vamos simplemente a describir la misma colisión desde otro sistema de referencia.

Conviene entender por qué eso puede servir de algo, porque a primera vista no debería. La cantidad $L_{\mu\nu}W^{\mu\nu}$ es un escalar de Lorentz: su valor numérico es idéntico en todos los marcos de referencia, y cambiar de marco no puede alterarlo. Lo que sí cambia con el marco es otra cosa, y es justamente lo que necesitamos: \emph{qué combinaciones de $F_1$ y $F_2$ aparecen agrupadas de manera natural}. En el laboratorio, donde el protón absorbe energía, la respuesta de carga y la de corriente se solapan y el álgebra produce el producto $F_1F_2$. Existe, sin embargo, un marco de referencia en el que el protón no absorbe energía en absoluto, y en él las dos respuestas se separan limpiamente en componentes distintas de la corriente. Las combinaciones que ese marco privilegiado señala como naturales son las que adoptaremos como variables, y---este es el punto crucial---una vez adoptadas, la ventaja se conserva al volver al laboratorio, porque un cambio de variables es válido en cualquier marco. Elegir bien el sistema de referencia es, en definitiva, una forma de elegir bien las variables.

\subsection{El Sistema de Breit («Brick-Wall Frame») y Espinores de Pauli}
La raíz del problema quedó identificada al final de la sección anterior: $q^0 \neq 0$. Preguntémonos entonces si existe un marco en el que la transferencia de energía se anule exactamente. La respuesta es afirmativa y la construcción es elemental: basta con situarse en el sistema donde el protón llega con un momento y sale con el momento exactamente opuesto, de modo que su energía cinética no cambie. En ese marco el protón rebota especularmente sin ganar ni perder energía, igual que una pelota contra una pared infinitamente pesada---de ahí el nombre pintoresco de «sistema de pared de ladrillos» con el que se conoce en la literatura anglosajona. La Figura~\ref{fig:marco_breit} contrasta ambas descripciones de la misma colisión.

\begin{figure}[t]
\centering
\begin{tikzpicture}[>=stealth, line width=0.7pt, font=\small]

%======================================================================
% PANEL (a): SISTEMA DEL LABORATORIO
%======================================================================
\begin{scope}[shift={(0,0)}]
  \node[font=\small\bfseries] at (0,2.30) {(a) Sistema del laboratorio};

  % eje del haz
  \draw[gray!45, dashed, line width=0.45pt] (-2.45,0) -- (2.35,0);
  % electrón incidente
  \draw[->] (-2.30,0) -- (-1.25,0);
  \draw     (-1.25,0) -- (-0.34,0);
  \node[anchor=south, font=\scriptsize, inner sep=2.5pt] at (-1.55,0.02) {$e^-$};
  % protón blanco en reposo
  \shade[inner color=orange!58, outer color=orange!8] (0,0) circle (0.30);
  \draw[orange!65!black, line width=0.45pt] (0,0) circle (0.30);
  \node[anchor=north east, font=\scriptsize, inner sep=2pt] at (-0.26,-0.30) {en reposo};
  % electrón dispersado
  \draw[->] (40:0.40) -- (40:1.45);
  \draw     (40:1.45) -- (40:2.15);
  \node[anchor=west, font=\scriptsize, inner sep=2pt] at (40:2.20) {$e^-$};
  % protón de retroceso
  \draw[->, line width=1.05pt] (-30:0.40) -- (-30:1.25);
  \draw[line width=1.05pt]     (-30:1.25) -- (-30:1.85);
  \node[anchor=west, font=\scriptsize, inner sep=2pt] at (-30:1.90) {$p'$};
  % ángulo
  \draw[gray!70, line width=0.45pt] (0.85,0) arc (0:40:0.85);
  \node[gray!85, font=\scriptsize] at (20:1.10) {$\theta$};

  % recuadro diagnóstico
  \node[anchor=north, font=\footnotesize, align=left, inner sep=5pt,
        fill=colordest, draw=colorborde, line width=0.4pt, text width=5.30cm] at (0,-1.48)
    {$q^0 = E - E' \neq 0$\\[0.4ex]
     {\scriptsize El blanco absorbe energía: la carga y la corriente responden a la vez y aparece $F_1 F_2$.}};
\end{scope}

%======================================================================
% PANEL (b): SISTEMA DE BREIT
%======================================================================
\begin{scope}[shift={(9.1,0)}]
  \node[font=\small\bfseries] at (0,2.30) {(b) Sistema de Breit («pared de ladrillos»)};

  % --- la "pared" ---
  \draw[gray!70, line width=0.6pt] (-1.62,-1.15) rectangle (-1.30,1.15);
  \foreach \yy in {-1.10,-0.90,...,1.10}
     \draw[gray!60, line width=0.4pt] (-1.62,\yy) -- (-1.30,\yy+0.18);
  \node[gray!70, font=\scriptsize, anchor=east, align=right, inner sep=4pt] at (-1.68,0)
       {pared\\ infinitamente\\ pesada};

  % --- protón entrante: p = -q/2 (viaja hacia la izquierda) ---
  \draw[line width=1.05pt]     ( 2.45, 0.42) -- ( 0.85, 0.42);
  \draw[->, line width=1.05pt] ( 0.85, 0.42) -- (-1.18, 0.42);
  \node[anchor=west, font=\scriptsize, inner sep=3pt] at (2.50, 0.42)
       {$\mathbf{p} = -\dfrac{\mathbf{q}}{2}$,\; $E_B$};

  % --- protón saliente: p' = +q/2 (viaja hacia la derecha) ---
  \draw[line width=1.05pt]     (-1.18,-0.42) -- ( 0.85,-0.42);
  \draw[->, line width=1.05pt] ( 0.85,-0.42) -- ( 2.45,-0.42);
  \node[anchor=west, font=\scriptsize, inner sep=3pt] at (2.50,-0.42)
       {$\mathbf{p}' = +\dfrac{\mathbf{q}}{2}$,\; $E_B$};

  % --- fotón virtual incidente sobre el punto de rebote ---
  \draw[decorate, decoration={coil, aspect=0, segment length=4.6pt, amplitude=2.6pt},
        line width=0.6pt, gray!85] (-1.04, 2.00) -- (-1.04, 0.66);
  \draw[->, gray!85, line width=0.6pt] (-1.04,0.66) -- (-1.04,0.10);
  \node[gray!85, font=\scriptsize, anchor=west, inner sep=3pt] at (-0.94,1.62)
       {$\gamma^*$,\; $q^\mu = (0,\mathbf{q})$};

  % recuadro diagnóstico
  \node[anchor=north, font=\footnotesize, align=left, inner sep=5pt,
        fill=colordest, draw=colorborde, line width=0.4pt, text width=5.30cm] at (0.55,-1.48)
    {$q^0 = 0$\\[0.4ex]
     {\scriptsize Sin transferencia de energía: la carga responde en $J^0$ y la magnetización en $\mathbf{J}$, por separado ($G_E$ y $G_M$).}};
\end{scope}

% separador
\draw[gray!35, line width=0.4pt] (4.35,-2.65) -- (4.35,2.50);

\end{tikzpicture}
\caption{La misma colisión elástica descrita en dos sistemas de referencia, y la razón por la cual uno de ellos permite separar la carga de la magnetización y el otro no. \textbf{(a)} En el sistema del laboratorio (versión esquemática de la Figura~\ref{fig:cinematica_lab}) el protón parte del reposo y retrocede, de modo que el fotón virtual le transfiere energía además de momento: $q^0 = E - E' \neq 0$. Esa transferencia de energía hace que el fotón excite simultáneamente la distribución de carga y las corrientes de magnetización, y el resultado algebraico es el término de interferencia $F_1F_2$ obtenido en la Ec.~\eqref{eq:contraccion_invariante_f1_f2}. \textbf{(b)} En el sistema de Breit se elige el marco de referencia de modo que el protón llegue con trimomento $\mathbf{p} = -\mathbf{q}/2$ y salga con $\mathbf{p}' = +\mathbf{q}/2$: su momento se invierte pero su módulo, y por tanto su energía $E_B = M\sqrt{1+\tau}$, permanecen inalterados. El protón rebota especularmente, como lo haría contra una pared infinitamente pesada, y el fotón virtual transfiere momento puro, $q^\mu = (0,\mathbf{q})$. Al anularse la componente temporal de la transferencia, la componente temporal de la corriente hadrónica queda acoplada exclusivamente a la densidad estática de carga y las componentes espaciales exclusivamente a la magnetización, sin mezcla. Las dos combinaciones de factores de forma que emergen de ese desacoplamiento son los factores de Sachs $G_E$ y $G_M$, y el resto de la sección consiste en calcularlas explícitamente.}
\label{fig:marco_breit}
\end{figure}

Para desacoplar físicamente las densidades espaciales de carga eléctrica y de corriente de magnetización, R. G. Sachs, F. J. Ernst y K. C. Wali (1960--1962) formularon la interacción en el \textbf{Sistema de Breit} o «sistema de pared de ladrillos» \cite{ernst1960,sachs1962,halzen1984,perdrisat2007}, definido unívocamente por la ausencia total de transferencia neta de energía:
\begin{equation}
q^0_{\text{Breit}} = 0 \iff E_B = E'_B.
\end{equation}
En este marco de referencia cinemático, el fotón virtual transfiere únicamente momento espacial puro:
\begin{equation}
q^\mu_{\text{Breit}} = (0, \mathbf{q}).
\end{equation}
Por conservación del cuadrimomento total $p'^\mu = p^\mu + q^\mu$, el protón simplemente invierte el sentido de su vector de trimomento incidente:
\begin{equation}
\mathbf{p} = -\frac{\mathbf{q}}{2}, \qquad \mathbf{p}' = +\frac{\mathbf{q}}{2}.
\end{equation}
El invariante de Lorentz $Q^2 \equiv -q^2$ coincide exactamente con el cuadrado del trimomento transferido:
\begin{equation}
Q^2 = - (0^2 - |\mathbf{q}|^2) = \mathbf{q}^2.
\end{equation}
Esta última igualdad tiene una consecuencia que resultará decisiva en la Sección~\ref{sec:radio_fourier}: en el marco de Breit, y solo en él, $Q^2$ es literalmente el cuadrado de un trimomento ordinario. Eso es lo que permitirá interpretar los factores de forma como transformadas de Fourier de densidades espaciales en el sentido habitual de la mecánica cuántica no relativista.

La energía relativista del protón en el sistema de Breit $E_B$ se calcula directamente a partir de la relación de Einstein:
\begin{equation}
E_B = \sqrt{M^2 + |\mathbf{p}|^2} = \sqrt{M^2 + \frac{|\mathbf{q}|^2}{4}} = \sqrt{M^2 + \frac{Q^2}{4}} = M \sqrt{1 + \frac{Q^2}{4M^2}} = M \sqrt{1 + \tau}.
\label{eq:EB_energy}
\end{equation}

Para clarificar el contraste operativo y la necesidad conceptual de este cambio de sistema, la Tabla~\ref{tab:marcos_cinematicos} compara las propiedades cinemáticas y dinámicas fundamentales entre el sistema del laboratorio y el sistema de Breit.

\begin{table}[t]
\centering
\small
\renewcommand{\arraystretch}{1.3}
\caption{Comparativa cinemática y dinámica entre el Sistema del Laboratorio y el Sistema de Breit («Brick-Wall Frame»)~\cite{ernst1960,sachs1962,perdrisat2007}. El marco del laboratorio corresponde a la configuración física experimental real del detector, mientras que el marco de Breit representa el sistema teórico privilegiado donde la ausencia de transferencia de energía ($q^0 = 0$) permite aislar limpiamente los momentos multipolares eléctricos y magnéticos sin interferencias cruzadas.}
\label{tab:marcos_cinematicos}
\resizebox{0.98\textwidth}{!}{%
\begin{tabular}{p{4.0cm} p{6.0cm} p{6.0cm}}
\toprule
\textbf{Propiedad / Magnitud Cinemática} & \textbf{Sistema del Laboratorio (Lab)} & \textbf{Sistema de Breit («Pared de Ladrillos»)} \\
\midrule
\textbf{Cuadrimomento inicial del blanco ($p^\mu$)} & $(M, \mathbf{0})$ (protón en reposo absoluto). & $(E_B, -\mathbf{q}/2)$ con $E_B = M\sqrt{1+\tau}$. \\
\addlinespace[0.6ex]
\textbf{Cuadrimomento final del blanco ($p'^\mu$)} & $(E_p', \mathbf{p}')$ con retroceso relativista no nulo. & $(E_B, +\mathbf{q}/2)$ (rebote especular elástico simétrico). \\
\addlinespace[0.6ex]
\textbf{Transferencia de cuadrimomento ($q^\mu$)} & $(E - E', \mathbf{k} - \mathbf{k}') \implies q^0 \neq 0$ (transfiere energía y momento). & $(0, \mathbf{q}) \implies q^0 \equiv 0$ (transferencia puramente espacial). \\
\addlinespace[0.6ex]
\textbf{Transferencia invariante ($Q^2$)} & $Q^2 = 4 E E' \sin^2(\theta/2)$ (acoplada a variables del detector). & $Q^2 = |\mathbf{q}|^2 = \mathbf{q}^2$ (trimomento estático puro). \\
\addlinespace[0.6ex]
\textbf{Estructura de la corriente hadrónica $J^\mu$} & Componentes $J^0$ y $\mathbf{J}$ entremezcladas por el retroceso $q^0 \neq 0$. & \textbf{Desacoplamiento multipolar exacto:}\newline $J^0 = 2M G_E(Q^2)\chi'^\dagger\chi$ (carga estática pura).\newline $\mathbf{J} = i(\boldsymbol{\sigma}\times\mathbf{q})G_M(Q^2)\chi'^\dagger\chi$ (magnetización pura). \\
\addlinespace[0.6ex]
\textbf{Interferencias en $L_{\mu\nu}W^{\mu\nu}$} & Aparece término cruzado $\mathcal{I}_{\text{cruz}} \propto +4\tau\sin^2(\theta/2)F_1 F_2$. & Cancelación idéntica $\pm 2\tau G_E G_M \equiv 0$; suma lineal de $G_E^2$ y $G_M^2$. \\
\addlinespace[0.6ex]
\textbf{Rol operacional y utilidad} & Marco indispensable para la adquisición experimental de datos ($d\sigma/d\Omega$). & Marco analítico óptimo para definir factores de forma y calcular transformadas de Fourier $\rho(r)$. \\
\bottomrule
\end{tabular}%
}%
\end{table}

El cálculo que sigue cambia de técnica respecto a la Sección~\ref{sec:cinematica_trazas}, y conviene anunciarlo. Allí evaluábamos trazas porque promediábamos sobre espines y no nos interesaba ninguna componente en particular. Aquí haremos lo contrario: evaluaremos explícitamente las componentes temporal y espaciales de la corriente $J^\mu = \bar{u}(p')\Gamma^\mu u(p)$, usando la forma concreta de los espinores de Dirac en el marco de Breit. La razón para descender al nivel de los espinores de dos componentes de Pauli es que en este marco tienen interpretación física inmediata: el producto $\chi'^\dagger\chi$ es el solapamiento entre el estado de espín inicial y el final, de modo que un término proporcional a $\chi'^\dagger\chi$ describe una interacción que \emph{no} voltea el espín---típica de la carga electrostática---mientras que un término proporcional a $\boldsymbol{\sigma}\times\mathbf{q}$ describe una que sí lo hace, como corresponde a un acoplamiento magnético. Lo que vamos a comprobar es que $J^0$ contiene exclusivamente lo primero y $\mathbf{J}$ exclusivamente lo segundo.

Para calcular rigurosamente los elementos matriciales de la corriente hadrónica $J^\mu = \bar{u}(p')\Gamma^\mu u(p)$, empleamos la solución exacta de espinores de cuatro componentes de Dirac para fermiones libres de energía positiva en la representación estándar de Dirac-Pauli:
\begin{equation}
u(p) = \sqrt{E_B + M} \begin{pmatrix} \chi \\ \frac{\boldsymbol{\sigma} \cdot \mathbf{p}}{E_B + M} \chi \end{pmatrix}, \qquad
u(p') = \sqrt{E_B + M} \begin{pmatrix} \chi' \\ \frac{\boldsymbol{\sigma} \cdot \mathbf{p}'}{E_B + M} \chi' \end{pmatrix},
\end{equation}
donde $\chi$ y $\chi'$ son espinores de Pauli de dos componentes normalizados a la unidad ($\chi^\dagger \chi = 1$).

Antes de usar estos espinores conviene verificar que están correctamente normalizados, pues de ello dependen todos los factores numéricos posteriores. Demostremos primero la \textbf{normalización covariante relativista} del espinor $\bar{u}(p) u(p) = 2M$:
\begin{align}
\bar{u}(p) u(p) &= u^\dagger(p) \gamma^0 u(p) = (E_B + M) \begin{pmatrix} \chi^\dagger, & \chi^\dagger \frac{\boldsymbol{\sigma} \cdot \mathbf{p}}{E_B + M} \end{pmatrix} \begin{pmatrix} \mathbb{I} & 0 \\ 0 & -\mathbb{I} \end{pmatrix} \begin{pmatrix} \chi \\ \frac{\boldsymbol{\sigma} \cdot \mathbf{p}}{E_B + M} \chi \end{pmatrix} \nonumber \\
&= (E_B + M) \begin{pmatrix} \chi^\dagger, & \chi^\dagger \frac{\boldsymbol{\sigma} \cdot \mathbf{p}}{E_B + M} \end{pmatrix} \begin{pmatrix} \chi \\ -\frac{\boldsymbol{\sigma} \cdot \mathbf{p}}{E_B + M} \chi \end{pmatrix} \nonumber \\
&= (E_B + M) \left[ \chi^\dagger \chi - \chi^\dagger \frac{(\boldsymbol{\sigma} \cdot \mathbf{p})^2}{(E_B + M)^2} \chi \right].
\end{align}
Utilizando $(\boldsymbol{\sigma} \cdot \mathbf{p})^2 = \mathbf{p}^2$:
\begin{equation}
\bar{u}(p) u(p) = (E_B + M) \left[ 1 - \frac{\mathbf{p}^2}{(E_B + M)^2} \right] \chi^\dagger \chi.
\end{equation}
Reduciendo a común denominador:
\begin{align}
1 - \frac{\mathbf{p}^2}{(E_B + M)^2} &= \frac{(E_B + M)^2 - \mathbf{p}^2}{(E_B + M)^2} = \frac{E_B^2 + 2 E_B M + M^2 - (E_B^2 - M^2)}{(E_B + M)^2} \nonumber \\
&= \frac{2 M^2 + 2 E_B M}{(E_B + M)^2} = \frac{2 M (E_B + M)}{(E_B + M)^2} = \frac{2 M}{E_B + M}.
\end{align}
Multiplicando por el prefactor $(E_B + M)$:
\begin{equation}
\dest{\bar{u}(p) u(p) = (E_B + M) \left( \frac{2 M}{E_B + M} \right) = 2 M}
\end{equation}
La combinación $[1 - \mathbf{p}^2/(E_B+M)^2] = 2M/(E_B+M)$ que acaba de aparecer volverá a surgir dos veces más en esta subsección; conviene retenerla, pues es el mecanismo por el cual todos los factores de energía se cancelan y los resultados quedan expresados solo en términos de $M$ y $\tau$.

El espinor adjunto de Dirac para el estado final $p'$ es:
\begin{align}
\bar{u}(p') &= u^\dagger(p') \gamma^0 = \sqrt{E_B + M} \begin{pmatrix} \chi'^\dagger, & \chi'^\dagger \frac{\boldsymbol{\sigma} \cdot \mathbf{p}'}{E_B + M} \end{pmatrix} \begin{pmatrix} \mathbb{I} & 0 \\ 0 & -\mathbb{I} \end{pmatrix} \nonumber \\
&= \sqrt{E_B + M} \begin{pmatrix} \chi'^\dagger, & -\chi'^\dagger \frac{\boldsymbol{\sigma} \cdot \mathbf{p}'}{E_B + M} \end{pmatrix}.
\end{align}

Evaluemos ahora minuciosamente las componentes temporal y espacial de la corriente de transición hadrónica en el sistema de Breit:

\textbf{1. Componente temporal ($J^0_{\text{Breit}}$):}
Esta es la componente que describe la densidad de carga. Si la separación anunciada ocurre realmente, debería resultar proporcional a $\chi'^\dagger\chi$ sin ningún operador de espín intercalado.
\begin{equation}
J^0_{\text{Breit}} = \bar{u}(p') \left[ \gamma^0 F_1(Q^2) + \frac{i \sigma^{0\nu} q_\nu}{2M} F_2(Q^2) \right] u(p).
\end{equation}
Evaluemos primero el elemento de matriz de Dirac $\bar{u}(p') \gamma^0 u(p)$:
\begin{align}
\bar{u}(p') \gamma^0 u(p) &= (E_B + M) \begin{pmatrix} \chi'^\dagger, & -\chi'^\dagger \frac{\boldsymbol{\sigma} \cdot \mathbf{p}'}{E_B + M} \end{pmatrix} \begin{pmatrix} \mathbb{I} & 0 \\ 0 & -\mathbb{I} \end{pmatrix} \begin{pmatrix} \chi \\ \frac{\boldsymbol{\sigma} \cdot \mathbf{p}}{E_B + M} \chi \end{pmatrix} \nonumber \\
&= (E_B + M) \begin{pmatrix} \chi'^\dagger, & -\chi'^\dagger \frac{\boldsymbol{\sigma} \cdot \mathbf{p}'}{E_B + M} \end{pmatrix} \begin{pmatrix} \chi \\ -\frac{\boldsymbol{\sigma} \cdot \mathbf{p}}{E_B + M} \chi \end{pmatrix} \nonumber \\
&= (E_B + M) \left[ \chi'^\dagger \chi + \chi'^\dagger \frac{(\boldsymbol{\sigma} \cdot \mathbf{p}')(\boldsymbol{\sigma} \cdot \mathbf{p})}{(E_B + M)^2} \chi \right].
\end{align}
Es aquí donde la cinemática particular del marco de Breit entra en juego por primera vez. Sustituyendo $\mathbf{p} = -\mathbf{q}/2$ y $\mathbf{p}' = +\mathbf{q}/2$:
\begin{equation}
(\boldsymbol{\sigma} \cdot \mathbf{p}')(\boldsymbol{\sigma} \cdot \mathbf{p}) = \left( \frac{\boldsymbol{\sigma} \cdot \mathbf{q}}{2} \right) \left( -\frac{\boldsymbol{\sigma} \cdot \mathbf{q}}{2} \right) = -\frac{1}{4} (\boldsymbol{\sigma} \cdot \mathbf{q})^2 = -\frac{1}{4} \mathbf{q}^2 = -\frac{Q^2}{4} = - M^2 \tau.
\end{equation}
Obsérvese lo ocurrido: como ambos momentos son paralelos y opuestos, el producto de los dos operadores de espín se ha reducido a un número, sin dejar ninguna matriz de Pauli residual. Esto no sucedería en un marco genérico, y es la primera manifestación concreta de la ventaja del sistema de Breit. Por consiguiente:
\begin{equation}
\bar{u}(p') \gamma^0 u(p) = (E_B + M) \left[ 1 - \frac{M^2 \tau}{(E_B + M)^2} \right] \chi'^\dagger \chi.
\end{equation}
Desarrollando la fracción interior mediante la relación $E_B^2 - M^2 = M^2 \tau$:
\begin{equation}
1 - \frac{M^2 \tau}{(E_B + M)^2} = \frac{(E_B + M)^2 - (E_B^2 - M^2)}{(E_B + M)^2} = \frac{2 M (E_B + M)}{(E_B + M)^2} = \frac{2 M}{E_B + M}.
\end{equation}
Multiplicando por $(E_B + M)$:
\begin{equation}
\dest{\bar{u}(p') \gamma^0 u(p) = 2 M \chi'^\dagger \chi}
\label{eq:breit_gamma0_exact}
\end{equation}

Evaluemos a continuación el término tensorial de Pauli para $\mu = 0$:
\begin{equation}
\sigma^{0\nu} q_\nu = \sigma^{00} q_0 - \sigma^{0j} q_j = 0 - \sigma^{0j} q^j.
\end{equation}
Por definición de las matrices de Dirac:
\begin{align}
\sigma^{0j} &= \frac{i}{2} [\gamma^0, \gamma^j] = \frac{i}{2} \left( \begin{pmatrix} \mathbb{I} & 0 \\ 0 & -\mathbb{I} \end{pmatrix} \begin{pmatrix} 0 & \sigma^j \\ -\sigma^j & 0 \end{pmatrix} - \begin{pmatrix} 0 & \sigma^j \\ -\sigma^j & 0 \end{pmatrix} \begin{pmatrix} \mathbb{I} & 0 \\ 0 & -\mathbb{I} \end{pmatrix} \right) \nonumber \\
&= \frac{i}{2} \left( \begin{pmatrix} 0 & \sigma^j \\ \sigma^j & 0 \end{pmatrix} - \begin{pmatrix} 0 & -\sigma^j \\ -\sigma^j & 0 \end{pmatrix} \right) = \frac{i}{2} \begin{pmatrix} 0 & 2\sigma^j \\ 2\sigma^j & 0 \end{pmatrix} = i \begin{pmatrix} 0 & \sigma^j \\ \sigma^j & 0 \end{pmatrix}.
\end{align}
Multiplicando por $+i$:
\begin{equation}
i \sigma^{0\nu} q_\nu = - i \sigma^{0j} q^j = - i \left[ i \begin{pmatrix} 0 & \sigma^j \\ \sigma^j & 0 \end{pmatrix} q^j \right] = - i^2 \begin{pmatrix} 0 & \boldsymbol{\sigma} \cdot \mathbf{q} \\ \boldsymbol{\sigma} \cdot \mathbf{q} & 0 \end{pmatrix} = \begin{pmatrix} 0 & \boldsymbol{\sigma} \cdot \mathbf{q} \\ \boldsymbol{\sigma} \cdot \mathbf{q} & 0 \end{pmatrix}.
\end{equation}
Nótese que el operador resultante es \emph{antidiagonal} por bloques: conecta la componente superior del espinor con la inferior. Esa estructura es la que hace que el término de Pauli contribuya a $J^0$ con signo opuesto al de Dirac, y de ella saldrá el signo menos que define $G_E$. Aplicamos este operador bloque sobre el espinor $u(p)$:
\begin{align}
(i \sigma^{0\nu} q_\nu) u(p) &= \sqrt{E_B + M} \begin{pmatrix} 0 & \boldsymbol{\sigma} \cdot \mathbf{q} \\ \boldsymbol{\sigma} \cdot \mathbf{q} & 0 \end{pmatrix} \begin{pmatrix} \chi \\ -\frac{\boldsymbol{\sigma} \cdot \mathbf{q}}{2(E_B + M)} \chi \end{pmatrix} \nonumber \\
&= \sqrt{E_B + M} \begin{pmatrix} -\frac{(\boldsymbol{\sigma} \cdot \mathbf{q})^2}{2(E_B + M)} \chi \\ (\boldsymbol{\sigma} \cdot \mathbf{q}) \chi \end{pmatrix} = \sqrt{E_B + M} \begin{pmatrix} -\frac{\mathbf{q}^2}{2(E_B + M)} \chi \\ (\boldsymbol{\sigma} \cdot \mathbf{q}) \chi \end{pmatrix}.
\end{align}
Multiplicando por la izquierda por el espinor adjunto $\bar{u}(p')$:
\begin{align}
\bar{u}(p')(i\sigma^{0\nu} q_\nu)u(p) &= (E_B + M) \begin{pmatrix} \chi'^\dagger, & -\chi'^\dagger \frac{\boldsymbol{\sigma} \cdot \mathbf{q}}{2(E_B + M)} \end{pmatrix} \begin{pmatrix} -\frac{\mathbf{q}^2}{2(E_B + M)} \chi \\ (\boldsymbol{\sigma} \cdot \mathbf{q}) \chi \end{pmatrix} \nonumber \\
&= (E_B + M) \left[ -\frac{\mathbf{q}^2}{2(E_B + M)} \chi'^\dagger \chi - \frac{\chi'^\dagger (\boldsymbol{\sigma} \cdot \mathbf{q})^2 \chi}{2(E_B + M)} \right] \nonumber \\
&= (E_B + M) \left[ -\frac{\mathbf{q}^2}{2(E_B + M)} - \frac{\mathbf{q}^2}{2(E_B + M)} \right] \chi'^\dagger \chi \nonumber \\
&= (E_B + M) \left[ -\frac{\mathbf{q}^2}{E_B + M} \right] \chi'^\dagger \chi = -\mathbf{q}^2 \chi'^\dagger \chi = -Q^2 \chi'^\dagger \chi.
\end{align}
Expresando $Q^2 = 4 M^2 \tau$ y dividiendo entre el factor normalizador de masa $2M$:
\begin{equation}
\dest{\frac{\bar{u}(p')(i\sigma^{0\nu} q_\nu)u(p)}{2M} = -\frac{4 M^2 \tau}{2M} \chi'^\dagger \chi = -2 M \tau \chi'^\dagger \chi}
\label{eq:breit_pauli0_exact}
\end{equation}

Sumando la contribución de Dirac de la Ec.~\eqref{eq:breit_gamma0_exact} y la de Pauli de la Ec.~\eqref{eq:breit_pauli0_exact}:
\begin{align}
J^0_{\text{Breit}} &= 2 M F_1(Q^2) \chi'^\dagger \chi - 2 M \tau F_2(Q^2) \chi'^\dagger \chi \nonumber \\
&= 2 M \left[ F_1(Q^2) - \tau F_2(Q^2) \right] \chi'^\dagger \chi \equiv 2 M G_E(Q^2) \chi'^\dagger \chi.
\label{eq:J0_breit_proved}
\end{align}
¡El factor de forma eléctrico de Sachs $G_E(Q^2) = F_1(Q^2) - \tau F_2(Q^2)$ surge rigurosamente como el acoplamiento directo a la densidad estática de carga en el sistema de Breit! Se ha cumplido la predicción anunciada: la componente temporal de la corriente resultó proporcional a $\chi'^\dagger\chi$, es decir, a una interacción que no voltea el espín, tal como corresponde a la respuesta de una densidad de carga estática. Obsérvese además que el momento magnético anómalo \emph{sí} contribuye a la densidad de carga, y lo hace con signo negativo y con peso $\tau$: un dipolo magnético en movimiento genera, por efecto relativista, una densidad de carga aparente. Ese es el origen físico del término $-\tau F_2$.

\textbf{2. Componente espacial ($\mathbf{J}_{\text{Breit}}$):}
Repetimos ahora el mismo procedimiento para las tres componentes espaciales, que describen la densidad de corriente. La expectativa simétrica a la anterior es que el resultado sea proporcional a $\boldsymbol{\sigma}\times\mathbf{q}$, la estructura característica de un acoplamiento magnético con volteo de espín.
\begin{equation}
\mathbf{J}_{\text{Breit}} = \bar{u}(p') \left[ \boldsymbol{\gamma} F_1(Q^2) + \frac{i \boldsymbol{\sigma}^\nu q_\nu}{2M} F_2(Q^2) \right] u(p).
\end{equation}
Evaluemos el término de Dirac para la componente vectorial espacial $\boldsymbol{\gamma} = \begin{pmatrix} 0 & \boldsymbol{\sigma} \\ -\boldsymbol{\sigma} & 0 \end{pmatrix}$:
\begin{align}
\bar{u}(p') \boldsymbol{\gamma} u(p) &= (E_B + M) \begin{pmatrix} \chi'^\dagger, & -\chi'^\dagger \frac{\boldsymbol{\sigma} \cdot \mathbf{p}'}{E_B + M} \end{pmatrix} \begin{pmatrix} 0 & \boldsymbol{\sigma} \\ -\boldsymbol{\sigma} & 0 \end{pmatrix} \begin{pmatrix} \chi \\ \frac{\boldsymbol{\sigma} \cdot \mathbf{p}}{E_B + M} \chi \end{pmatrix} \nonumber \\
&= (E_B + M) \begin{pmatrix} \chi'^\dagger, & -\chi'^\dagger \frac{\boldsymbol{\sigma} \cdot \mathbf{p}'}{E_B + M} \end{pmatrix} \begin{pmatrix} \boldsymbol{\sigma} \frac{\boldsymbol{\sigma} \cdot \mathbf{p}}{E_B + M} \chi \\ -\boldsymbol{\sigma} \chi \end{pmatrix} \nonumber \\
&= (E_B + M) \left[ \chi'^\dagger \boldsymbol{\sigma} \frac{\boldsymbol{\sigma} \cdot \mathbf{p}}{E_B + M} \chi + \chi'^\dagger \frac{\boldsymbol{\sigma} \cdot \mathbf{p}'}{E_B + M} \boldsymbol{\sigma} \chi \right].
\end{align}
Sustituyendo $\mathbf{p} = -\mathbf{q}/2$ y $\mathbf{p}' = +\mathbf{q}/2$:
\begin{equation}
\bar{u}(p') \boldsymbol{\gamma} u(p) = \frac{1}{2} \chi'^\dagger \left[ -\boldsymbol{\sigma}(\boldsymbol{\sigma} \cdot \mathbf{q}) + (\boldsymbol{\sigma} \cdot \mathbf{q})\boldsymbol{\sigma} \right] \chi = \frac{1}{2} \chi'^\dagger \left[ (\boldsymbol{\sigma} \cdot \mathbf{q})\boldsymbol{\sigma} - \boldsymbol{\sigma}(\boldsymbol{\sigma} \cdot \mathbf{q}) \right] \chi.
\end{equation}
El corchete es exactamente un conmutador, y ahí está la clave: en un marco genérico los dos sumandos tendrían prefactores distintos y no se combinarían así, mientras que la antisimetría $\mathbf{p}' = -\mathbf{p}$ propia del marco de Breit los enfrenta con signos opuestos y produce el conmutador. Evaluemos la componente $k$-ésima del corchete vectorial:
\begin{equation}
\left[ (\boldsymbol{\sigma} \cdot \mathbf{q})\boldsymbol{\sigma} - \boldsymbol{\sigma}(\boldsymbol{\sigma} \cdot \mathbf{q}) \right]_k = q_j (\sigma_j \sigma_k - \sigma_k \sigma_j) = q_j [\sigma_j, \sigma_k].
\end{equation}
Utilizando el conmutador de Pauli $[\sigma_j, \sigma_k] = 2i \epsilon_{jkl} \sigma_l$:
\begin{equation}
q_j [\sigma_j, \sigma_k] = 2i \epsilon_{jkl} q_j \sigma_l = -2i \epsilon_{kjl} \sigma_l q_j = 2i (\boldsymbol{\sigma} \times \mathbf{q})_k.
\end{equation}
Multiplicando por el factor prefactor $1/2$:
\begin{equation}
\dest{\bar{u}(p') \boldsymbol{\gamma} u(p) = i (\boldsymbol{\sigma} \times \mathbf{q}) \chi'^\dagger \chi}
\label{eq:breit_gamma_vec_exact}
\end{equation}

Para el operador dipolar espacial $\frac{i \sigma^{k\nu} q_\nu}{2M}$, puesto que $q_0 = 0$, la contracción con $q_\nu = -q_j$ involucra el conmutador espacial de matrices de Dirac:
\begin{align}
\sigma^{kj} &= \frac{i}{2} [\gamma^k, \gamma^j] = \frac{i}{2} \left[ \begin{pmatrix} 0 & \sigma^k \\ -\sigma^k & 0 \end{pmatrix} \begin{pmatrix} 0 & \sigma^j \\ -\sigma^j & 0 \end{pmatrix} - \begin{pmatrix} 0 & \sigma^j \\ -\sigma^j & 0 \end{pmatrix} \begin{pmatrix} 0 & \sigma^k \\ -\sigma^k & 0 \end{pmatrix} \right] \nonumber \\
&= \frac{i}{2} \left[ \begin{pmatrix} -\sigma^k \sigma^j & 0 \\ 0 & -\sigma^k \sigma^j \end{pmatrix} - \begin{pmatrix} -\sigma^j \sigma^k & 0 \\ 0 & -\sigma^j \sigma^k \end{pmatrix} \right] = -\frac{i}{2} \begin{pmatrix} [\sigma^k, \sigma^j] & 0 \\ 0 & [\sigma^k, \sigma^j] \end{pmatrix} \nonumber \\
&= -\frac{i}{2} (2i \epsilon^{kjl}) \begin{pmatrix} \sigma^l & 0 \\ 0 & \sigma^l \end{pmatrix} = \epsilon^{kjl} \begin{pmatrix} \sigma^l & 0 \\ 0 & \sigma^l \end{pmatrix}.
\end{align}
Contrayendo con $q_\nu = (0, -q^1, -q^2, -q^3)$:
\begin{equation}
\sigma^{k\nu} q_\nu = - \sigma^{kj} q^j = - \epsilon^{kjl} q^j \begin{pmatrix} \sigma^l & 0 \\ 0 & \sigma^l \end{pmatrix} = \epsilon^{klj} q^j \begin{pmatrix} \sigma^l & 0 \\ 0 & \sigma^l \end{pmatrix} = (\boldsymbol{\sigma} \times \mathbf{q})^k \begin{pmatrix} \mathbb{I}_{2\times 2} & 0 \\ 0 & \mathbb{I}_{2\times 2} \end{pmatrix}.
\end{equation}
Multiplicando por la unidad imaginaria $+i$:
\begin{equation}
i \sigma^{k\nu} q_\nu = i (\boldsymbol{\sigma} \times \mathbf{q})^k \begin{pmatrix} \mathbb{I}_{2\times 2} & 0 \\ 0 & \mathbb{I}_{2\times 2} \end{pmatrix}.
\end{equation}
A diferencia del caso temporal, este operador es \emph{diagonal} por bloques, y por eso el término de Pauli contribuirá a la corriente espacial con el mismo signo que el de Dirac. Esta diferencia de estructura entre el caso temporal y el espacial es, en última instancia, la razón de que $G_E$ lleve un signo menos y $G_M$ un signo más. Insertando este operador diagonal por bloques entre los espinores $\bar{u}(p')$ y $u(p)$:
\begin{align}
\bar{u}(p') (i \sigma^{k\nu} q_\nu) u(p) &= (E_B + M) \begin{pmatrix} \chi'^\dagger, & -\chi'^\dagger \frac{\boldsymbol{\sigma} \cdot \mathbf{p}'}{E_B + M} \end{pmatrix} \begin{pmatrix} i(\boldsymbol{\sigma} \times \mathbf{q})^k & 0 \\ 0 & i(\boldsymbol{\sigma} \times \mathbf{q})^k \end{pmatrix} \begin{pmatrix} \chi \\ \frac{\boldsymbol{\sigma} \cdot \mathbf{p}}{E_B + M} \chi \end{pmatrix} \nonumber \\
&= (E_B + M) \left[ \chi'^\dagger i(\boldsymbol{\sigma} \times \mathbf{q})^k \chi - \chi'^\dagger \frac{(\boldsymbol{\sigma} \cdot \mathbf{p}') i(\boldsymbol{\sigma} \times \mathbf{q})^k (\boldsymbol{\sigma} \cdot \mathbf{p})}{(E_B + M)^2} \chi \right].
\end{align}
Sustituyendo $\mathbf{p}' = +\mathbf{q}/2$ y $\mathbf{p} = -\mathbf{q}/2$ en el segundo término:
\begin{equation}
(\boldsymbol{\sigma} \cdot \mathbf{p}') i(\boldsymbol{\sigma} \times \mathbf{q})^k (\boldsymbol{\sigma} \cdot \mathbf{p}) = \left( \frac{\boldsymbol{\sigma} \cdot \mathbf{q}}{2} \right) i(\boldsymbol{\sigma} \times \mathbf{q})^k \left( -\frac{\boldsymbol{\sigma} \cdot \mathbf{q}}{2} \right) = - \frac{i}{4} (\boldsymbol{\sigma} \cdot \mathbf{q}) (\boldsymbol{\sigma} \times \mathbf{q})^k (\boldsymbol{\sigma} \cdot \mathbf{q}).
\end{equation}
Evaluemos minuciosamente el producto de Pauli $(\boldsymbol{\sigma} \cdot \mathbf{q}) (\boldsymbol{\sigma} \times \mathbf{q})^k$:
\begin{align}
(\boldsymbol{\sigma} \cdot \mathbf{q}) (\boldsymbol{\sigma} \times \mathbf{q})^k &= q_j \epsilon^{klm} q^m \sigma^j \sigma^l = q_j q^m \epsilon^{klm} (\delta^{jl} \mathbb{I} + i \epsilon^{jln}\sigma^n) \nonumber \\
&= q_j q^m \epsilon^{kjm} \mathbb{I} + i q_j q^m \epsilon^{klm} \epsilon^{jln} \sigma^n.
\end{align}
El término escalar se anula idénticamente por contracción de los índices simétricos $j, m$ con el tensor antisimétrico de Levi-Civita: $q_j q^m \epsilon^{kjm} = 0$. Utilizando la contracción $\epsilon^{klm}\epsilon^{jln} = - \epsilon^{lmk}\epsilon^{lnj} = -(\delta^{mn}\delta^{kj} - \delta^{mj}\delta^{kn}) = \delta^{mj}\delta^{kn} - \delta^{mn}\delta^{kj}$:
\begin{align}
(\boldsymbol{\sigma} \cdot \mathbf{q}) (\boldsymbol{\sigma} \times \mathbf{q})^k &= i q_j q^m (\delta^{mj}\delta^{kn} - \delta^{mn}\delta^{kj}) \sigma^n \nonumber \\
&= i \left[ q_j q^j \sigma^k - q^k (q_m \sigma^m) \right] = i \left[ \mathbf{q}^2 \sigma^k - q^k (\boldsymbol{\sigma} \cdot \mathbf{q}) \right].
\end{align}
Multiplicando ahora por $(\boldsymbol{\sigma} \cdot \mathbf{q})$ por la derecha:
\begin{align}
(\boldsymbol{\sigma} \cdot \mathbf{q}) (\boldsymbol{\sigma} \times \mathbf{q})^k (\boldsymbol{\sigma} \cdot \mathbf{q}) &= i \left[ \mathbf{q}^2 \sigma^k (\boldsymbol{\sigma} \cdot \mathbf{q}) - q^k (\boldsymbol{\sigma} \cdot \mathbf{q})^2 \right] \nonumber \\
&= i \left[ \mathbf{q}^2 (q^k \mathbb{I} + i (\boldsymbol{\sigma} \times \mathbf{q})^k) - q^k \mathbf{q}^2 \mathbb{I} \right] \nonumber \\
&= i \left[ i \mathbf{q}^2 (\boldsymbol{\sigma} \times \mathbf{q})^k \right] = - \mathbf{q}^2 (\boldsymbol{\sigma} \times \mathbf{q})^k.
\end{align}
El operador $(\boldsymbol{\sigma}\times\mathbf{q})^k$ ha reaparecido intacto, multiplicado por un escalar: es un autovector de esta operación, y por eso toda la componente espacial de la corriente mantiene la misma estructura tensorial. Sustituyendo este resultado fundamental en el numerador del segundo término:
\begin{equation}
- \chi'^\dagger \frac{(\boldsymbol{\sigma} \cdot \mathbf{p}') i(\boldsymbol{\sigma} \times \mathbf{q})^k (\boldsymbol{\sigma} \cdot \mathbf{p})}{(E_B + M)^2} \chi = - \frac{- \frac{i}{4} [ - \mathbf{q}^2 (\boldsymbol{\sigma} \times \mathbf{q})^k ]}{(E_B + M)^2} \chi'^\dagger \chi = - \frac{\frac{1}{4} \mathbf{q}^2}{(E_B + M)^2} i (\boldsymbol{\sigma} \times \mathbf{q})^k \chi'^\dagger \chi.
\end{equation}
Como $\frac{1}{4}\mathbf{q}^2 = M^2 \tau$:
\begin{align}
\bar{u}(p') (i \sigma^{k\nu} q_\nu) u(p) &= (E_B + M) \left[ 1 - \frac{M^2 \tau}{(E_B + M)^2} \right] i (\boldsymbol{\sigma} \times \mathbf{q})^k \chi'^\dagger \chi \nonumber \\
&= (E_B + M) \left[ \frac{2M}{E_B + M} \right] i (\boldsymbol{\sigma} \times \mathbf{q})^k \chi'^\dagger \chi = 2 M i (\boldsymbol{\sigma} \times \mathbf{q})^k \chi'^\dagger \chi.
\end{align}
Dividiendo entre la escala de masa $2M$:
\begin{equation}
\dest{\frac{\bar{u}(p') (i \sigma^{k\nu} q_\nu) u(p)}{2M} = i (\boldsymbol{\sigma} \times \mathbf{q})^k \chi'^\dagger \chi}
\label{eq:breit_pauli_vec_exact}
\end{equation}

Sumando ambas contribuciones espaciales:
\begin{equation}
\mathbf{J}_{\text{Breit}} = i (\boldsymbol{\sigma} \times \mathbf{q}) \left[ F_1(Q^2) + F_2(Q^2) \right] \chi'^\dagger \chi \equiv i (\boldsymbol{\sigma} \times \mathbf{q}) G_M(Q^2) \chi'^\dagger \chi.
\label{eq:Jvec_breit_proved}
\end{equation}
¡El factor de forma magnético de Sachs $G_M(Q^2) = F_1(Q^2) + F_2(Q^2)$ surge rigurosamente como el acoplamiento directo a la densidad de magnetización o corriente dipolar espacial!

Merece la pena detenerse en el balance de esta subsección, porque contiene el argumento central de toda la sección. Las Ecs.~\eqref{eq:J0_breit_proved} y \eqref{eq:Jvec_breit_proved} muestran que, en el sistema de Breit, la corriente hadrónica se parte en dos piezas que no se comunican entre sí: la componente temporal depende exclusivamente de la combinación $F_1 - \tau F_2$ y describe la densidad estática de carga, sin volteo de espín; las componentes espaciales dependen exclusivamente de $F_1 + F_2$ y describen la magnetización, con la estructura $\boldsymbol{\sigma}\times\mathbf{q}$ característica del volteo de espín. Ningún término contiene ambas combinaciones. Conviene subrayar que esas dos agrupaciones no las hemos elegido nosotros por conveniencia de notación: han emergido del cálculo, y son la respuesta a la pregunta de qué combinaciones de $F_1$ y $F_2$ corresponden a observables físicos separables. Bautizarlas es lo que hacemos a continuación.

\subsection{Definición y Demostración de Cancelación de Sachs}
Dar un nombre a las dos combinaciones que acaban de emerger no es, sin embargo, una demostración. Lo que todavía falta por probar es que ese cambio de variables---sugerido por un marco de referencia particular---elimina efectivamente el término de interferencia de la expresión \emph{covariante}, Ec.~\eqref{eq:contraccion_invariante_f1_f2}, que fue calculada en el laboratorio y es válida en cualquier marco. Solo entonces habremos ganado algo experimentalmente útil. Esa comprobación es el objeto de esta subsección y consiste, una vez planteada, en un ejercicio de álgebra elemental cuyo desenlace no es en absoluto evidente de antemano.

Se definen de forma unívoca los \textbf{Factores de Forma de Sachs}:
\begin{align}
\boxed{G_E(Q^2) \equiv F_1(Q^2) - \tau F_2(Q^2)} \quad & (\text{Factor Eléctrico}), \label{eq:GE_def} \\
\boxed{G_M(Q^2) \equiv F_1(Q^2) + F_2(Q^2)} \quad & (\text{Factor Magnético}). \label{eq:GM_def}
\end{align}
Como esta definición es una transformación lineal invertible, podemos expresar $F_1$ y $F_2$ en términos de $G_E$ y $G_M$ y sustituir en cualquier expresión previa. Invirtiendo algebraicamente el sistema lineal $2 \times 2$:
\begin{equation}
G_M - G_E = (F_1 + F_2) - (F_1 - \tau F_2) = (1 + \tau) F_2,
\end{equation}
luego:
\begin{equation}
\dest{F_2(Q^2) = \frac{G_M(Q^2) - G_E(Q^2)}{1 + \tau}}
\label{eq:F2_inverted}
\end{equation}
Sustituyendo $F_2$ en la expresión de $G_M$:
\begin{equation}
F_1 = G_M - F_2 = G_M - \frac{G_M - G_E}{1 + \tau} = \frac{G_M(1+\tau) - G_M + G_E}{1 + \tau},
\end{equation}
de donde:
\begin{equation}
\dest{F_1(Q^2) = \frac{G_E(Q^2) + \tau G_M(Q^2)}{1 + \tau}}
\label{eq:F1_inverted}
\end{equation}

Ahora sustituyamos minuciosamente las Ecs.~\eqref{eq:F1_inverted} y \eqref{eq:F2_inverted} en el primer término de la contracción tensorial de la Ec.~\eqref{eq:contraccion_invariante_f1_f2}:
\begin{equation}
\mathcal{S} \equiv F_1^2 + \tau F_2^2 = \left[ \frac{G_E + \tau G_M}{1 + \tau} \right]^2 + \tau \left[ \frac{G_M - G_E}{1 + \tau} \right]^2.
\end{equation}
Desarrollando los numeradores con el binomio al cuadrado:
\begin{align}
(G_E + \tau G_M)^2 &= G_E^2 + 2\tau G_E G_M + \tau^2 G_M^2, \\
\tau (G_M - G_E)^2 &= \tau (G_M^2 - 2 G_E G_M + G_E^2) \nonumber \\
&= \tau G_M^2 - 2\tau G_E G_M + \tau G_E^2.
\end{align}
El punto decisivo está a la vista: cada binomio ha generado un término cruzado $G_EG_M$, y ambos tienen el mismo coeficiente $2\tau$ pero signos opuestos. Sumando algebraicamente ambos numeradores sobre el denominador común $(1+\tau)^2$:
\begin{align}
\mathcal{S} &= \frac{G_E^2 + 2\tau G_E G_M + \tau^2 G_M^2 + \tau G_E^2 - 2\tau G_E G_M + \tau G_M^2}{(1+\tau)^2} \nonumber \\
&= \frac{G_E^2(1 + \tau) + \tau G_M^2(1 + \tau) + (2\tau - 2\tau) G_E G_M}{(1+\tau)^2}.
\end{align}
¡Los términos de interferencia de espín $+2\tau G_E G_M$ y $-2\tau G_E G_M$ se cancelan con absoluta exactitud matemática! Factorizando $(1+\tau)$:
\begin{equation}
\mathcal{S} = \frac{(1+\tau) \left[ G_E^2 + \tau G_M^2 \right]}{(1+\tau)^2} = \frac{G_E^2 + \tau G_M^2}{1 + \tau}.
\end{equation}
Por tanto, hemos demostrado analíticamente la identidad exacta:
\begin{equation}
\dest{F_1^2(Q^2) + \tau F_2^2(Q^2) = \frac{G_E^2(Q^2) + \tau G_M^2(Q^2)}{1 + \tau}}
\label{eq:sachs_identity_proved}
\end{equation}
Nótese que la cancelación depende del factor $\tau$ que acompaña a $F_2$ en la definición de $G_E$: con cualquier otra combinación lineal de $F_1$ y $F_2$ los dos términos cruzados no tendrían el mismo módulo y sobreviviría un residuo. En este sentido, $G_E$ y $G_M$ no son una elección conveniente entre muchas posibles, sino las únicas variables que producen el desacoplamiento.

Sustituyendo la Ec.~\eqref{eq:sachs_identity_proved} y $[F_1 + F_2]^2 = G_M^2$ en la Ec.~\eqref{eq:contraccion_invariante_f1_f2}:
\begin{align}
L_{\mu\nu} W^{\mu\nu} &= 16 M^2 E E' \left\{ \cos^2\left(\frac{\theta}{2}\right) \left[ \frac{G_E^2 + \tau G_M^2}{1 + \tau} \right] + 2\tau \sin^2\left(\frac{\theta}{2}\right) G_M^2 \right\} \nonumber \\
&= 16 M^2 E E' \cos^2\left(\frac{\theta}{2}\right) \left[ \frac{G_E^2 + \tau G_M^2}{1 + \tau} + 2\tau G_M^2 \frac{\sin^2(\theta/2)}{\cos^2(\theta/2)} \right].
\end{align}
Utilizando la identidad $\tan^2(\theta/2) = \sin^2(\theta/2)/\cos^2(\theta/2)$, obtenemos la contracción en su forma desacoplada definitiva:
\begin{equation}
\dest{
L_{\mu\nu} W^{\mu\nu} = 16 M^2 E E' \cos^2\left(\frac{\theta}{2}\right) \left[ \frac{G_E^2 + \tau G_M^2}{1 + \tau} + 2\tau G_M^2 \tan^2\left(\frac{\theta}{2}\right) \right]
}
\label{eq:LW_sachs_final}
\end{equation}

Compárese esta expresión con la Ec.~\eqref{eq:contraccion_invariante_f1_f2}, que es la misma cantidad escrita en las variables anteriores: allí aparecía el producto $F_1F_2$ y aquí no aparece ningún producto $G_EG_M$. Conviene ser explícito sobre lo que esto significa y lo que no. No hemos hecho ninguna aproximación ni hemos descartado ningún término: las dos expresiones son idénticas, y describen exactamente la misma física. Lo que ha cambiado es únicamente la elección de las dos funciones incógnita, y con ella la posibilidad de determinarlas. En las variables $G_E$ y $G_M$, y en ningún otro par de combinaciones, la sección eficaz es una \emph{función lineal de dos cuadrados independientes}, y un ajuste de datos puede por tanto separarlos. Obsérvese también que la ventaja, una vez ganada, es permanente: la Ec.~\eqref{eq:LW_sachs_final} está escrita en variables del laboratorio y es válida allí. El sistema de Breit no ha sido el marco donde se hace el experimento, sino el marco que nos indicó qué preguntas tienen respuesta separable.


Con la formulación de los factores de forma de Sachs ($G_E, G_M$) y la demostración de la cancelación idéntica de las interferencias cruzadas, el elemento de matriz invariante $\overline{|\mathcal{M}|^2}$ ha adquirido una estructura matemática limpia y desacoplada. Sin embargo, una probabilidad de transición cuántica no constituye todavía un observable físico directo: para predecir la tasa de dispersión angular registrada por los detectores en el laboratorio, es imprescindible dividir por el flujo incidente de electrones e integrar el espacio de fases relativista de dos cuerpos invariante de Lorentz ($d\Phi_2$). Este último paso es de naturaleza más rutinaria que los anteriores---no requiere ninguna idea nueva, solo contabilidad cuidadosa de estados finales---pero es el que convierte una amplitud en un número comparable con un experimento. Completará formalmente la demostración de la sección eficaz de Rosenbluth anunciada en la Ec.~\eqref{eq:rosenbluth_master_intro} y revelará la técnica experimental para desacoplar y extraer cuantitativamente $G_E$ y $G_M$.
\section{Demostración Formal de la Sección Eficaz y Separación Experimental}
\label{sec:seccion_rosenbluth}

Esta sección tiene dos mitades de carácter muy distinto. La primera es contabilidad: convertir la probabilidad de transición $\overline{|\mathcal{M}|^2}$, que ya tenemos calculada y desacoplada, en una tasa de conteo por unidad de ángulo sólido. No requiere ninguna idea física nueva, solo cuidado con los factores de normalización. La segunda mitad es el punto culminante del artículo. En ella comprobaremos que toda la maquinaria covariante que hemos levantado desde la Sección~\ref{sec:vertice_gordon}---vértice general, identidad de Gordon, identidad de Ward, trazas de Dirac, marco de Breit---reproduce, término por término y sin margen de ajuste, la ecuación que Marshall Rosenbluth dedujo a mano en 1950 y que presentamos en la Ec.~\eqref{eq:rosenbluth_1950_original} sin justificar. Cerrado ese círculo, mostraremos cómo la fórmula se convierte en un método de medida.

\subsection{Integración del Espacio de Fases de Dos Cuerpos}
El paso de $\overline{|\mathcal{M}|^2}$ a $d\sigma/d\Omega$ requiere dos ingredientes que no tienen nada de cuántico. El primero es dividir por el flujo de partículas incidentes, porque una sección eficaz es una probabilidad \emph{por partícula incidente y por centro dispersor}. El segundo es contar cuántos estados finales son accesibles en el intervalo de ángulo sólido considerado, lo que se conoce como espacio de fases. La única dificultad técnica de esta subsección reside en que la conservación de energía impone una delta de Dirac cuyo argumento depende de $E'$ de manera no trivial---a través de la energía del protón de retroceso---y evaluarla correctamente exige derivar ese argumento. A eso dedicaremos la mayor parte del cálculo.

La sección eficaz diferencial invariante se obtiene dividiendo la probabilidad de transición por el flujo incidente $\Phi$ e integrando sobre el espacio de fases de Lorentz de dos cuerpos $d\Phi_2$~\cite{halzen1984,griffiths2008}:
\begin{equation}
d\sigma = \frac{\overline{|\mathcal{M}|^2}}{\Phi} d\Phi_2.
\end{equation}
El flujo incidente en el sistema de laboratorio (con el protón blanco en reposo) es:
\begin{equation}
\Phi = 4 \sqrt{(k \cdot p)^2 - m_e^2 M^2} \approx 4 (k \cdot p) = 4 M E.
\end{equation}
El espacio de fases relativista de dos cuerpos es:
\begin{equation}
d\Phi_2 = \frac{d^3\mathbf{k}'}{(2\pi)^3 2E'} \frac{d^3\mathbf{p}'}{(2\pi)^3 2E'_p} (2\pi)^4 \delta^{(4)}(k + p - k' - p').
\end{equation}
La delta cuadridimensional impone las cuatro leyes de conservación. Tres de ellas---las de momento---se usan de inmediato para eliminar la integración sobre el protón final, que no se detecta. Integrando trivialmente sobre el trimomento del protón $d^3\mathbf{p}'$ mediante la delta tridimensional, fijamos $\mathbf{p}' = \mathbf{k} - \mathbf{k}'$:
\begin{equation}
d\Phi_2 = \frac{1}{16\pi^2} \frac{d^3\mathbf{k}'}{E' E'_p} \delta(E + M - E' - E'_p).
\end{equation}
Queda una sola delta, la de la energía, y es la que exige trabajo. En coordenadas esféricas: $d^3\mathbf{k}' = E'^2 dE' d\Omega$. Para evaluar la delta unidimensional de conservación de energía:
\begin{equation}
f(E') \equiv E + M - E' - \sqrt{M^2 + E^2 + E'^2 - 2 E E' \cos\theta}.
\end{equation}
Derivando $f(E')$ respecto a $E'$:
\begin{align}
\frac{df}{dE'} &= -1 - \frac{2E' - 2E\cos\theta}{2\sqrt{M^2 + (\mathbf{k}-\mathbf{k}')^2}} = -1 - \frac{E' - E\cos\theta}{E'_p} \nonumber \\
&= -\frac{E'_p + E' - E\cos\theta}{E'_p}.
\end{align}
Utilizando la conservación de energía $E'_p = E + M - E'$:
\begin{equation}
E'_p + E' - E\cos\theta = E + M - E\cos\theta = M + E(1 - \cos\theta).
\end{equation}
De la cinemática de Compton de la Ec.~\eqref{eq:recoil_energy_exact}:
\begin{equation}
\frac{E}{E'} = 1 + \frac{E}{M}(1 - \cos\theta) \implies M \frac{E}{E'} = M + E(1 - \cos\theta).
\end{equation}
Por tanto:
\begin{equation}
\left| \frac{df}{dE'} \right| = \frac{M E}{E' E'_p}.
\end{equation}
Aplicando la propiedad de la delta de Dirac $\int dE' g(E') \delta(f(E')) = \frac{g(E'_0)}{|df/dE'|}$:
\begin{equation}
d\Phi_2 = \frac{1}{16\pi^2} \frac{E'^2 d\Omega}{E' E'_p} \left( \frac{E' E'_p}{M E} \right) = \frac{1}{16\pi^2 M} \left( \frac{E'}{E} \right) E' d\Omega.
\end{equation}
Conviene señalar la aparición del cociente $E'/E$ en este resultado, porque es el mismo factor que veremos dentro de la sección eficaz de Mott y que suele llamarse «factor de retroceso». Su origen queda ahora claro y no es dinámico: no proviene de la interacción sino del recuento de estados finales, es decir, del hecho de que el protón se lleva una parte del momento y comprime el espacio de fases disponible para el electrón.

\subsection{Deducción de la Fórmula de Rosenbluth}
Todas las piezas están sobre la mesa. Multiplicando el flujo, el espacio de fases y la contracción desacoplada de la Ec.~\eqref{eq:LW_sachs_final}, la sección eficaz queda determinada sin ninguna libertad restante. Multiplicando todos los factores:
\begin{align}
\frac{d\sigma}{d\Omega} &= \frac{1}{4 M E} \frac{1}{16\pi^2 M} \left(\frac{E'}{E}\right) E' \left[ \frac{e^4}{Q^4} L_{\mu\nu} W^{\mu\nu} \right] \nonumber \\
&= \frac{e^4}{64 \pi^2 M^2 E^2} \left(\frac{E'}{E}\right) E' \frac{1}{Q^4} L_{\mu\nu} W^{\mu\nu}.
\end{align}
Reemplazando $e^2 = 4\pi\alpha$ y $Q^2 = 4 E E' \sin^2(\theta/2) \implies Q^4 = 16 E^2 E'^2 \sin^4(\theta/2)$:
\begin{align}
\frac{d\sigma}{d\Omega} &= \frac{16 \pi^2 \alpha^2}{64 \pi^2 M^2 E^2} \left(\frac{E'}{E}\right) E' \frac{1}{16 E^2 E'^2 \sin^4(\theta/2)} L_{\mu\nu} W^{\mu\nu} \nonumber \\
&= \frac{\alpha^2}{64 M^2 E^4 E' \sin^4(\theta/2)} \left(\frac{E'}{E}\right) L_{\mu\nu} W^{\mu\nu}.
\end{align}
Insertando la contracción covariante de Sachs de la Ec.~\eqref{eq:LW_sachs_final}:
\begin{align}
\frac{d\sigma}{d\Omega} &= \frac{\alpha^2}{64 M^2 E^4 E' \sin^4(\theta/2)} \left(\frac{E'}{E}\right) \left[ 16 M^2 E E' \cos^2\left(\frac{\theta}{2}\right) \right] \nonumber \\
&\quad \times \left[ \frac{G_E^2 + \tau G_M^2}{1 + \tau} + 2\tau G_M^2 \tan^2\left(\frac{\theta}{2}\right) \right].
\end{align}
Simplificando los factores constantes $16 M^2 E E' / (64 M^2 E^4 E') = \frac{1}{4 E^2}$:
\begin{align}
\frac{d\sigma}{d\Omega} &= \left[ \frac{\alpha^2 \cos^2(\theta/2)}{4 E^2 \sin^4(\theta/2)} \left(\frac{E'}{E}\right) \right] \nonumber \\
&\quad \times \left[ \frac{G_E^2(Q^2) + \tau G_M^2(Q^2)}{1+\tau} + 2\tau G_M^2(Q^2) \tan^2\left(\frac{\theta}{2}\right) \right].
\end{align}
El resultado se ha factorizado espontáneamente en dos corchetes de significado muy distinto, y esa separación es la que da a la fórmula su valor práctico. El primero depende únicamente de la cinemática y de la carga del electrón: es lo que se obtendría si el protón fuese puntual. El segundo contiene toda la estructura del nucleón. Identificando la sección eficaz de Mott con retroceso cinemático:
\begin{equation}
\left( \frac{d\sigma}{d\Omega} \right)_{\text{Mott}} \equiv \frac{\alpha^2 \cos^2(\theta/2)}{4 E^2 \sin^4(\theta/2)} \left( \frac{E'}{E} \right),
\label{eq:mott_def}
\end{equation}
arribamos a la célebre \textbf{Fórmula de Rosenbluth} (1950) \cite{rosenbluth1950}, la cual escribimos tanto en su forma canónica abreviada como en sus \textbf{dos formas completas y explícitas}:
\begin{equation}
\boxed{
\left( \frac{d\sigma}{d\Omega} \right) = \left( \frac{d\sigma}{d\Omega} \right)_{\text{Mott}} \left[ \frac{G_E^2(Q^2) + \tau G_M^2(Q^2)}{1+\tau} + 2\tau G_M^2(Q^2) \tan^2\left(\frac{\theta}{2}\right) \right]
}
\label{eq:rosenbluth_formula_box}
\end{equation}
Esta ecuación es el objetivo anunciado en la hoja de ruta de la Sección~\ref{sec:punto_de_partida}, y conviene notar que nada en ella ha sido postulado: la estructura del corchete, con sus dos cuadrados independientes y su dependencia angular, es consecuencia forzosa de la cadena que arranca en la invariancia de Lorentz. Al desplegar completamente el factor de Mott con retroceso elástico $(E'/E) = [1 + (2E/M)\sin^2(\theta/2)]^{-1}$, la \textbf{fórmula completa de Sachs} queda desglosada en su totalidad sin factores abreviados:
\begin{equation}
\dest{
\left(\frac{d\sigma}{d\Omega}\right) = \frac{\alpha^2 \cos^2\left(\dfrac{\theta}{2}\right)}{4 E^2 \sin^4\left(\dfrac{\theta}{2}\right) \left[ 1 + \dfrac{2E}{M}\sin^2\left(\dfrac{\theta}{2}\right) \right]} \left[ \frac{G_E^2(Q^2) + \tau G_M^2(Q^2)}{1+\tau} + 2\tau G_M^2(Q^2) \tan^2\left(\frac{\theta}{2}\right) \right]
}
\label{eq:rosenbluth_complete_sachs_sec5}
\end{equation}
Asimismo, expresada en términos de los factores de forma de Dirac $F_1(Q^2)$ y Pauli $F_2(Q^2)$, la \textbf{fórmula completa de Dirac-Pauli} se escribe como:
\begin{equation}
\dest{
\left(\frac{d\sigma}{d\Omega}\right) = \frac{\alpha^2 \cos^2\left(\dfrac{\theta}{2}\right)}{4 E^2 \sin^4\left(\dfrac{\theta}{2}\right) \left[ 1 + \dfrac{2E}{M}\sin^2\left(\dfrac{\theta}{2}\right) \right]} \left\{ F_1^2(Q^2) + \tau \left[ F_2^2(Q^2) + 2\Big( F_1(Q^2) + F_2(Q^2) \Big)^2 \tan^2\left(\frac{\theta}{2}\right) \right] \right\}
}
\label{eq:rosenbluth_complete_dirac_sec5}
\end{equation}
donde $\tau = \frac{E^2 \sin^2(\theta/2)}{M^2 [1 + (2E/M)\sin^2(\theta/2)]}$. Esta última forma, escrita en las variables originales $F_1$ y $F_2$ en lugar de las de Sachs, es la que permite el contraste histórico que emprendemos a continuación: es la versión que puede compararse directamente con lo que Rosenbluth escribió, porque en 1950 los factores de Sachs no existían todavía.

\paragraph{Conexión Exacta con la Ecuación (1) de Rosenbluth (1950):}
Hemos llegado al momento que da sentido a la estructura entera de este artículo. La Ec.~\eqref{eq:rosenbluth_1950_original}, con la que abrimos la Sección~\ref{sec:punto_de_partida}, fue presentada allí sin justificación alguna: era simplemente la fórmula que Rosenbluth publicó en \textit{Physical Review} en agosto de 1950, escrita en una notación hoy en desuso, con dos cantidades efectivas $e'$ y $\kappa'$ cuyo significado microscópico él no podía conocer. Todo lo que hemos hecho desde entonces---construir el vértice más general compatible con las simetrías, demostrar la identidad de Gordon, aplicar la identidad de Ward, evaluar trazas de Dirac, cambiar al marco de Breit, integrar el espacio de fases---ha sido, visto retrospectivamente, un único cálculo cuyo resultado debería coincidir con aquella ecuación. Verifiquemos que coincide. Nótese que no se trata de una comprobación aproximada ni de un límite: la afirmación es que ambas expresiones son \emph{algebraicamente idénticas} una vez traducida la notación.

Es profundamente instructivo verificar cómo la formulación moderna de Dirac-Pauli, Ec.~\eqref{eq:rosenbluth_complete_dirac_sec5}, se reduce exactamente a la histórica Ecuación (1) publicada por Marshall Rosenbluth \cite{rosenbluth1950}. En la convención original de Rosenbluth, el factor de forma de Dirac se normalizó como la carga efectiva $F_1 \to e'/e$ y el de Pauli como el momento anómalo efectivo $F_2 \to \kappa' e'/e$. Sustituyendo estas cantidades e incorporando la corrección radiativa de Schwinger $(e''/e)^2$:
\begin{align}
\left(\frac{d\sigma}{d\Omega}\right) &= \left(\frac{e^2}{4\pi}\right)^2 \frac{\cos^2(\theta/2)}{4E^2 \sin^4(\theta/2) [1 + 2(E/M)\sin^2(\theta/2)]} \left(\frac{e' e''}{e^2}\right)^2 \nonumber \\
&\quad \times \left\{ 1 + \frac{(E^2/M^2)\sin^2(\theta/2)}{1 + 2(E/M)\sin^2(\theta/2)} \left[ {\kappa'}^2 + 2(1+\kappa')^2 \tan^2\left(\frac{\theta}{2}\right) \right] \right\}.
\end{align}
Notando que $\frac{\cos^2(\theta/2)}{\sin^4(\theta/2)} = \cot^2(\theta/2) \csc^2(\theta/2)$ y distribuyendo el denominador cinemático de retroceso $[1 + 2(E/M)\sin^2(\theta/2)]$ en el corchete:
\begin{equation}
\begin{aligned}
\sigma \, d\Omega &= \left(\frac{e' e''}{2E}\right)^2 \cot^2\left(\frac{\theta}{2}\right) \csc^2\left(\frac{\theta}{2}\right) \\
&\quad \times \left\{ \frac{1 + 2\left(\dfrac{E}{M}\right)\sin^2\left(\dfrac{\theta}{2}\right) + \left(\dfrac{E^2}{M^2}\right)\left[ 2(1+\kappa')^2 \tan^2\left(\dfrac{\theta}{2}\right)\sin^2\left(\dfrac{\theta}{2}\right) + {\kappa'}^2 \sin^2\left(\dfrac{\theta}{2}\right) \right]}{\left[1 + 2\left(\dfrac{E}{M}\right)\sin^2\left(\dfrac{\theta}{2}\right)\right]^2} \right\} d\Omega.
\end{aligned}
\end{equation}
¡Esta es con total fidelidad la célebre Ecuación (1) de Rosenbluth de 1950, Ec.~\eqref{eq:rosenbluth_1950_original}! Queda así rigurosamente demostrada la equivalencia analítica entre la teoría de perturbaciones covariante moderna y el cálculo fundacional de Rosenbluth. El círculo abierto en la primera página está cerrado: la ecuación que allí era un dato se ha convertido en un teorema. Merece la pena valorar lo que esa coincidencia significa en ambas direcciones. Por un lado valida nuestro desarrollo, pues una cadena tan larga de manipulaciones algebraicas difícilmente reproduciría por azar la estructura completa de una expresión ajena. Por otro, y más interesante, muestra hasta qué punto el contenido de la fórmula estaba ya determinado por las simetrías: Rosenbluth no disponía en 1950 ni del formalismo de los factores de forma de Sachs, ni del marco de Breit como herramienta sistemática, ni de la identidad de Ward en la forma que hemos usado, y sin embargo la estructura a la que llegó es exactamente la misma. Sus dos parámetros efectivos, $e'$ y $\kappa'$, eran los valores de $F_1$ y $F_2$ que la teoría no podía predecir entonces ni puede predecir ahora.

\subsection{Técnica de Separación Lineal y Reconstrucción Cinemática}
Disponer de la fórmula correcta no equivale todavía a poder medir. El problema experimental es el siguiente: cada medida de la sección eficaz a unas condiciones $(E,\theta)$ dadas produce \emph{un} número, y ese número depende de \emph{dos} incógnitas, $G_E(Q^2)$ y $G_M(Q^2)$. Una sola medida no puede determinar dos funciones. La idea de Rosenbluth---y esta es su contribución metodológica, independiente de la fórmula misma---consiste en explotar el hecho de que las dos incógnitas entran con dependencias angulares \emph{distintas}: $G_E^2$ entra sin ningún factor angular adicional y $G_M^2$ entra multiplicado por $\tan^2(\theta/2)$. Si se consigue variar el ángulo manteniendo $Q^2$ fijo---lo cual exige mover simultáneamente la energía del haz---las dos contribuciones se separan como el intercepto y la pendiente de una recta.

Para separar experimentalmente $G_E^2$ y $G_M^2$, se define la polarización longitudinal del fotón virtual:
\begin{equation}
\epsilon \equiv \left[ 1 + 2(1+\tau)\tan^2\left(\frac{\theta}{2}\right) \right]^{-1}, \qquad 0 < \epsilon \le 1.
\end{equation}
Esta variable, que a primera vista parece una abreviatura arbitraria, tiene un significado físico preciso: mide en qué grado el fotón virtual intercambiado está polarizado longitudinalmente. A ángulos frontales ($\theta \to 0$) se tiene $\epsilon \to 1$ y el fotón es máximamente longitudinal, favoreciendo la respuesta eléctrica; a ángulos traseros ($\theta \to 180^\circ$) se tiene $\epsilon \to 0$ y domina la componente transversal, es decir, la magnética. Barrer $\epsilon$ es, por tanto, barrer el carácter del sondeo.

Multiplicando la Ec.~\eqref{eq:rosenbluth_formula_box} por $(1+\tau)$ y dividiendo entre la sección de Mott, definimos la \textbf{Sección Eficaz Reducida}:
\begin{equation}
\boxed{\sigma_{\text{red}} \equiv \left( \frac{d\sigma}{d\Omega} \right) \left[ \left( \frac{d\sigma}{d\Omega} \right)_{\text{Mott}} \frac{1}{1+\tau} \right]^{-1} = G_E^2(Q^2) + x \, G_M^2(Q^2)},
\label{eq:sigma_red_linear}
\end{equation}
donde la variable angular de Rosenbluth es:
\begin{equation}
x \equiv \tau \left[ 1 + 2(1+\tau)\tan^2\left(\frac{\theta}{2}\right) \right] = \frac{\tau}{\epsilon}.
\end{equation}
La Ec.~\eqref{eq:sigma_red_linear} es la forma operativa de todo lo anterior y la razón de ser de esta sección: dividiendo el dato bruto por factores puramente cinemáticos y conocidos, la sección eficaz medida se convierte en una función \emph{lineal} de una variable controlable. Toda la física desconocida queda en los dos coeficientes de esa recta.

Despejando el ángulo de laboratorio $\theta$ en términos de $\epsilon$:
\begin{equation}
\frac{1}{\epsilon} - 1 = 2(1+\tau)\tan^2\left(\frac{\theta}{2}\right) \implies \tan^2\left(\frac{\theta}{2}\right) = \frac{1 - \epsilon}{2(1+\tau)\epsilon}.
\end{equation}
Para mantener fijado el valor de $Q^2$ mientras se barre $\epsilon$, la energía del haz incidente $E$ debe ajustarse según la ecuación cuadrática $Q^2 = 4 E E' \sin^2(\theta/2)$:
\begin{equation}
E = \frac{Q^2}{4 M \sin^2(\theta/2)} \left[ 1 + \sqrt{1 + \frac{4 M^2 \sin^2(\theta/2)}{Q^2}} \right].
\end{equation}
Esta ecuación es la receta experimental concreta: dice a qué energía hay que poner el acelerador para cada ángulo del espectrómetro si se quiere permanecer sobre el mismo valor de $Q^2$. Una separación de Rosenbluth completa consiste, pues, en repetir la medida a varias parejas $(E,\theta)$ ligadas por esta relación. La gráfica de $\sigma_{\text{red}}$ en función de $x$ define entonces una recta rigurosa:
\begin{itemize}
    \item \textbf{Intercepto ($x = 0$):} Determina directamente $G_E^2(Q^2)$.
    \item \textbf{Pendiente:} Determina directamente $G_M^2(Q^2)$.
\end{itemize}

\begin{figure}[t]
\centering
\includegraphics[width=0.92\textwidth]{figuras/rosenbluth_separation.pdf}
\caption{Simulación computacional de la separación lineal de Rosenbluth para cuatro valores fijos: $Q^2 = 0.5, 1.0, 2.0, 3.0\text{ GeV}^2$. Puntos sintéticos simulados con incertidumbres experimentales del $1.5\%$. Las rectas de mínimos cuadrados ponderados desacoplan el intercepto eléctrico $G_E^2$ de la pendiente magnética $G_M^2$. A medida que $Q^2$ crece, $\tau$ magnifica la pendiente, reduciendo drásticamente la precisión relativa del intercepto.}
\label{fig:rosenbluth_separation}
\end{figure}

\subsection{Resultados Cuantitativos de la Simulación}
Para comprobar que el procedimiento funciona---y, más importante, para cuantificar hasta dónde funciona---hemos simulado numéricamente el proceso completo: generar puntos sintéticos de $\sigma_{\text{red}}$ a partir de valores conocidos de $G_E$ y $G_M$, añadirles incertidumbres realistas del $1.5\%$, ajustar la recta por mínimos cuadrados ponderados y comparar los factores de forma recuperados con los de partida. En la Fig.~\ref{fig:rosenbluth_separation} se presentan las rectas de ajuste obtenidas mediante regresión lineal ponderada con matriz de covarianza. En la Tabla~\ref{tab:rosenbluth_results} se reportan los valores teóricos frente a los extraídos numéricamente.

\begin{table}[h]
\centering
\small
\caption{Resultados numéricos cuantitativos de la separación de Rosenbluth a distintas transferencias de momento.}
\label{tab:rosenbluth_results}
\begin{tabular}{ccccccc}
\toprule
$Q^2$ [$\text{GeV}^2$] & $\tau$ & $G_E^{\text{teor}}$ & $G_E^{\text{extr}}$ & $G_M^{\text{teor}}$ & $G_M^{\text{extr}}$ & $\chi^2/\text{ndf}$ \\
\midrule
0.50 & 0.142 & 0.344 & $0.354 \pm 0.007$ & 0.962 & $0.952 \pm 0.009$ & 0.77 \\
1.00 & 0.284 & 0.172 & $0.167 \pm 0.006$ & 0.481 & $0.486 \pm 0.004$ & 0.23 \\
2.00 & 0.568 & 0.069 & $0.073 \pm 0.004$ & 0.192 & $0.190 \pm 0.001$ & 1.24 \\
3.00 & 0.852 & 0.037 & $0.031 \pm 0.004$ & 0.102 & $0.103 \pm 0.001$ & 1.26 \\
\bottomrule
\end{tabular}
\end{table}

\textbf{Análisis de incertidumbre:} A $Q^2 = 0.5\text{ GeV}^2$, el error relativo en $G_E$ es de un $2.0\%$. A $Q^2 = 3.0\text{ GeV}^2$, el error relativo se dispara al $12.9\%$. Esto evidencia la pérdida intrínseca de sensibilidad de la técnica de Rosenbluth para medir $G_E$ a transferencias moderadas y altas debido al factor multiplicativo $\tau \approx 0.85$~\cite{arrington2007,perdrisat2007}.

Conviene entender el mecanismo de esta degradación, porque no es un defecto de la simulación ni de la calidad de los datos, sino una limitación estructural del método que tendrá consecuencias históricas. En la Ec.~\eqref{eq:sigma_red_linear}, la contribución magnética entra multiplicada por la variable $x = \tau/\epsilon$. Al crecer $Q^2$, el parámetro $\tau$ aumenta proporcionalmente y la pendiente de la recta se magnifica, mientras que el intercepto---que es $G_E^2$---se vuelve una fracción cada vez menor de la señal total. Dicho de otro modo: a $Q^2$ alto se está intentando determinar un intercepto muy pequeño extrapolando una recta de pendiente muy pronunciada, y en esa situación una incertidumbre sistemática modesta en la normalización de la sección eficaz se traduce en un error enorme en $G_E$. Los números de la Tabla~\ref{tab:rosenbluth_results} lo muestran con claridad: un $1.5\%$ de error en los datos produce un $2\%$ de error en $G_E$ a $Q^2 = 0.5\text{ GeV}^2$, pero casi un $13\%$ a $Q^2 = 3\text{ GeV}^2$, y la tendencia empeora rápidamente. Esta es exactamente la razón por la que, a partir del año 2000, los experimentos de Jefferson Lab abandonaron la separación de Rosenbluth como método de referencia a $Q^2$ elevado y adoptaron una técnica completamente distinta, la transferencia de polarización~\cite{jones2000,punjabi2005,gayou2002}, que mide el cociente $G_E/G_M$ de forma directa y no está sujeta a esta supresión. Esa técnica, y la sorpresa que deparó al compararse con los resultados de Rosenbluth, ocuparán la Sección~\ref{sec:radio_fourier}.


La separación lineal de Rosenbluth permite extraer con precisión cuantitativa los factores de Sachs $G_E(Q^2)$ y $G_M(Q^2)$ como funciones de la transferencia de momento. Sin embargo, el objetivo último de la física hadrónica no se agota en tabular funciones en el espacio recíproco de momentos, sino en desvelar la distribución geométrica real de carga y magnetización en el interior del nucleón. Para conectar $G_E(Q^2)$ con la densidad espacial estática $\rho(\mathbf{r})$ y deducir analíticamente una de las cantidades más fundamentales de la física subatómica---el radio cuadrático medio de carga del protón $\langle r_E^2 \rangle^{1/2}$---debemos tender un puente matemático hacia el espacio de coordenadas mediante la transformada inversa de Fourier tridimensional y el análisis complejo de residuos en el plano de Cauchy.
\section{Transformada de Fourier, Teorema de Residuos y el Radio de Carga del Protón}
\label{sec:radio_fourier}

Hasta este punto hemos aprendido a medir dos funciones de $Q^2$, pero una función de la transferencia de momento no es todavía lo que la pregunta original pedía. Lo que motivaba el artículo entero---la anomalía de Stern, el protón que ocupa volumen---era una cuestión sobre el espacio real: cuánto mide el protón. Esta sección tiende el puente entre ambos lenguajes. Consta de dos mitades. La primera es matemática y clásica: demostrar que $G_E(Q^2)$ es la transformada de Fourier de la densidad espacial de carga, extraer de ahí el radio cuadrático medio como una derivada en el origen, e invertir analíticamente el modelo dipolar mediante el teorema de residuos de Cauchy para obtener la forma explícita de $\rho(r)$. La segunda es contemporánea y menos plácida: la técnica de Rosenbluth, cuya demostración acabamos de completar, resultó entrar en conflicto con un método experimental alternativo desarrollado en Jefferson Lab a partir del año 2000, y la resolución de ese conflicto obligó a revisar una aproximación que habíamos dado por buena desde la Sección~\ref{sec:vertice_gordon}: la de un solo fotón intercambiado.

\subsection{Deducción Analítica del Radio Cuadrático Medio}
La relación entre un factor de forma y una distribución espacial no es exclusiva de la física de partículas: es el mismo fenómeno que gobierna la difracción de rayos X por un cristal o la difracción de la luz por una rendija. Siempre que se sondea un objeto extendido con una onda, la amplitud dispersada en función del momento transferido resulta ser la transformada de Fourier de la densidad del objeto, y la razón es puramente geométrica: las contribuciones procedentes de puntos distintos del objeto llegan al detector con fases distintas, y esa diferencia de fase es exactamente $e^{i\mathbf{q}\cdot\mathbf{r}}$. Lo que hace especial al caso presente es que la identificación requiere un marco de referencia concreto. Solo en el sistema de Breit, donde según vimos $Q^2 = \mathbf{q}^2$ es el cuadrado de un trimomento ordinario y no hay transferencia de energía, tiene sentido hablar de una densidad de carga \emph{estática}; por eso la interpretación que sigue es rigurosa en el límite no relativista $Q^2 \ll M^2$, que es precisamente el régimen en el que se mide el radio.

En el límite no relativista ($Q^2 \ll M^2$), $G_E(Q^2)$ representa la transformada de Fourier tridimensional de la densidad estática de carga eléctrica $\rho(\mathbf{r})$ \cite{hofstadter1956,hofstadter1961,das2003,perdrisat2007}:
\begin{equation}
G_E(Q^2) = \int d^3\mathbf{r} \, \rho(r) e^{i \mathbf{q} \cdot \mathbf{r}}.
\end{equation}
Efectuando la integración angular en coordenadas esféricas con $\mathbf{q} \cdot \mathbf{r} = q r \cos\theta_r$:
\begin{align}
\int d\Omega_r e^{i q r \cos\theta_r} &= 2\pi \int_{-1}^1 e^{i q r u} du = 2\pi \left[ \frac{e^{i q r u}}{i q r} \right]_{-1}^1 \nonumber \\
&= 2\pi \frac{e^{i q r} - e^{-i q r}}{i q r} = 4\pi \frac{\sin(q r)}{q r}.
\end{align}
Por lo tanto:
\begin{equation}
G_E(Q^2) = 4\pi \int_0^\infty r^2 \rho(r) \frac{\sin(q r)}{q r} dr.
\end{equation}
La estrategia para extraer un radio de esta expresión es la siguiente. Medir $G_E$ en todo el rango de $Q^2$ y hacer la transformada inversa daría $\rho(r)$ completa, pero eso exigiría datos hasta $Q^2$ infinito. Si en cambio nos conformamos con una sola característica de la distribución---su anchura---basta con examinar la región $Q^2 \to 0$, porque allí la exponencial apenas oscila y la integral solo es sensible a los primeros momentos de $\rho$. Formalizamos esa idea con un desarrollo en serie. Efectuamos la expansión de Maclaurin de la función $\sin(x)/x$:
\begin{equation}
\frac{\sin(q r)}{q r} = 1 - \frac{1}{3!} (q r)^2 + \frac{1}{5!} (q r)^4 - \dots = 1 - \frac{1}{6} Q^2 r^2 + \mathcal{O}(Q^4).
\end{equation}
Sustituyendo en la integral:
\begin{align}
G_E(Q^2) &= 4\pi \int_0^\infty r^2 \rho(r) \left[ 1 - \frac{1}{6} Q^2 r^2 + \mathcal{O}(Q^4) \right] dr \nonumber \\
&= \int d^3\mathbf{r} \, \rho(r) - \frac{1}{6} Q^2 \int d^3\mathbf{r} \, r^2 \rho(r) + \mathcal{O}(Q^4).
\end{align}
Los dos primeros coeficientes del desarrollo han resultado ser precisamente los dos primeros momentos de la distribución: la carga total y el radio cuadrático medio. Imponiendo la normalización de la carga total $\int d^3\mathbf{r} \rho(r) = 1$ y definiendo el radio de carga cuadrático medio $\langle r_E^2 \rangle \equiv \int d^3\mathbf{r} \, r^2 \rho(r)$:
\begin{equation}
G_E(Q^2) = 1 - \frac{1}{6} \langle r_E^2 \rangle Q^2 + \mathcal{O}(Q^4).
\end{equation}
Derivando formalmente respecto a $Q^2$ en el origen cinemático:
\begin{equation}
\left. \frac{d G_E(Q^2)}{d Q^2} \right|_{Q^2 = 0} = -\frac{1}{6} \langle r_E^2 \rangle,
\end{equation}
de donde se concluye unívocamente:
\begin{equation}
\boxed{\langle r_E^2 \rangle = -6 \left. \frac{d G_E(Q^2)}{d Q^2} \right|_{Q^2 = 0}}
\label{eq:proton_radius_analytic_box}
\end{equation}

Conviene ser preciso sobre lo que esta ecuación afirma, porque es un resultado de notable economía. No dice que haga falta conocer $\rho(r)$ para obtener el radio; dice exactamente lo contrario. El radio cuadrático medio del protón es, ni más ni menos, la pendiente de la curva experimental $G_E(Q^2)$ evaluada en el origen, multiplicada por $-6$. No se necesita modelo alguno de la estructura interna, ni una hipótesis sobre la forma funcional de la distribución: el tamaño del protón es una derivada de datos medibles. Aquí se cierra por fin la cadena que arrancó con la anomalía de Stern. Esa anomalía indicaba que el protón tiene estructura; las Secciones~\ref{sec:vertice_gordon} a~\ref{sec:marco_breit} demostraron que la estructura cabe en dos funciones; la Sección~\ref{sec:seccion_rosenbluth} mostró cómo medirlas; y la Ec.~\eqref{eq:proton_radius_analytic_box} convierte una de ellas en una longitud.

\begin{figure}[t]
\centering
\includegraphics[width=\columnwidth]{figuras/form_factors_dipole.pdf}
\caption{(a) Factores de forma de Sachs normalizados frente al modelo dipolar estándar $G_D(Q^2) = (1 + Q^2/0.71)^{-2}$~\cite{hofstadter1956}. (b) Desviaciones relativas respecto al dipolo comparadas con la parametrización de Kelly~\cite{kelly2004}.}
\label{fig:form_factors_dipole}
\end{figure}

\subsection{El Modelo Dipolar y Cálculo de Residuos en el Plano Complejo}
La subsección anterior extrajo un número---el radio---de la pendiente en el origen. Podemos ser más ambiciosos y preguntar por la función completa $\rho(r)$, pero para ello hace falta conocer $G_E(Q^2)$ en todo el rango de transferencias, y eso obliga a adoptar una parametrización de los datos. La que se emplea históricamente, y que se ajusta a las medidas con notable fidelidad en un rango amplio, es la forma dipolar de Hofstadter. Invertir su transformada de Fourier tiene interés por partida doble: proporciona la forma espacial explícita de la distribución de carga del protón y, de paso, permite verificar la Ec.~\eqref{eq:proton_radius_analytic_box} por una vía completamente independiente. El cálculo requiere integración de contorno en el plano complejo, que resumimos con detalle porque el resultado es especialmente limpio.

Empíricamente, los factores de forma escalan según la parametrización dipolar de Hofstadter \cite{hofstadter1956}:
\begin{equation}
G_D(Q^2) = \frac{1}{\left( 1 + \frac{Q^2}{M_V^2} \right)^2} = \frac{M_V^4}{(Q^2 + M_V^2)^2}, \quad M_V^2 \approx 0.71\text{ GeV}^2.
\end{equation}
Para deducir rigurosamente la distribución espacial $\rho(r)$, calculamos la transformada inversa de Fourier:
\begin{equation}
\rho(r) = \frac{1}{(2\pi)^3} \int d^3\mathbf{q} \, G_D(\mathbf{q}^2) e^{-i \mathbf{q} \cdot \mathbf{r}}.
\end{equation}
Integrando la variable angular en coordenadas esféricas:
\begin{equation}
\rho(r) = \frac{1}{2\pi^2 r} \int_0^\infty q \sin(qr) \frac{M_V^4}{(q^2 + M_V^2)^2} dq.
\end{equation}
Esta integral no admite una primitiva elemental, pero sí puede evaluarse exactamente por métodos de variable compleja. El procedimiento estándar consta de tres pasos: extender el dominio de integración a toda la recta real aprovechando una paridad, sustituir el seno por una exponencial compleja, y cerrar el contorno en el semiplano donde la exponencial decae. Notemos que el integrando $f(q) = \frac{q \sin(qr)}{(q^2 + M_V^2)^2}$ es estrictamente una \textbf{función par} de $q$, ya que $(-q)\sin(-qr) = q\sin(qr)$. Por lo tanto:
\begin{equation}
\int_0^\infty f(q) dq = \frac{1}{2} \int_{-\infty}^\infty f(q) dq.
\end{equation}
Escribiendo $\sin(qr) = \text{Im}[e^{iqr}]$:
\begin{equation}
\rho(r) = \frac{M_V^4}{4\pi^2 r} \text{Im} \left[ \int_{-\infty}^\infty \frac{q e^{iqr}}{(q^2 + M_V^2)^2} dq \right].
\end{equation}
Definimos la función analítica $g(z) \equiv \frac{z e^{izr}}{(z^2 + M_V^2)^2} = \frac{z e^{izr}}{(z - iM_V)^2 (z + iM_V)^2}$. Para $r > 0$, cerramos el contorno de integración en el semiplano superior $\text{Im}(z) > 0$ mediante un semicírculo $C_R$ de radio $R \to \infty$. Por el Lema de Jordan, la integral sobre $C_R$ se anula.

La única singularidad en el semiplano superior es un \textbf{polo doble} en $z = +i M_V$. Que el polo sea doble, y no simple, es consecuencia directa del exponente $2$ del modelo dipolar, y es lo que producirá el factor lineal en $r$ que aparecerá en el residuo. El residuo en un polo de orden 2 se calcula mediante:
\begin{equation}
\text{Res}[g, iM_V] = \lim_{z \to iM_V} \frac{d}{dz} \left[ (z - iM_V)^2 g(z) \right] = \lim_{z \to iM_V} \frac{d}{dz} \left[ \frac{z e^{izr}}{(z + iM_V)^2} \right].
\end{equation}
Evaluemos la derivada del cociente minuciosamente:
\begin{align}
\frac{d}{dz} \left[ \frac{z e^{izr}}{(z + iM_V)^2} \right] &= \frac{\left( e^{izr} + i z r e^{izr} \right)(z + iM_V)^2 - z e^{izr} \cdot 2(z + iM_V)}{(z + iM_V)^4} \nonumber \\
&= \frac{(1 + i z r)(z + iM_V) - 2z}{(z + iM_V)^3} e^{izr}.
\end{align}
Evaluando en $z = i M_V$:
\begin{align}
z + iM_V &= 2i M_V, \qquad (z + iM_V)^3 = (2i M_V)^3 = -8i M_V^3, \\
e^{izr} &= e^{i(iM_V)r} = e^{-M_V r}.
\end{align}
Obsérvese que la exponencial compleja se ha convertido en una exponencial real decreciente al evaluarse sobre el eje imaginario: ahí está, ya visible, el origen del perfil exponencial que obtendremos. El numerador evaluado en $z = i M_V$ es:
\begin{align}
(1 + i(iM_V)r)(2i M_V) - 2(iM_V) &= (1 - M_V r)(2i M_V) - 2i M_V \nonumber \\
&= 2i M_V - 2i M_V^2 r - 2i M_V \nonumber \\
&= - 2i M_V^2 r.
\end{align}
Por lo tanto, el residuo es:
\begin{equation}
\text{Res}[g, iM_V] = \frac{-2i M_V^2 r}{-8i M_V^3} e^{-M_V r} = \frac{r}{4 M_V} e^{-M_V r}.
\end{equation}
Aplicando el Teorema de los Residuos de Cauchy:
\begin{equation}
\oint g(z) dz = 2\pi i \, \text{Res}[g, iM_V] = 2\pi i \left( \frac{r}{4 M_V} e^{-M_V r} \right) = \frac{i \pi r}{2 M_V} e^{-M_V r}.
\end{equation}
Tomando la parte imaginaria:
\begin{equation}
\text{Im} \left[ \oint g(z) dz \right] = \frac{\pi r}{2 M_V} e^{-M_V r}.
\end{equation}
Sustituyendo en la expresión para $\rho(r)$:
\begin{equation}
\rho(r) = \frac{M_V^4}{4\pi^2 r} \left[ \frac{\pi r}{2 M_V} e^{-M_V r} \right] = \frac{M_V^3}{8\pi} e^{-M_V r}.
\end{equation}
¡El modelo dipolar describe una \textbf{distribución espacial exponencial decreciente}!
\begin{equation}
\dest{\rho(r) = \frac{M_V^3}{8\pi} e^{-M_V r} \equiv \rho_0 e^{-M_V r}}
\end{equation}
El resultado tiene una lectura física inmediata. El protón no posee una superficie definida ni un borde abrupto: su carga se desvanece de forma exponencial con una longitud característica $1/M_V$, de manera análoga a como decae la función de onda de un estado ligado fuera del pozo de potencial que lo confina. Un dipolo en el espacio de momentos equivale a una exponencial en el espacio de coordenadas, del mismo modo que una gaussiana equivale a otra gaussiana; la correspondencia entre la forma funcional medida y la geometría real es, en este sentido, directa.

Verifiquemos la normalización integrando en todo el espacio tridimensional:
\begin{equation}
\int d^3\mathbf{r} \rho(r) = 4\pi \int_0^\infty r^2 \left( \frac{M_V^3}{8\pi} e^{-M_V r} \right) dr = \frac{M_V^3}{2} \int_0^\infty r^2 e^{-M_V r} dr.
\end{equation}
Recordando que $\int_0^\infty r^n e^{-\lambda r} dr = \frac{n!}{\lambda^{n+1}}$, para $n=2$ y $\lambda = M_V$:
\begin{equation}
\int d^3\mathbf{r} \rho(r) = \frac{M_V^3}{2} \left( \frac{2!}{M_V^3} \right) = 1.
\end{equation}
La carga total sale exactamente la unidad, como debía, lo que confirma que la inversión por residuos se ha hecho con los factores correctos.

Calculamos la derivada analítica del dipolo respecto a $Q^2$:
\begin{equation}
\frac{d G_D}{d Q^2} = -\frac{2}{M_V^2 \left( 1 + \frac{Q^2}{M_V^2} \right)^3} \implies \left. \frac{d G_D}{d Q^2} \right|_{Q^2 = 0} = -\frac{2}{M_V^2}.
\end{equation}
Sustituyendo en la Ec.~\eqref{eq:proton_radius_analytic_box}:
\begin{equation}
\langle r_E^2 \rangle_{\text{dipolo}} = -6 \left( -\frac{2}{M_V^2} \right) = \frac{12}{M_V^2}.
\label{eq:rE2_momento_deriv}
\end{equation}

\textbf{Comprobación directa por integración en el espacio de coordenadas:}\\
Disponemos ahora de la densidad explícita, de modo que podemos calcular el mismo radio por el camino opuesto---integrando en el espacio real en lugar de derivar en el espacio de momentos---y comprobar que ambos coinciden. Es una verificación no trivial de toda la cadena. Evaluemos de forma independiente el valor esperado del radio al cuadrado integrando directamente la densidad exponencial $\rho(r)$:
\begin{equation}
\langle r_E^2 \rangle_{\text{espacial}} = \int d^3\mathbf{r} \, r^2 \rho(r) = 4\pi \int_0^\infty r^4 \left( \frac{M_V^3}{8\pi} e^{-M_V r} \right) dr = \frac{M_V^3}{2} \int_0^\infty r^4 e^{-M_V r} dr.
\end{equation}
Aplicando la fórmula integral canónica para potencias de la función gamma:
\begin{equation}
\int_0^\infty r^n e^{-\lambda r} dr = \frac{n!}{\lambda^{n+1}},
\end{equation}
para $n = 4$ y $\lambda = M_V$:
\begin{equation}
\int_0^\infty r^4 e^{-M_V r} dr = \frac{4!}{M_V^5} = \frac{24}{M_V^5}.
\end{equation}
Multiplicando por el factor prefactor:
\begin{equation}
\langle r_E^2 \rangle_{\text{espacial}} = \frac{M_V^3}{2} \left( \frac{24}{M_V^5} \right) = \frac{12}{M_V^2}.
\end{equation}
¡Ambos métodos (la derivada analítica en el espacio de momentos y la integración directa en el espacio de coordenadas) coinciden con absoluta exactitud matemática! La coincidencia confirma simultáneamente la inversión por residuos y la identidad $\langle r_E^2\rangle = -6\,dG_E/dQ^2|_0$, que fueron obtenidas por caminos enteramente independientes.

Evaluando numéricamente con el parámetro estándar $M_V^2 = 0.71\text{ GeV}^2$:
\begin{equation}
\langle r_E^2 \rangle = \frac{12}{0.71\text{ GeV}^2} \approx 16.9014\text{ GeV}^{-2}.
\end{equation}
Empleando el factor de conversión relativista exacto con $\hbar c = 0.197327\text{ GeV}\cdot\text{fm}$:
\begin{equation}
(\hbar c)^2 = (0.197327)^2\text{ GeV}^2\cdot\text{fm}^2 = 0.038938\text{ GeV}^2\cdot\text{fm}^2,
\end{equation}
\begin{equation}
\langle r_E^2 \rangle = 16.9014\text{ GeV}^{-2} \times 0.038938\text{ GeV}^2\cdot\text{fm}^2 = 0.65811\text{ fm}^2,
\end{equation}
extrayendo la raíz cuadrada:
\begin{equation}
\dest{r_p^{\text{dipolo}} = \sqrt{\langle r_E^2 \rangle} = \sqrt{0.65811\text{ fm}^2} \approx \mathbf{0.8112\text{ fm}}}
\end{equation}

Este número merece un comentario honesto, porque no coincide exactamente con el valor aceptado que citaremos al final de la sección, y la diferencia es instructiva. El valor $0.8112\text{ fm}$ no es una medida del radio del protón: es el radio que tendría un protón cuyo factor de forma fuese \emph{exactamente} un dipolo con $M_V^2 = 0.71\text{ GeV}^2$. El dipolo es una parametrización sencilla y notablemente buena en un rango amplio de $Q^2$, pero los datos reales se apartan de él en un pequeño porcentaje, como muestra el panel (b) de la Figura~\ref{fig:form_factors_dipole}, y la Ec.~\eqref{eq:proton_radius_analytic_box} es sensible precisamente a la pendiente en el origen, que es donde esas desviaciones pesan más. La discrepancia de un $4\%$ entre $0.811\text{ fm}$ y el valor de consenso no cuestiona, por tanto, ninguno de los desarrollos anteriores; cuantifica la calidad del modelo dipolar. Determinar el radio con precisión real exige ajustar la pendiente a los datos sin imponer una forma funcional, y eso nos lleva directamente a las dificultades experimentales de las dos últimas subsecciones.

\subsection{Deducción de la Transferencia de Polarización (\texorpdfstring{$P_t/P_l$}{Pt/Pl})}
La Sección~\ref{sec:seccion_rosenbluth} terminó con un diagnóstico incómodo: a transferencias de momento elevadas, la separación de Rosenbluth pierde sensibilidad a $G_E$ porque el término magnético domina la recta y el intercepto se vuelve indistinguible de cero. La solución no consiste en mejorar la estadística sino en medir otra cosa. Si en lugar de la sección eficaz total se mide la \emph{polarización} del protón de retroceso cuando el haz incidente está longitudinalmente polarizado, resulta que las dos componentes de esa polarización dependen de $G_E$ y $G_M$ de forma distinta, y su cociente es directamente proporcional a $G_E/G_M$. La ventaja es doble: el cociente elimina las normalizaciones absolutas del haz y del detector, que son la principal fuente de error sistemático, y no está sujeto a la supresión por $\tau$ que arruinaba la extracción de $G_E$. El precio es experimental---hace falta un haz polarizado y un polarímetro para el protón de retroceso---y por eso la técnica no fue viable hasta finales de los años noventa.

En los experimentos de \textbf{transferencia de polarización} en Jefferson Lab ($\vec{e} p \to e \vec{p}$) \cite{jones2000,punjabi2005}, electrones longitudinalmente polarizados con helicidad $h = \pm 1$ inciden sobre un blanco no polarizado, transfiriendo su polarización al protón en retroceso. El tensor leptónico polarizado adquiere una componente antisimétrica:
\begin{equation}
L_{\mu\nu}^{(h)} = L_{\mu\nu} + i h \epsilon_{\mu\nu\alpha\beta} k^\alpha k'^\beta.
\end{equation}
Nótese que este término antisimétrico es precisamente el que se anulaba al promediar sobre helicidades en la Sección~\ref{sec:cinematica_trazas}; polarizar el haz equivale a no promediar, y el término reaparece. Al contraer con el tensor hadrónico de espín polarizado $W^{\mu\nu}(s)$, surgen dos componentes de polarización en el plano de dispersión (formalismo de Dombey y Akhiezer-Rekalo \cite{perdrisat2007}):
\begin{align}
I_0 P_t &= - 2 \sqrt{\tau(1+\tau)} G_E(Q^2) G_M(Q^2) \tan\left(\frac{\theta}{2}\right), \label{eq:Pt_eval} \\
I_0 P_l &= \frac{E + E'}{M} \sqrt{\tau(1+\tau)} G_M^2(Q^2) \tan^2\left(\frac{\theta}{2}\right), \label{eq:Pl_eval}
\end{align}
donde $I_0 = G_E^2 + \frac{\tau}{\epsilon} G_M^2$ es la sección diferencial no polarizada. Obsérvese el hecho clave: la componente transversal contiene el producto $G_EG_M$ y la longitudinal contiene $G_M^2$. Lo que en la sección eficaz no polarizada era un inconveniente---la aparición de productos cruzados---aquí se vuelve una ventaja, porque al dividir una por otra el factor $G_M$ se cancela parcialmente y queda el cociente que buscamos. Dividiendo miembro a miembro la Ec.~\eqref{eq:Pt_eval} entre la Ec.~\eqref{eq:Pl_eval}, el término $I_0$ y el factor común $\sqrt{\tau(1+\tau)}$ se cancelan exactamente:
\begin{equation}
\frac{P_t}{P_l} = \frac{- 2 \sqrt{\tau(1+\tau)} G_E G_M \tan(\theta/2)}{\frac{E + E'}{M} \sqrt{\tau(1+\tau)} G_M^2 \tan^2(\theta/2)} = - \frac{2 M}{(E + E')} \frac{1}{\tan(\theta/2)} \frac{G_E(Q^2)}{G_M(Q^2)}.
\label{eq:Pt_Pl_ratio_intermediate}
\end{equation}
La cancelación de $I_0$ es lo que hace poderosa a esta técnica: $I_0$ contiene toda la dependencia en la luminosidad, la eficiencia del detector y la normalización absoluta, y desaparece sin dejar rastro. Despejando el cociente de factores de forma:
\begin{equation}
\frac{G_E(Q^2)}{G_M(Q^2)} = - \frac{P_t}{P_l} \left[ \frac{E + E'}{2M} \tan\left(\frac{\theta}{2}\right) \right].
\label{eq:GE_GM_recoil_kin}
\end{equation}

\textbf{Deducción microscópica del factor cinemático en términos de $\tau$ y $\epsilon$:}\\
El corchete de la Ec.~\eqref{eq:GE_GM_recoil_kin} está escrito en variables de laboratorio, lo cual es incómodo para comparar experimentos realizados a energías distintas. Demostremos que puede reescribirse enteramente en términos de los dos invariantes $\tau$ y $\epsilon$, con lo que el resultado adquiere una forma universal. Demostremos analíticamente cómo se expresa la combinación $\frac{E+E'}{2M}\tan(\theta/2)$ en términos de los invariantes relativistas. De la conservación de energía-momento:
\begin{equation}
E - E' = \frac{Q^2}{2M} = 2 M \tau \implies (E - E')^2 = 4 M^2 \tau^2.
\end{equation}
Por definición de la transferencia de momento:
\begin{equation}
Q^2 = 4 E E' \sin^2\left(\frac{\theta}{2}\right) = 4 M^2 \tau \implies 4 E E' = \frac{4 M^2 \tau}{\sin^2(\theta/2)}.
\end{equation}
Escribiendo el cuadrado de la suma de energías como $(E+E')^2 = (E-E')^2 + 4EE'$:
\begin{equation}
(E + E')^2 = 4 M^2 \tau^2 + \frac{4 M^2 \tau}{\sin^2(\theta/2)} = 4 M^2 \tau \left[ \tau + \frac{1}{\sin^2(\theta/2)} \right].
\end{equation}
Dividiendo entre $4M^2$:
\begin{equation}
\left( \frac{E + E'}{2M} \right)^2 = \tau \left[ \tau + \frac{1}{\sin^2(\theta/2)} \right].
\label{eq:energy_sum_squared}
\end{equation}
Relacionemos ahora $\sin^2(\theta/2)$ con la polarización longitudinal del fotón virtual $\epsilon$. Por definición:
\begin{equation}
\frac{1}{\epsilon} = 1 + 2(1+\tau)\tan^2\left(\frac{\theta}{2}\right) \implies \tan^2\left(\frac{\theta}{2}\right) = \frac{1 - \epsilon}{2(1+\tau)\epsilon}.
\label{eq:tan2_epsilon_rel}
\end{equation}
Utilizando la identidad trigonométrica $\frac{1}{\sin^2(\theta/2)} = 1 + \cot^2(\theta/2) = 1 + \frac{1}{\tan^2(\theta/2)}$:
\begin{equation}
\frac{1}{\sin^2(\theta/2)} = 1 + \frac{2(1+\tau)\epsilon}{1 - \epsilon} = \frac{(1 - \epsilon) + 2\epsilon + 2\tau\epsilon}{1 - \epsilon} = \frac{1 + \epsilon + 2\tau\epsilon}{1 - \epsilon}.
\end{equation}
Sustituyendo en el corchete de la Ec.~\eqref{eq:energy_sum_squared}:
\begin{align}
\tau + \frac{1}{\sin^2(\theta/2)} &= \tau + \frac{1 + \epsilon + 2\tau\epsilon}{1 - \epsilon} = \frac{\tau(1 - \epsilon) + 1 + \epsilon + 2\tau\epsilon}{1 - \epsilon} \nonumber \\
&= \frac{\tau - \tau\epsilon + 1 + \epsilon + 2\tau\epsilon}{1 - \epsilon} = \frac{1 + \tau + \epsilon(1 + \tau)}{1 - \epsilon} \nonumber \\
&= \frac{(1 + \tau)(1 + \epsilon)}{1 - \epsilon}.
\end{align}
Por lo tanto:
\begin{equation}
\left( \frac{E + E'}{2M} \right)^2 = \tau \left[ \frac{(1 + \tau)(1 + \epsilon)}{1 - \epsilon} \right].
\end{equation}
Multiplicando miembro a miembro por $\tan^2(\theta/2)$ dado por la Ec.~\eqref{eq:tan2_epsilon_rel}:
\begin{align}
\left[ \frac{E + E'}{2M} \tan\left(\frac{\theta}{2}\right) \right]^2 &= \left[ \tau \frac{(1 + \tau)(1 + \epsilon)}{1 - \epsilon} \right] \left[ \frac{1 - \epsilon}{2(1 + \tau)\epsilon} \right] \nonumber \\
&= \frac{\tau(1 + \epsilon)}{2\epsilon}.
\end{align}
¡Los factores $(1-\epsilon)$ y $(1+\tau)$ se cancelan simultánea y rigurosamente! Extrayendo la raíz cuadrada:
\begin{equation}
\dest{\frac{E + E'}{2M} \tan\left(\frac{\theta}{2}\right) = \sqrt{\frac{\tau(1+\epsilon)}{2\epsilon}}}
\label{eq:factor_cinematico_polarizacion_exact}
\end{equation}
Sustituyendo la Ec.~\eqref{eq:factor_cinematico_polarizacion_exact} en la Ec.~\eqref{eq:GE_GM_recoil_kin}, arribamos a la fórmula canónica exacta:
\begin{equation}
\dest{\frac{G_E(Q^2)}{G_M(Q^2)} = - \frac{P_t}{P_l} \sqrt{\frac{\tau(1+\epsilon)}{2\epsilon}}} \iff \dest{\frac{P_t}{P_l} = - \sqrt{\frac{2\epsilon}{\tau(1+\epsilon)}} \frac{G_E(Q^2)}{G_M(Q^2)}}
\label{eq:polarization_ratio_exact}
\end{equation}
¡Esta técnica mide directamente el cociente $G_E/G_M$ sin requerir la resta de un término magnético abrumador ($\tau G_M^2$) ni depender de la normalización absoluta del haz! Compárese la situación con la de la separación de Rosenbluth: allí $G_E^2$ había de extraerse como un intercepto pequeño bajo una pendiente grande, mientras que aquí $G_E/G_M$ se obtiene de un cociente de dos cantidades medidas en el mismo suceso. Dos métodos independientes, ambos rigurosos, midiendo la misma cantidad física: la situación ideal para una comprobación cruzada. El resultado de esa comprobación, sin embargo, fue el contrario del esperado.

\begin{figure}[t]
\centering
\includegraphics[width=\columnwidth]{figuras/polarization_tpe_discrepancy.pdf}
\caption{Discrepancia entre la separación de Rosenbluth tradicional (cuadrados rojos) y los datos de polarización de JLab Hall A (círculos azules). Las correcciones de intercambio de dos fotones (TPE, línea discontinua verde) introducen una dependencia en $\epsilon$ que resuelve la aparente contradicción.}
\label{fig:polarization_tpe}
\end{figure}

\subsection{El Intercambio de Dos Fotones (TPE) y el Puzzle del Radio}
Cuando a partir del año 2000 los experimentos de transferencia de polarización de Jefferson Lab publicaron sus primeros resultados a $Q^2$ elevado, el cociente $G_E/G_M$ medido resultó ser sistemáticamente menor que el que se deducía de las separaciones de Rosenbluth existentes, y la discrepancia crecía con $Q^2$ hasta hacerse muy superior a las barras de error de ambos conjuntos de datos (Figura~\ref{fig:polarization_tpe}). Puesto que las dos técnicas miden la misma cantidad y ambas derivaciones son algebraicamente correctas---las hemos hecho las dos en este artículo---el desacuerdo no podía resolverse revisando el álgebra. Tenía que fallar una hipótesis compartida por ambas, y solo había una candidata: la aproximación de un único fotón intercambiado que adoptamos en la Sección~\ref{sec:vertice_gordon} al escribir la amplitud de Born, Ec.~\eqref{eq:feynman_amplitude_full}, y que nunca volvimos a cuestionar.

La inclusión de amplitudes de \textbf{Intercambio de Dos Fotones} (TPE) a orden $\mathcal{O}(\alpha^2)$ \cite{blunden2003,arrington2007,carlson2015} añade una corrección $\delta_{2\gamma}(\epsilon, Q^2)$ a la sección de Rosenbluth:
\begin{equation}
\sigma_{\text{red}} = G_E^2 + \frac{\tau}{\epsilon} G_M^2 + 2 G_M^2 \text{Re}(\delta_{2\gamma}).
\end{equation}
La pendiente en $\epsilon$ distorsiona la extracción lineal de Rosenbluth inflando artificialmente $G_E^2$, pero deja inmune la razón de polarizaciones $P_t/P_l$. El detalle decisivo es que $\delta_{2\gamma}$ depende de $\epsilon$, es decir, del ángulo: una corrección que dependiera solo de $Q^2$ se absorbería en la normalización y pasaría inadvertida, pero una que varía con el ángulo se disfraza de pendiente y contamina exactamente la magnitud que la separación de Rosenbluth intenta medir. Resulta además que la contaminación afecta de forma desproporcionada al término pequeño, $G_E^2$, y de forma despreciable al grande, lo que explica que el efecto solo se hiciera visible a $Q^2$ alto y que las medidas de polarización, al ser cocientes en un mismo suceso, quedaran esencialmente inmunes. Al incorporar teóricamente las amplitudes de dos fotones, ambas técnicas vuelven a coincidir.

Por otra parte, la medición del desplazamiento de Lamb en \textbf{hidrógeno muónico} ($\mu^- p$) por Pohl et al. (2010) \cite{pohl2010,antognini2013}:
\begin{equation}
\Delta E_{\text{Lamb}} = 206.2949(32)\text{ meV} - 5.2275(10) r_p^2 / \text{fm}^2,
\end{equation}
arrojó $r_p^\mu = 0.84087(39)\text{ fm}$, frente a $0.8775(51)\text{ fm}$ de la dispersión electrónica clásica ($7\sigma$ de discrepancia). El razonamiento que hace tan potente a esta medida es el siguiente: el radio de Bohr muónico ($a_0^\mu \approx a_0^e / 186$) incrementa la densidad de probabilidad en el origen en un factor $|\psi(0)|^2 \propto m_{\text{red}}^3 \approx 10^7$, haciendo al muón hipersensible a la estructura finita del protón. Un muón pasa una fracción muchísimo mayor de su tiempo dentro del volumen del protón que un electrón, de modo que el corrimiento de sus niveles de energía depende del radio de forma mucho más acusada, y basta una espectroscopía de precisión para leerlo. Experimentos modernos sin ventanas como PRad (2019) confirmaron el valor menor~\cite{xiong2019}, convergiendo en el consenso internacional CODATA-2022 \cite{tiesinga2021,pdg2022}:
\begin{equation}
\boxed{r_p = 0.8414 \pm 0.0019\text{ fm}}
\end{equation}

Para sintetizar la evolución metodológica e histórica que condujo a este valor, la Tabla~\ref{tab:proton_radius_experiments} resume las determinaciones experimentales del radio del protón a través de las distintas técnicas físicas discutidas.

\begin{table}[t]
\centering
\small
\renewcommand{\arraystretch}{1.3}
\caption{Panorama comparativo de las determinaciones experimentales del radio de carga del protón $r_p$. La confrontación entre la dispersión elástica $ep$ tradicional, las medidas de transferencia de polarización en JLab, la espectroscopía de hidrógeno muónico ($\mu^- p$) y la dispersión moderna a ultra-bajo $Q^2$ (PRad) desató el «puzzle del radio del protón», resuelto a favor del valor menor mediante correcciones de intercambio de dos fotones (TPE).}
\label{tab:proton_radius_experiments}
\resizebox{0.98\textwidth}{!}{%
\begin{tabular}{p{3.8cm} p{3.2cm} p{3.6cm} p{2.8cm} p{3.6cm}}
\toprule
\textbf{Técnica / Experimento} & \textbf{Sonda / Blanco} & \textbf{Observable Físico} & \textbf{Radio $r_p$ [fm]} & \textbf{Impacto Metodológico / Clave} \\
\midrule
\textbf{Dispersión Elástica Clásica (Rosenbluth)}~\cite{rosenbluth1950,hofstadter1956} & Haz $e^-$ sin polarizar / Protón líquido & Pendiente en el origen:\newline $\langle r_E^2 \rangle = -6\left.\frac{dG_E}{dQ^2}\right|_0$ & $0.877 \pm 0.007$ & Sesgado por amplitudes de dos fotones (TPE) no corregidas a $Q^2$ finito; originó el valor «grande». \\
\addlinespace[0.6ex]
\textbf{Transferencia de Polarización (JLab)}~\cite{jones2000,punjabi2005,gayou2002} & Electrones polarizados / Blanco móvil & Razón de polarizaciones:\newline $P_t/P_l \propto G_E/G_M$ & Determina la caída de $G_E/G_M$ & Inmune a errores sistemáticos de normalización del haz; reveló el quiebre del escalamiento dipolar. \\
\addlinespace[0.6ex]
\textbf{Espectroscopía en Hidrógeno Muónico ($\mu^- p$)}~\cite{pohl2010,antognini2013} & Átomo muónico exótico (CREMA 2010) & Desplazamiento de Lamb:\newline $\Delta E(2S_{1/2} - 2P_{3/2})$ & $0.84087 \pm 0.00039$ & Densidad en el origen $|\psi(0)|^2$ multiplicada por $(m_\mu/m_e)^3 \sim 8\times 10^6$; desató el puzzle a $7\sigma$. \\
\addlinespace[0.6ex]
\textbf{Dispersión $ep$ sin Ventanas (PRad, 2019)}~\cite{xiong2019} & Haz $e^-$ / Gas $H_2$ criogénico sin celdas & Secciones eficaces a ultra-bajo $Q^2 \sim 10^{-4}\text{ GeV}^2$ & $0.831 \pm 0.014$ & Confirmó el radio «pequeño» en dispersión leptónica pura, descartando anomalías exclusivas del muón. \\
\addlinespace[0.6ex]
\textbf{Consenso Internacional CODATA (2022)}~\cite{tiesinga2021,pdg2022} & Promedio ponderado global y ajuste de constantes & Combinación global de espectroscopía y dispersión con TPE & $\mathbf{0.8414 \pm 0.0019}$ & Establece el valor oficial internacional, convergiendo definitivamente en el radio del nucleón moderno. \\
\bottomrule
\end{tabular}%
}%
\end{table}

La resolución de la discrepancia de polarización mediante el intercambio de dos fotones (TPE) y la convergencia de las mediciones del radio del protón hacia el valor de consenso CODATA ilustran cómo la fórmula de Rosenbluth, nacida en 1950, condujo a la física nuclear moderna a cuestionar y refinar sus propios fundamentos perturbativos. A continuación, sintetizamos las conclusiones globales y el balance epistemológico de este desarrollo.

\section{Conclusiones y Balance Epistemológico}
\label{sec:conclusiones}

En este trabajo hemos desarrollado la derivación analítica microscópica y exhaustiva de la estructura electromagnética del protón, adoptando como punto de partida y eje estructurador la sección eficaz diferencial de Rosenbluth formulada en 1950. A lo largo del recorrido analítico formal:
\begin{enumerate}
    \item Demostramos rigurosamente cómo el principio de invarianza de Lorentz y la conservación de la corriente de Noether de QED (identidad de Ward) restringen de forma unívoca el operador de vértice más general de un barión de espín $1/2$ a exactamente dos factores de forma escalares independientes: el factor de Dirac $F_1(q^2)$ y el factor de Pauli $F_2(q^2)$, eliminando el término longitudinal a través de la identidad de Gordon on-shell.
    \item Evaluamos algebraicamente las trazas de Dirac para los tensores leptónico $L_{\mu\nu}$ y hadrónico $W^{\mu\nu}$, mostrando el origen cinemático de los términos de interferencia $F_1 F_2$ en el sistema del laboratorio causados por el retroceso del blanco ($q^0 \neq 0$).
    \item Introdujimos el marco de Breit ($q^0 = 0$) para aislar las densidades estáticas mediante espinores de Pauli bidimensionales, demostrando la cancelación algebraica exacta de los términos cruzados y justificando la introducción de los factores de Sachs $G_E$ y $G_M$.
    \item Integramos el espacio de fases relativista de dos cuerpos invariante de Lorentz ($d\Phi_2$) para cerrar la demostración formal de la fórmula de Rosenbluth, y validamos mediante simulaciones computacionales la técnica de separación lineal $y = mx + b$, evidenciando la pérdida progresiva de sensibilidad en $G_E$ a transferencias elevadas ($Q^2 \ge 2\text{ GeV}^2$).
    \item Tendimos el puente hacia el espacio de coordenadas espaciales reales probando que $G_E(\mathbf{q}^2)$ es la transformada de Fourier tridimensional de la densidad de carga $\rho(\mathbf{r})$, derivamos formalmente la relación de la pendiente en el origen con el radio cuadrático medio $\langle r_E^2 \rangle = -6 \left.\frac{dG_E}{dQ^2}\right|_0$, y mediante integración de contorno y el teorema de residuos de Cauchy en el plano complejo invertimos el modelo dipolar para deducir analíticamente la densidad exponencial y el radio $r_p \approx 0.811\text{ fm}$.
    \item Finalmente, examinamos la discrepancia histórica entre la separación de Rosenbluth y las mediciones de transferencia de polarización en Jefferson Lab ($P_t/P_l$), explicando cómo las correcciones radiativas duras por intercambio de dos fotones (TPE) reconcilian ambos métodos y enmarcan la solución del célebre puzzle del radio del protón.
\end{enumerate}

Conviene extraer de este recorrido dos lecciones que trascienden el caso particular del nucleón. La primera concierne al alcance de los argumentos de simetría. La afirmación de que el protón se describe con exactamente dos factores de forma no dependió en ningún momento de saber de qué está hecho: la establecieron la invariancia de Lorentz, la conservación de paridad y la conservación de la carga, tres principios que sobrevivirán a cualquier revisión futura de la dinámica interna del nucleón. Los argumentos de simetría no predicen valores, pero acotan de manera irrevocable el espacio de lo posible, y esa modestia es justamente su fuerza. La segunda lección es de signo opuesto y tiene que ver con las aproximaciones que se dan por buenas. Durante medio siglo, la aproximación de un solo fotón intercambiado fue una hipótesis de trabajo que nadie consideraba problemática, hasta que dos técnicas experimentales independientes---ambas correctamente derivadas---dieron resultados incompatibles. No falló la teoría, ni fallaron los experimentos: falló una suposición tácita compartida por ambos análisis. Que la fórmula de 1950 siga generando, setenta y cinco años después, preguntas capaces de obligar a revisar los fundamentos perturbativos del cálculo es la mejor medida de su profundidad.

Volviendo al punto donde empezamos: la anomalía que Otto Stern midió en 1933 ha quedado convertida, a lo largo de estas páginas, en el valor en el origen de una función bien definida, $F_2(0) = \kappa_p$; esa función y su compañera han resultado ser medibles mediante la separación lineal de una recta; y la pendiente de una de ellas en el origen ha entregado finalmente una longitud, $r_p = 0.8414 \pm 0.0019\text{ fm}$. Este marco teórico unificado demuestra la vigencia y profundidad de la formulación de Rosenbluth, confirmando que la dispersión elástica leptón-hadrón continúa siendo uno de los pilares más fructíferos y rigurosos de la física subatómica contemporánea.

\appendix
\section{Álgebra de Clifford y Teoremas de Trazas de Dirac}
\label{sec:apendice_dirac}

Este apéndice no es un anexo opcional: contiene las demostraciones de todas las identidades algebraicas que el cuerpo del artículo ha utilizado sin probar, en particular las que hicieron posible evaluar los tensores $L_{\mu\nu}$ y $W^{\mu\nu}$ en la Sección~\ref{sec:cinematica_trazas}. Está escrito para ser autocontenido y puede leerse antes o después del texto principal. El lector que no haya trabajado previamente con matrices gamma encontrará aquí todo lo necesario, y comprobará que un único punto de partida---la relación de anticonmutación de Clifford---genera por manipulación puramente algebraica el conjunto completo de resultados. No se usará ninguna representación matricial explícita: todo se deduce de la regla de anticonmutación y de la propiedad cíclica de la traza~\cite{griffiths2008,halzen1984,thomson2013}.

Las matrices gamma de Dirac satisfacen el álgebra de Clifford en espacio-tiempo de Minkowski ($d=4$):
\begin{equation}
\{\gamma^\mu, \gamma^\nu\} \equiv \gamma^\mu \gamma^\nu + \gamma^\nu \gamma^\mu = 2 g^{\mu\nu} \mathbb{I}_{4\times 4}, \qquad g_{\mu\nu} = \text{diag}(+1, -1, -1, -1).
\end{equation}
La traza de la matriz identidad en cuatro dimensiones es:
\begin{equation}
\text{Tr}[\mathbb{I}_{4\times 4}] = 4.
\end{equation}
Definiendo la matriz quiral $\gamma^5 \equiv i \gamma^0 \gamma^1 \gamma^2 \gamma^3$, se satisfacen las identidades:
\begin{equation}
(\gamma^5)^2 = \mathbb{I}_{4\times 4}, \qquad \{\gamma^\mu, \gamma^5\} = 0, \qquad \gamma^{5\dagger} = \gamma^5.
\end{equation}

\subsection*{A.1 Teoremas Fundamentales de Contracción de Clifford}
Demostremos rigurosamente las cuatro identidades de contracción de matrices de Dirac con índices de Lorentz mudos. Todas ellas siguen el mismo patrón: se anticonmuta una matriz gamma a través del producto usando la relación de Clifford para llevarla junto a su pareja contraída, y se aplica el resultado del caso anterior. Cada identidad se apoya, pues, en la precedente.

\textbf{1. Contracción escalar ($\gamma_\mu \gamma^\mu$):}
\begin{align}
\gamma_\mu \gamma^\mu &= g_{\mu\nu} \gamma^\mu \gamma^\nu = \frac{1}{2} g_{\mu\nu} \{\gamma^\mu, \gamma^\nu\} = \frac{1}{2} g_{\mu\nu} (2 g^{\mu\nu} \mathbb{I}) = g_{\mu\nu} g^{\mu\nu} \mathbb{I} \nonumber \\
&= \delta_\mu^\mu \mathbb{I} = (1 + 1 + 1 + 1)\mathbb{I} = 4 \mathbb{I}.
\end{align}
\begin{equation}
\dest{\gamma_\mu \gamma^\mu = 4 \mathbb{I}}
\label{eq:clifford_contract_0}
\end{equation}

\textbf{2. Contracción con un momento ($\gamma_\mu \slashed{a} \gamma^\mu$):}\\
Utilizando la relación de anticonmutación $\gamma^\mu \slashed{a} = 2 a^\mu \mathbb{I} - \slashed{a} \gamma^\mu$:
\begin{align}
\gamma_\mu \slashed{a} \gamma^\mu &= \gamma_\mu (2 a^\mu \mathbb{I} - \slashed{a} \gamma^\mu) = 2 \gamma_\mu a^\mu - (\gamma_\mu \gamma^\mu) \slashed{a} \nonumber \\
&= 2 \slashed{a} - 4 \slashed{a} = -2 \slashed{a}.
\end{align}
\begin{equation}
\dest{\gamma_\mu \slashed{a} \gamma^\mu = -2 \slashed{a}}
\label{eq:clifford_contract_1}
\end{equation}

\textbf{3. Contracción con dos momentos ($\gamma_\mu \slashed{a} \slashed{b} \gamma^\mu$):}\\
Conmutando $\slashed{b}$ a través de $\gamma^\mu$ usando $\slashed{b} \gamma^\mu = 2 b^\mu \mathbb{I} - \gamma^\mu \slashed{b}$:
\begin{align}
\gamma_\mu \slashed{a} \slashed{b} \gamma^\mu &= \gamma_\mu \slashed{a} (2 b^\mu \mathbb{I} - \gamma^\mu \slashed{b}) = 2 (\gamma_\mu b^\mu) \slashed{a} - (\gamma_\mu \slashed{a} \gamma^\mu) \slashed{b} \nonumber \\
&= 2 \slashed{b} \slashed{a} - (-2 \slashed{a}) \slashed{b} = 2 \slashed{b} \slashed{a} + 2 \slashed{a} \slashed{b} \nonumber \\
&= 2 \{\slashed{a}, \slashed{b}\} = 2 [2 (a \cdot b) \mathbb{I}] = 4 (a \cdot b) \mathbb{I}.
\end{align}
\begin{equation}
\dest{\gamma_\mu \slashed{a} \slashed{b} \gamma^\mu = 4 (a \cdot b) \mathbb{I}}
\label{eq:clifford_contract_2}
\end{equation}

\textbf{4. Contracción con tres momentos ($\gamma_\mu \slashed{a} \slashed{b} \slashed{c} \gamma^\mu$):}\\
Conmutando $\slashed{c}$ usando $\slashed{c} \gamma^\mu = 2 c^\mu \mathbb{I} - \gamma^\mu \slashed{c}$:
\begin{align}
\gamma_\mu \slashed{a} \slashed{b} \slashed{c} \gamma^\mu &= \gamma_\mu \slashed{a} \slashed{b} (2 c^\mu \mathbb{I} - \gamma^\mu \slashed{c}) = 2 (\gamma_\mu c^\mu) \slashed{a} \slashed{b} - (\gamma_\mu \slashed{a} \slashed{b} \gamma^\mu) \slashed{c} \nonumber \\
&= 2 \slashed{c} \slashed{a} \slashed{b} - [4 (a \cdot b) \mathbb{I}] \slashed{c} = 2 \slashed{c} \slashed{a} \slashed{b} - 4 (a \cdot b) \slashed{c}.
\end{align}
Por otra parte, evaluemos la expresión $-2 \slashed{c} \slashed{b} \slashed{a}$:
\begin{align}
-2 \slashed{c} \slashed{b} \slashed{a} &= -2 \slashed{c} [2 (a \cdot b) \mathbb{I} - \slashed{a} \slashed{b}] = -4 (a \cdot b) \slashed{c} + 2 \slashed{c} \slashed{a} \slashed{b}.
\end{align}
Como ambos lados son idénticos término a término:
\begin{equation}
\dest{\gamma_\mu \slashed{a} \slashed{b} \slashed{c} \gamma^\mu = -2 \slashed{c} \slashed{b} \slashed{a}}
\label{eq:clifford_contract_3}
\end{equation}

\subsection*{A.2 Teoremas de Trazas de Dirac}
Las identidades que siguen son las que permitieron convertir el promedio sobre espines en expresiones algebraicas cerradas en la Sección~\ref{sec:cinematica_trazas}. La primera de ellas es, con diferencia, la más rentable: al eliminar de un golpe todas las trazas con un número impar de matrices, reduce a la mitad el número de términos que hay que evaluar en cualquier cálculo.

\textbf{1. Anulación de trazas de un número impar de matrices:}\\
Sea $\mathcal{A} \equiv \gamma^{\mu_1} \gamma^{\mu_2} \dots \gamma^{\mu_{2n+1}}$ el producto de un número impar de matrices. Insertando $(\gamma^5)^2 = \mathbb{I}$ dentro de la traza:
\begin{equation}
\text{Tr}[\mathcal{A}] = \text{Tr}[\mathcal{A} \gamma^5 \gamma^5] = \text{Tr}[\gamma^5 \mathcal{A} \gamma^5].
\end{equation}
Conmutando $\gamma^5$ a través de cada una de las $(2n+1)$ matrices gamma, usando $\gamma^5 \gamma^\mu = -\gamma^\mu \gamma^5$:
\begin{equation}
\gamma^5 \mathcal{A} = (-1)^{2n+1} \mathcal{A} \gamma^5 = - \mathcal{A} \gamma^5.
\end{equation}
Por tanto:
\begin{equation}
\text{Tr}[\mathcal{A}] = \text{Tr}[(-\mathcal{A} \gamma^5)\gamma^5] = -\text{Tr}[\mathcal{A} (\gamma^5)^2] = -\text{Tr}[\mathcal{A}].
\end{equation}
Como $\text{Tr}[\mathcal{A}] = -\text{Tr}[\mathcal{A}]$, se concluye de forma categórica:
\begin{equation}
\dest{ \text{Tr}[\gamma^{\mu_1} \gamma^{\mu_2} \dots \gamma^{\mu_{2n+1}}] \equiv 0 }
\end{equation}

\textbf{2. Traza de dos matrices gamma:}\\
El procedimiento que sigue---y que se repetirá en el caso de cuatro matrices---consiste en usar la relación de Clifford para invertir el orden del producto dentro de la traza, y aprovechar después la invariancia cíclica para reconocer la misma traza en ambos miembros, lo que permite despejarla. Utilizando $\slashed{a}\slashed{b} = 2(a\cdot b)\mathbb{I} - \slashed{b}\slashed{a}$:
\begin{equation}
\text{Tr}[\slashed{a}\slashed{b}] = 2(a\cdot b)\text{Tr}[\mathbb{I}] - \text{Tr}[\slashed{b}\slashed{a}] = 8(a\cdot b) - \text{Tr}[\slashed{a}\slashed{b}].
\end{equation}
Sumando $\text{Tr}[\slashed{a}\slashed{b}]$ a ambos miembros:
\begin{equation}
2\text{Tr}[\slashed{a}\slashed{b}] = 8(a\cdot b) \implies \dest{\text{Tr}[\slashed{a}\slashed{b}] = 4 (a \cdot b)}
\end{equation}
En componentes covariantes:
\begin{equation}
\dest{\text{Tr}[\gamma^\mu \gamma^\nu] = 4 g^{\mu\nu}}
\end{equation}

\textbf{3. Traza de cuatro matrices gamma:}\\
Esta es la identidad que se empleó para evaluar tanto el tensor leptónico como el hadrónico, y por tanto la pieza algebraica más utilizada de todo el artículo. Evaluemos $\text{Tr}[\slashed{a}\slashed{b}\slashed{c}\slashed{d}]$. Conmutando $\slashed{a}$ sucesivamente hacia la derecha mediante $\slashed{a}\slashed{x} = 2(a\cdot x) - \slashed{x}\slashed{a}$:
\begin{align}
\text{Tr}[\slashed{a}\slashed{b}\slashed{c}\slashed{d}] &= 2(a\cdot b)\text{Tr}[\slashed{c}\slashed{d}] - \text{Tr}[\slashed{b}\slashed{a}\slashed{c}\slashed{d}] \nonumber \\
&= 8(a\cdot b)(c\cdot d) - 2(a\cdot c)\text{Tr}[\slashed{b}\slashed{d}] + \text{Tr}[\slashed{b}\slashed{c}\slashed{a}\slashed{d}] \nonumber \\
&= 8(a\cdot b)(c\cdot d) - 8(a\cdot c)(b\cdot d) + 2(a\cdot d)\text{Tr}[\slashed{b}\slashed{c}] - \text{Tr}[\slashed{b}\slashed{c}\slashed{d}\slashed{a}].
\end{align}
Por la propiedad de invariancia cíclica de la traza, $\text{Tr}[\slashed{b}\slashed{c}\slashed{d}\slashed{a}] = \text{Tr}[\slashed{a}\slashed{b}\slashed{c}\slashed{d}]$. Sumando a la izquierda:
\begin{equation}
2\text{Tr}[\slashed{a}\slashed{b}\slashed{c}\slashed{d}] = 8(a\cdot b)(c\cdot d) - 8(a\cdot c)(b\cdot d) + 8(a\cdot d)(b\cdot c).
\end{equation}
Dividiendo entre 2:
\begin{equation}
\dest{\text{Tr}[\slashed{a}\slashed{b}\slashed{c}\slashed{d}] = 4 \left[ (a \cdot b)(c \cdot d) - (a \cdot c)(b \cdot d) + (a \cdot d)(b \cdot c) \right]}
\label{eq:trace_4gamma_proved}
\end{equation}
Nótese la alternancia de signos y la estructura de emparejamientos: los tres términos corresponden a las tres formas distintas de agrupar cuatro objetos en dos parejas, y el signo negativo acompaña al emparejamiento «cruzado». Esta regla mnemotécnica es útil para recordar la identidad sin rederivarla.

\textbf{4. Trazas quirales con $\gamma^5$:}\\
Las identidades que siguen no intervienen en el cálculo no polarizado del cuerpo principal, pero son necesarias para tratar el tensor leptónico polarizado que aparece en la técnica de transferencia de polarización de la Sección~\ref{sec:radio_fourier}. Dado que $\gamma^5 = i \gamma^0 \gamma^1 \gamma^2 \gamma^3$, la traza con menos de cuatro matrices gamma se anula idénticamente:
\begin{equation}
\text{Tr}[\gamma^5] = 0, \qquad \text{Tr}[\gamma^5 \gamma^\mu \gamma^\nu] = 0.
\end{equation}
Para cuatro matrices gamma, la traza produce el tensor totalmente antisimétrico de Levi-Civita:
\begin{equation}
\dest{\text{Tr}[\gamma^5 \gamma^\mu \gamma^\nu \gamma^\rho \gamma^\sigma] = -4i \epsilon^{\mu\nu\rho\sigma}}
\end{equation}
donde adoptamos la convención estándar $\epsilon^{0123} = +1 \implies \epsilon_{0123} = -1$. Las contracciones de dos tensores de Levi-Civita satisfacen:
\begin{equation}
\epsilon^{\mu\nu\alpha\beta} \epsilon_{\mu\nu\rho\sigma} = -2 (\delta_\rho^\alpha \delta_\sigma^\beta - \delta_\sigma^\alpha \delta_\rho^\beta), \qquad \epsilon^{\mu\nu\alpha\beta} \epsilon_{\mu\nu\alpha\sigma} = -6 \delta_\sigma^\beta, \qquad \epsilon^{\mu\nu\alpha\beta} \epsilon_{\mu\nu\alpha\beta} = -24.
\end{equation}

\bibliographystyle{unsrt}
\bibliography{referencias}

\end{document}
